\documentclass[10pt,a4paper,openright]{book}

%% Formateo del título del documento
\title{\Huge Cancionero}
\author{} %\\ es salto de linea
\date{\today}

%% Formateo del estilo de escritura y de la pagina
\pagestyle{plain}
\setlength{\parskip}{0.35cm} %edicion de espaciado
\setlength{\parindent}{0cm} %edicion de sangría
\clubpenalty=10000 %líneas viudas NO
\widowpenalty=10000 %líneas viudas NO
\usepackage[top=2.5cm, bottom=2.5cm, left=1cm, right=1cm]{geometry} % para establecer las medidas de los margenes

\usepackage[spanish]{babel} %Para que el idioma por defecto sea emaketitletraspañol
\usepackage{fontspec}

\usepackage[chorded]{songs}
\newindex{canciones}{idxfile}
% \indexsongsas{canciones}{\thesongnum} 					% Esto es para que se indique el número de canción en lugar de la página
\indexsongsas{canciones}{\thesongnum\ Pag(\thepage)}	% Esto imprime el número de canción junto a la página en la que está

\begin{document}
	% Title
	\maketitle
	\showindex[2]{Indice}{canciones}

	% Songs settings
\noversenumbers
\MultiwordChords
% \renewcommand{\stitlefon}{}
\renewcommand{\lyricfont}{
    % \fontsize{12}{10}\selectfont
    \fontspec{Arial}
}
\renewcommand{\chorusfont}{\bf}
\renewcommand{\printchord}[1]{\color{Blue}\fontspec{Arial}\small #1}
\versesep=16pt                          % separacion entre estrofas
\baselineadj=+8pt plus 1pt minus 0pt   % separacion entre versos
\renewcommand{\clineparams}{            % separación entre acorde y verso
    \baselineskip=14pt
    \lineskiplimit=1pt
    \lineskip=1pt
}
\setlength{\parindent}{0pt}
% \spenalty=-1000
\songpos{1}
\setlength{\cbarwidth}{0pt}
\setlength{\sbarheight}{0pt}
\notenamesin{La}{Si}{Do}{Re}{Mi}{Fa}{Sol} %TODO
\notenamesout{La}{Si}{Do}{Re}{Mi}{Fa}{Sol}%TODO
\renewcommand{\sharpsymbol}{\#}

\renewcommand\makeprelude{
    \resettitles
    {\Large\bfseries\color{Blue}\thesongnum. \songtitle\par
    \nexttitle\foreachtitle{(\songtitle)\par}}
    \vspace{5pt}
    \hrule
    \vspace{5pt}
    \extendprelude
}

% \sepindexesfalse % ???
	
	\songcolumns{2}
	\begin{songs}{canciones}
		% File containing songs that will be included
		\beginsong{A tu lado Señor}[]
	\beginverse
        {\nolyrics Intro \[Do] \[Fa] \[Do]}
    \endverse

    \beginverse
        Jesu\[Do]cristo, yo \[Fa]siento tu \[Do Do7]voz.
        Tú me has dicho: \[Fa]¡Ven y \[Sol]sígue\[Do Do7]me!
        Déjalo \[Fa]todo y \[Sol]dalo a los \[Do Do7]pobres.
        \[Fa]Quiero que \[Sol]seas sal y \[Do Lam]luz
        confía \[Fa]siempre porq\[Sol]ue a tu lado \[Do Fa Do]estoy.
    \endverse

    \beginverse
        Aquí, Se\[Lam]ñor tie\[Mim]nes mi \[Fa]vida,
        que \[Sol]quiere ser presencia de tu \[Do Do7]amor.
        Sé que no es \[Fa]fácil se\[Sol]guir tus \[Do Do7]huellas,
        pero con tu \[Fa]fuerza seré \[Sol]fiel.
    \endverse

    \beginverse
        Te serv\[Lam]iré entre los \[Mim]hombres,
        tu \[Fa]reino anunciaré,
        porque a tu \[Rem]lado \[Fa]quiero cami\[Sol]nar.
    \endverse

    \beginverse
        Te servi\[Lam]ré entre los \[Mim]hombres,
        tu \[Fa]cruz abrazaré.
        Si no re\[Rem]spondo \[Fa]vuélveme a lla\[Sol]mar. \[Do Fa Do]Amén.
    \endverse
\endsong

\beginsong{Dejad las miserias en la puerta}[]
\beginverse
	De\[Do]jad las \[Mim]miserias en la pu\[Lam]erta, \[Fa] 
	y a\[Sol]brid el \[Do]cora\[Fa]zón. \[Sol] 
	\[Do]Unid vues\[Mim]tras ma\[Lam]nos en la mesa, \[Fa]
	sois el \[Sol]cuerpo del Se\[Do...7]ñor.
\endverse

\beginchorus
	\newchords{estribillo}
	Y COMO \[Fa]DIOS {T\[Sol]E} {A\[Do]LABA\[Lam]MOS} HOY,
	Y COMO \[Fa]REY TE \[Sol]\[Do]ADORAMOS \[Lam]YA. 
	Y \[Fa]\[Sol]POR TODO, Y EN \[Do]TI, SE\[Lam]ÑOR, 
	NOS REU\[Fa]NIMOS \[Sol]EN TU HO\[Do]NOR.
\endchorus
\endsong
\beginsong{Comenzar a andar}[]
    \beginverse
        \[Sol Mim Do Re]
    \endverse   
    \beginverse 
        Padre nuestro, ¿qué quieres de mi?
        Ante tus ojos no sé qué decir.
        Padre, lo siento, he vuelto a fallar.
        Arrepentido, quiero cantar:
        Toma mi vida, hazme soñar,
        sólo a tu lado puedo olvidar.
    \endverse
    \beginchorus
        Y al perderme en tus caminos
        Cristo, danos Tú la paz...
        Y al dudar de tus palabras.
        Nunca nos dejes de amar.
    \endchorus
    \beginchorus
        Y, mientras tanto Tú me haces
        callar, me dices:
        ¡Venga hay que actuar!
    \endchorus
    \beginchorus
        Que ya está bien,
        que ya está bien de tanto hablar.
    \endchorus
    \beginchorus
        Y solo entonces Tú me dices:
        ¡Ven, dejalo todo y sígueme!
        Que sabes bien
        Que queda mucho por hacer.
    \endchorus
    \beginchorus
        Señor, yo sé el camino,
        pero es duro seguirlo,
        y ya no sé si, al fin,
        podré volver atrás,
        volver atrás,
        y comenzar a andar,
        volver atrás.
        (Perdona, Señor, yo sé
        el camino y no sé si
        comenzar a andar
        volver atrás.)
    \endchorus
    \beginverse
        Padre nuestro, quiero preguntar, tesiento lejos, he vuelto a dudar. Te necesito para ser feliz. He comprendido, sólo vivo en ti. Nunca te vayas lejos de mí, sólo a tu lado quiero vivir.
    \endverse
    \beginverse
        Y el amor que nos concedes nunca lo dejes de dar. Si el amor nos hace libres, danos tú la libertad.
    \endverse
    \beginverse
        Y el amor que nos concedes nunca lo dejes de dar. Si el amor nos hace libres, danos tú la libertad.
    \endverse
    \beginverse
        Y sólo entonces tú me dices: “¡ven, déjalo todo y sígueme!", que sabes bien que queda mucho por hacer. Señor, yo sé...
    \endverse
\endsong
\beginsong{El Faraón}[]

\beginverse
En E\[Do]gipto esclavo \[Sol]fui sí, sí, \[Do]oh sí
en E\[Do]gipto esclavo fui del \[Sol]rey fara\[Do]ón
triste, muy tris\[Do7]te estaba \[Fa]mi corazón gritaba:
"haz\[Sol]me libre el Señor"\[Do].    Aleluya (bis)
\endverse

\beginverse
Y Moisés fue al faraón, sí, sí, oh sí
y Moisés fue al faraón y le dijo así:
Deja ir a mi pueblo para rendirme culto,
ha dicho el Señor (bis)
\endverse

\beginverse
Faraón se enfureció, sí, sí, oh sí
faraón se enfureció no le dejó salir.
Dios envió diez plagas, desenvainó
la espada, nos dio la libertad. (bis).
\endverse

\beginverse
Libre, Tú me hiciste libre;
Tú me hiciste libre, libre Señor. Aleluya.
\endverse

\beginverse
Ahora que libre estoy, sí, sí, oh sí
ahora que libre estoy, libre para ti
quiero rendirte culto, quiero cantarte
gloria, gloria a Ti ,Señor.
\endverse

\beginverse
Gloria, gloria aleluya (bis) 
gloria, aleluya, aleluya.
\endverse
\endsong
\beginsong{El amor del Señor}[]
\beginchorus
\[Re]EL AMOR DEL SEÑOR
ES TAN MARAVILLOSO.
\[Mi]EL AMOR D\[La7]EL SEÑOR ES \[Mi]TAN
MARAVILLO\[La7]SO.
EL AM\[Re]OR DEL SEÑOR ES TAN
MARAVILLOSO,
\[La7]MARAVILLOSO AM\[Re]OR.
\endchorus
\beginverse
Es tan alto que se sale por arriba,
es tan bajo que se sale por abajo,
es tan ancho que
se sale por los lados, maravilloso amor.
\endverse
        \endsong

        \beginsong{Noé}[]
            \beginchorus
            \[La]ALZAD LAS MANOS Y DAD\[Re]LE LA
            GLO\[Mi]RIA A DIOS (bis)
            ALZA\[Fa#m]D LAS MANOS Y\[Re] ALABADLE
            COMO \[Mi]NIÑOS DEL SEÑOR\[La]
            \endchorus
            \beginverse
            \[La]Dios dijo a Noé, \[Re]construye una b\[Mi]arca.
            \[Fa#m]toda ella hecha, \[Re]de madera como \[Mi]niños del
            \[La]Señor.
            \endverse
            \beginverse
            Los animalitos subían de dos en dos
            el elefante y el canguro
            como niños del Señor.
            \endverse
            \beginverse
            Los animalitos subían de dos en dos.
            los pajarillos y el pingüino
            como niños del Señor.
            \endverse
            \beginverse
            Si tú vas al cielo primero que yo.
            Dile a todos esos angelitos
            que también yo iré.
            \endverse
            
        \endsong
\beginsong{Mueve los pies}[]
\beginverse
\[Re]Mueve tus pies\[La], ¡qué fácil es!
ellos te sirven \[Re]para andar,
\[Re7]para corr\[Sol]er, para saltar,
y acer\[Re]carte a\[la] los de\[Re]más. (bis)
\endverse

\beginchorus
Da las \[Sol]gracias a dios por tus \[Re/La/Re7]pies (3)
y utiliza \[Sol]siempre tus pies... \[Re]
\[La7 (2ª Re)]para acercarte a los demás. (2)
\endchorus

\beginverse
Mueve tus manos, siempre al compás,
ellas te sirven parar jugar,
para comer y acariciar,
y ayudar a los demás. (bis)
\endverse

\beginchorus
Da las gracias... por tus manos (3)
y utiliza siempre tus manos
para ayudar a los demás. (2)
\endchorus

\beginverse
Grita con fuerza ¡hey!,
ésta es tu voz, ¡hey, hey!
ella te sirve para hablar,
para reír (ja, ja)
para cantar,
y a veces también para callar. Pssssshhhh
(2)
\endverse

\beginchorus
ESTRIBILLO
\endchorus
\endsong

\beginsong{Tu amor me llena}[]
    \beginverse
        \[Sol]Cuando estoy en tu pr\[Do]esencia
        m\[Re]e lleno de tu \[Sol]paz.
        \[Do]Cuando miro en tu grandeza
        \[Re]entiendo que eres \[Sol]real
    \endverse
    \beginverse
        Cuando veo tu esencia
        comprendo tu bondad. 
        Cuando admiro tu belleza
        entiendo tu verdad. 
    \endverse
    \beginchorus
        \[Sol]Y TU AM\[Do]OR ME LLE\[Re]NA UNA VEZ \[Sol]MÁS
        \[Mim]AL POD\[Do]ERTE \[Re]HOY CANTAR\[Mim].
        \[Sol]Y TU AM\[Do]OR ME LLE\[Re]NA UNA VEZ \[Sol]MÁS 
        \[Mim]AL POD\[Do]ERTE \[Re]ADORA\[Sol]R.
    \endchorus
\endsong
\beginsong{Estamos todos con Jesús}[]
    \beginverse
        \[Do]Estamos todos con Jesús,
        todos ju\[Fa]ntos con Jes\[Do]ús,
        estamos \[Fa]todos con \[Sol]Jesús \[Do]aquí, oh yeah.
        \[Sol]Viva el am\[Do]or y la ale\[Fa]gría, viva la gente de
        co\[Do]razón, viva la marcha que todos tenemo\[Fa]s,
        viva por \[Sol]siempre nuestra ilusión\[Do]
    \endverse
\endsong
\beginsong{Alegría Amigo}[]
\beginchorus
\[Sol]NO HAS NACIDO, AMIGO,
PARA ESTAR \[Do]TRISTE,
\[Re]AUNQUE LLUEVA EN TU CORA\[Sol]ZÓN
\endchorus

\beginverse
\[Sol]Dime \[Do]dónde tiene \[Re7]ojos, tiene ojos tu
tris\[Sol]teza. (bis)
\endverse

\beginverse
Si tu alma suena, es que lleva, es que lleva
mucho Dios. (bis)
\endverse

\beginverse
Todos tenemos una estrella una estrella
de luz blanca. (bis)
\endverse

\endsong
\beginsong{Si yo tuviera un martillo}[]
\beginverse
\[Do]Si yo tuviera un martillo, tocaría la ma\[Lam]ñana,
tocaría la \[Do]noche, por todo el p\[Sol7]aís.
Alerta el pe\[Do]ligro, debemos un\[Lam]irnos, para
de\[Fa]fend\[Sol]er la \[Do]paz.

Si yo tuviera una campana,
tocaría la mañana, tocaría la noche,
por todo el país.
Alerta el peligro, debemos unirnos,
para defender la paz.

Si yo tuviera una canción,
cantaría a la mañana, cantaría a la noche,
por todo el país.
Alerta el peligro, debemos unirnos,
para defender la paz.

Yo ya tengo el marti\[Do]llo, y yo tengo la
\[Lam]campana, yo ya tengo una canc\[Do]ión con la que
cantar por todo el país\[Sol7] 
Martillo de jus\[Do]ticia, campana de liber\[Lam]tad
y una canción \[Fa]de \[Sol]amo\[Do]r.
\[Do]Hermanos, ohhhh, hermanos, y una \[Fa]canción
\[Sol]de \[Do]amor.
\endverse
\endsong

        \beginsong{Si en verdad}[]
            \beginchorus
            \[Do]Si en verdad Dios te ama \[Sol Do]di amén.
            \[Do]Si en verdad Dios te ama \[Sol Do]di amén.
            Si en ver\[Fa]dad Dios te ama \[Do]testifica con tu
            vida,
            si en ver\[Sol]dad Dios te ama d\[Do]i amén.
            \endchorus
            \beginverse
            . ... da las palmas. . di Jesús.. ... ... taconea..
            ... da las cuatro.
            \endverse
        \endsong
\beginsong{El diablo está enojado}[]
\beginverse
\[Re]El diablo está \[Sol]enojado hay \[La]una razón\[Re].
\[Re]El diablo está \[Sol]enojado hay \[La]una razón\[Re].
\endverse

\beginchorus
El diablo está eno\[Sol]jado hay una razón.
Cristo \[Re]vive en tu corazón y tus pecados\[Sol] 
\[La]son per\[Re]donados. (x2)
\endchorus

\beginverse
Yo leo el Evangelio hay una razón...
\endverse

\beginchorus
Yo leo el Evangelio hay una razón
Cristo vive en tu corazón y tus pecados son 
perdonados. (x2)
\endchorus

\beginverse
Yo amo a María hay una razón...
\endverse

\beginchorus
Yo amo a María hay una razón 
Cristo vive en tu corazón y tus pecados 
son perdonados. (x2)
\endchorus

\beginverse
Me están saliendo alitas hay una razón...
\endverse

\beginchorus
Me están saliendo alitas hay una razón
Cristo vive en tu corazón y tus pecados 
son perdonados. (x2)
\endchorus
\endsong

\beginsong{Feliz en tu día}[]
    \beginverse
        \[La]F\[Re]eliz, fe\[Mi7]liz en tu d\[La]ía
        \[La]amiguito que \[Re]Dios te ben\[Mi]diga
        \[Re]que reine la \[La]paz en tu día
        y que \[Mi7]cumplas muchos \[La]más.\[Mi La]
    \endverse
\endsong
\beginsong{Nada es imposible para Ti}[]
\beginchorus
\[Do]Por qué tengo miedo, \[Fa]si nada es im\[Sol]posible
pa\[Do]ra Ti.
Por qué tengo \[Lam]miedo, si nada \[Fa]es im\[Sol]posible
pa\[Do]ra Ti.
Por qué tengo \[Fa]miedo, si nada \[Sol]es imposible
pa\[Do]ra Ti.
Por qué tengo \[Lam]miedo, si na\[Fa]da es im\[Sol]posible
pa\[Do]ra Ti.
\endchorus

\beginverse
Por qué tengo tristeza si nada es
imposible para ti (4)
Nada es imposible para ti, nada es
imposible para ti.
Por qué tengo dudas...(4)
Enséñame a amar porque...(2)
Enséñame a perdonar porque....(2)
Nada es im posible para ti, nada es...
Tú te hiciste hombre porque... (4)
Tú venciste a la muerte, porque nada...(2)
Tú estás entre nosotros, porque nada ...(2)
Nada es imposible para ti, nada es....
\endverse

\endsong
        \beginsong{Poneos en pie}[]
\beginchorus
\[Re]PONEOS EN \[La]PIE
\[Mim]Y ALZAD LA \[Sim]CABEZA
\[Re]MIRAD QUE \[La]LLEGA
EL \[Sol]HIJO DEL HOMBRE A LA TIE\[La]RRA
\[Mim]Y VIENE A HABITAR VUESTRO HO\[La]GAR.
\[Re]PONEOS EN \[La]PIE Y EN\[Mim]CENDED LA
L\[Sim]ÁMPA\[Re]RA, NO OS HAL\[La]LE DORMIDAS
\[Sol]QUE ARDA VUES\[Re]TRA LLAMA.
\endchorus
\beginverse
\[Fa#m]Velad y \[Do#m]estad atentas
\[Re]leed los signos \[Fa#m]de los tiempos:
\[Sol]El Reino de Dios está cer\[La]ca.

Velad y estad alerta, cuidad la luz en
vuestras velas, el Hijo del Hombre llama a
la puerta.
\endverse
        \endsong

\beginsong{Si os podéis amar ahora}[]
            \beginverse
            \[Do]Si os podéis a\[Mim]mar ah\[Lam]ora con un 
            \[Fa]entrañable \[Sol]amor.
            Si sent\[Do]ís que en \[Sol]vuestro \[Lam]pecho siempre
            est\[Fa]á el amor de Di\[Sol]os.
            \endverse
            \beginverse
            Si donde ha\[Do]bía soledad ahora hay 
            \[Sol]paz y hay per\[Lam]dón.
            Si te\[Fa]néis sus bendiciones, 
            \[Rem]recordad al que os a\[Sol]mó.
            \endverse
            \beginverse
            Se entre\[La]gó \[Mim]para daros \[Fa]vida, 
            \[Rem]sanó las heridas de vuestro co\[Sol]razón.
            \endverse
            \beginverse
            Y en el \[Rem]fondo de la prue\[Sol]ba o en la \[Do]hora de
            l\[Lam]a victoria, ca\[Fa]da día recor\[Sol]dad al que os a\[Do]mó
            ca\[Fa]da día recor\[Sol]dad al que \[Fa Do]os amó.
            \endverse
        \endsong
\beginsong{Abrid las puertas}[]
    \beginverse
        \[Re]Abrid las puertas de par \[Sim]en par
        que va a entrar \[Sol]el Rey de la Glor\[La]ia.
    \endverse
    \beginverse
        Abrid las puertas de par en par
        que va a entrar el Hijo de Dios.
    \endverse
    \beginverse
        Para ce\[Sim]lebrar 2000 años de lu\[Sol]z
        \[Mim]2000 años de miser\[La]icordia.
    \endverse
    \beginverse
        Para celebrar 2000 años de luz
        2000 años de liberación
    \endverse
    \beginverse
        Cerrad los ojos y contemplad
        la presencia del cielo en nosotras
        vivamos juntas, en comunión
        santas, ángeles e hijas de Dios.
    \endverse
    \beginverse
        Cantad con júbilo ante el Señor
        que proclama un año de gracia.
        Alzad las voces para cantar
        y alabar el trono de Dios.
    \endverse
    \beginverse
        Dejad que fluya su gran amor
        y que sane todas las heridas.
        Dejad que os cubra con su gran paz
        y que os llene de gozo y perdón.
    \endverse
\endsong
\beginsong{Abrid las puertas}[]
\beginverse
\[Re]Abrid las puertas de par\[Sim] en par
que va a entrar \[Sol]el Rey de la Glo\[La]ria.
\endverse

\beginverse
Abrid las puertas de par en par
que va a entrar el Hijo de Dios.
\endverse

\beginverse
Para ce\[Sim]lebrar 2000 años de \[Sol]luz
\[Mim]2000 años de mise\[La]ricordia.
\endverse

\beginverse
Para celebrar 2000 años de luz
2000 años de liberación
\endverse

\beginverse
Cerrad los ojos y contemplad
la presencia del cielo en nosotras
vivamos juntas, en comunión
santas, ángeles e hijas de Dios.
\endverse

\beginverse
Cantad con júbilo ante el Señor
que proclama un año de gracia.
Alzad las voces para cantar
y alabar el trono de Dios.
\endverse

\beginverse
Dejad que fluya su gran amor
y que sane todas las heridas.
Dejad que os cubra con su gran paz
y que os llene de gozo y perdón.
\endverse
\endsong

\beginsong{Gracias}[]
\beginverse
\[Do]Hoy Se\[Sol]ñor \[Do7]te daré \[Fa]las gracias \[Do]por \[Sol]mi vivir
por la tie\[Do]rra y \[Fa]mis amigos \[Do]
\[Fa]porque \[Do]siempre fui feliz\[Rem Sol7].
\endverse

\beginverse
Por el \[Do]tronco en \[Sol]que nací \[Do7]
y la \[Fa]savia que \[Do]encontré \[Sol] 
y en los \[Do]brotes \[Fa]que nacie\[Do]ron
\[Fa]portado\[Do]res de \[Sol7]tu fe \[Do]
\endverse

\beginverse
\[Do7]Por las veces \[Fa]que caí y \[Rem]las que me \[Sol]levanté
porque \[Lam]siempre en ellas siempre ví \[Fa]
el a\[Sol]mor de tu poder
\endverse

\beginverse
\[Fa]Por lo bueno que viví y en \[Sol]lo que sentí \[Do]dolor
siempre en todo yo te vi \[Rem]
te doy gra\[Fa]cias \[Do]mi Señor\[Sol].
\endverse
\endsong
\beginsong{Abro mi ser}[]
\beginverse
\[Sol]Abro mi ser y \[Mim]alzo mis manos
\[Do]y mi voz hac\[Re]ia ti.
\[Sol]Quiero rendirme y \[Sim]entregarme
a \[Do]tu voluntad.
\endverse
\beginchorus
\[Lam]MI CORAZÓN SE \[Re]ABRE A LA VOZ
DE TU E\[Sim]SPÍRITU DE A\[Mi]MOR
\[Do]MI ALMA CANTA DE \[Re]GOZO
\[Si7]MI ALMA CANTA DE \[Mim]GOZO
\[Do]MI ALMA CANTA DE \[Lam]GOZO \[Re]EN TI
\[Sol]SEÑOR.
\endchorus
\beginverse
Calma la sed de tu palabra que hay en mi,
oh Jesús, no me abandones ni me dejes
pues confío en ti.
Dame la luz de tu mirada, mírame, oh
Jesús, cambia mi vida con tu fuerza y yo
te alabaré.
\endverse
        \endsong

        \beginsong{Cosas de locos}[]
            \beginverse
            Si pienso \[Re]lo que quieras que yo \[Sol]haga
            si tengo \[Re]que hacer caso a tus pala\[Sol]bras,
            si quieres \[Mim]que te diga lo que pien\[La]so
            que es de \[Sol Re]locos.
            \endverse
            \beginverse
            Si quieres que me quite las cadenas
            que me hacen sentir seguro aquí abajo
            es como si me vaciaras las venas.
            Tú estás loco.
            \endverse
            \beginchorus
            \[Sol]Sólo te pido \[La]fuerzas para hacer
            de mi debili\[Sol]dad un \[La]férreo venda\[Re Re7]val
            \[Sol]desde el convenci\[La]miento que tal \[Re]vez
            hoy todo puede \[Sol]ser de nuevo reali\[Mim]dad
            que ya \[La]está al lle\[Re]gar.
            \endchorus
            \beginverse
            De todas formas sé que es necesario
            andar contra corriente en esta tierra
            y que en el fondo merece la pena estar
            loco.
            \endverse
        \endsong
 \beginsong{Lámpara es tu palabra}[]
\beginchorus
\[Re]LÁMPARA ES TU PA\[Sol]LABRA 
PARA MIS PA\[Re]SOS, 
\[Sol]LUZ EN MI SEN\[RE]DERO. 
(\[Re]LUZ, TU \[Sol]PALABRA \[La]ES LA LUZ)
\endchorus
\beginverse
Yo guardaré tus justos mandamientos. 
Señor, dame vida, según tu promesa. 
\endverse
\beginverse
Mi vida, Señor, está siempre en peligro,
pero no olvido tu voluntad.
\endverse
\beginverse
Inclino mi corazón a cumplir tus leyes. 
Sé tú la alegría de mi corazón.
\endverse
   \endsong

\beginsong{Viviré alabándote}[by={Maite López}]

\beginchorus
\[Sol]Viviré ala\[Re]bándote, ado\[Do]rándote y sir\[Sol]viéndote
\[Mim]toda mi ca\[Sim]pacidad de \[Do]amar es pa\[Re]ra ti (bis).
\endchorus

\beginverse
\[Lam]Todo lo que tengo es \[Re]tuyo,
\[Lam]en todo puedo encon\[Re]trarte,
\[Si7]haz que sepa utili\[Do]zarlo,
\[Sol]sólo si me ayuda a \[Re]amarte.
\endverse

\beginchorus
Viviré alabándote…
\endchorus
\endsong
\beginsong{Me has seducido}[]
    \beginverse
        \[Lam]Señor, no soy n\[Rem]ada.
        ¿Por \[Fa]qué me has llam\[Mi7]ado?
        Has p\[Rem]asado por mi pue\[Sol]rta
        y bien \[Do]sabes que soy \[Mi7]pobre y soy d\[Lam]ébil.
        ¿Por \[Fa]qué te has fijado en \[Mi7]mí?
    \endverse
    \beginchorus
        ME HAS SEDUCIDO, SEÑ\[Lam]OR,
        CON TU MIR\[Fa]ADA.
        ME HAS HAB\[Sol]LADO AL CORAZÓN
        Y ME HAS QUE\[Do]RIDO.
        ES IMPOSIB\[Mi7]LE CONOCERTE 
        Y NO AM\[Lam]ARTE. 
        ES IMPO\[Fa]SIBLE AMARTE 
        Y NO SEGU\[Mi7]IRTE.
        ¡ME HAS SEDUCIDO, SE\[Lam]ÑOR! 
    \endchorus
    \beginverse
        Señor, yo te sigo, y quiero darte lo que pides; 
        aunque hay veces que me cuesta 
        darlo todo. Tú lo sabes, yo soy tuyo. 
        Camina, Señor, junto a mi.
    \endverse
    \beginverse
        Señor, hoy tu nombre es más que palabras: 
        es tu voz que hoy resuena en mi interior, y 
        me habla en el silencio. ¿Qué quieres que haga por Ti?

    \endverse
\endsong
\beginsong{Ven a la fiesta}[]
\beginchorus
\[Mi]VEN \[Si]A LA \[Do#m]FIESTA, (ven tú, podrás)
ES EL \[La]MOMENTO DE (cantar)
RE\[Fa#m]ZAR Y DE \[Si]CANTAR. (vamos a la fiesta)
\[Mi]HOY CE\[Si]LEBRA\[Do#m]MOS (Jesús está)
QUE EN NUES\[La]TRAS VIDAS
DIOS \[Si]VIVIENDO SIEM\[Mi]PRE ESTÁ.
\endchorus

\beginverse
1. \[La]Ven a la \[Si]fiesta a \[Mi]partici\[Do]par;
(uh,uh,ah,gua,gua)
\[La]Nos hace \[Si]falta tu ca\[Do#m]lor.
(uh,uh, nos hace falta tu calo-or ah,gua,gua)
\[Mim]Jesús te in\[Si]vita para celebrar \[Do#m]su amor.
\[Re]Atento tú esta\[Si]rás a responder por eso…
\endverse

\beginverse
2. Ven a la fiesta para recordar
Milagros que renacen hoy:
Jesús hoy viene para con todos estar
Y su vida y su cruz es nuestra luz por
eso…
\endverse

\beginverse
3. Un cielo nuevo hoy vamos a hacer:
amigos, comunidad.
Puertas abiertas queremos siempre tener.
Todos pueden entrar a compartir por eso..
\endverse
\endsong
 \beginsong{Caminaré}[]
\beginverse
{\nolyrics Intro: \[Sim] \[Sol] \[La] x2}
\[Sim]Solo \[Sol]hay cegu\[La]era, oscu\[Sim]ridad \[Sol]de rut\[La]ina 
A pe\[Sim]sar de \[Sol]que mis \[La]ojos no \[Sim]se \[Sol]cierra\[La]n 
\[Sim]Bajo \[Sol]mis pies\[La], arena \[Sim]suave \[Sol]y frí\[La]a 
De des\[Sim]iertos \[Sol]que no me \[La]deja\[Sim]n \[Sol]camina\[La]r 

\[Sol]Pero si cie\[La]rras los ojos \[Sol]hoy 
\[Sol]Las tinie\[La]blas se rin\[Fa#m]den ante su \[Sim]luz 
\[Sol]Y a su pa\[La]so el desier\[Fa#m]to llega a su \[Sim]fin 
La are\[Sol]na con su agua bend\[Re]ita 
En ceme\[Sol]nto se convierte 
Y con paso fir\[La]me; 
\endverse

\beginchorus
CAMINA\[Sim]RÉ, CAM\[Sol]INA\[La]RÉ, 
CAMIN\[Sim]ARÉ, CAM\[Sol]INA\[La]RÉ 
CAMINA\[Sim]RÉ, CAM\[Sol]INA\[La]RÉ 
CAMI\[Sim]NARÉ\[Sol] \[La] 

\[Sol]HABRÁ \[La]DÍAS QUE CREA 
QUE \[Fa#m]TODO MARCHA BI\[Sim]EN 
\[Sol]BAJANDO EL \[La]RITMO 
Y OLVIDA\[Fa#m]NDO TU VERD\[Sim]AD 
\[Sol]Y AVERGONZA\[Re]DO 
\[Sol]POR HABERTE FALLA\[Re]DO 
\[Sol]TE PEDI\[Re]RÉ PERD\[Sim]ÓN 
ME ENGA\[La]ÑÓ LA RUTI\[Sol]NA
\endchorus

\beginverse
{\nolyrics Instrumental: \[Sim] \[Sol] \[La] x2}
Cometí un error creyendo 
que tu camino se recorre 
en días o en meses como mucho 
Ahora he entendido 
que tu camino es mi vida 
y tu vida sirve de guía hacia la vida 

Pero si cierras los ojos hoy 
un soplo de aire renueva todo a tu alrededor 
sacando fuerzas de donde no había fe 
donde todo era ceguera 
su camino me da vida y con paso firme: 
\endverse
\beginchorus
CAMINARÉ…
\endchorus
        \endsong

\beginsong{Mi Entorno}[]

\capo{3}

\beginverse
No \[Lam]sé donde co\[Sol]mienzas \[Do]tú 
no \[Lam]sé dónde \[Sol]termino \[Fa]yo
mi en\[Lam]torno se lle\[Sol]nó de \[Do]ti 
y mi inte\[Sol]rior
\endverse

\beginverse
Y \[Lam]mi manera \[Sol]de pen\[Do]sar 
y \[Lam]mi mane\[Sol]ra de sen\[Fa]tir 
cam\[Lam]biaron al cre\[Sol]cer tu a\[Do]mor 
dentro de \[Sol]mi
\endverse

\beginverse
\[Rem]Y nada cambia\[Lam]ría quiero estar a\[Sol]sí
\[Rem]muriendo cada \[Lam]día para verte en \[Sol]mí
\endverse

\beginchorus
\[Do]Llenas mi es\[Sol]pacio llenas mi \[Lam]vida 
me ilumi\[Fa]naste \[Sol]como el \[Do]sol
quiero que \[Sol]crezcas y todos \[Lam]los días 
ser un re\[Fa]flejo de tu a\[Sol]mor 
\endchorus

\beginverse
Con el sonido de mi voz 
el viaje emprenderé por ti 
mis pies me llevarán a dónde 
quieras ir 
\endverse

\beginverse
Y todo lo que soy Señor 
mis planes y mi corazón
hoy giran a tu alrededor
rendido estoy
\endverse

\beginverse
Y nada cambiaría quiero estas así
muriendo cada día para verte en mí
\endverse

\beginchorus
Estribillo
\endchorus
\endsong
\beginsong{Volar}[]
	\beginverse
		\[Do]Ellos dicen que no se \[Sol]oye, \[Fa]que no sabemos \[Sol]gritar
		\[Do]que no merece \[Mim]la pena \[Lam]luchar\[Fa] por nuestro ideal. \[Sol]
		\[Do]Dicen que son todo p\[Sol]alabr\[Fa]as, nos limitamos a hablar\[Sol]
		\[Do]que entre d\[Mim]ecir y llegar\[Lam] a hacer n\[Fa]os queda mucho que \[Sol]andar
	\endverse

	\beginchorus
		VEN CONMIGO A \[Do(Re)]VOLAR Y EN\[Sol(La)]SEÑAR A LOS \[Lam(Sim)]DEMÁS
		QUE HAY ALG\[Mim]O MÁ\[Fa#m]S DENTRO DE \[Fa(Sol)]TI
		QUE CRISTO TE HA\[Sol(La)]CE VIVIR,
		Y VOLAR Y VOLAR…
	\endchorus

	\beginverse
		Y volar y v\[Do]olar \[Mim]como u\[Lam]n pájaro si\[Fa]n miedo al f\[Sol]racas\[Do]o
		que no di\[Mim]stingue front\[Lam]eras \[Fa]donde se deba frenar, \[Sol]
		\[Fa]para él no hay cielo \[Sol]con \[Fa]nubes pues \[Sol]sólo busca vola\[Do]r.
		\[Do]Vuela rápido hacia lo a\[Sol]lto, \[Fa]que no te quedes \[Sol]atrás
		\[Do]anuncia\[Mim]remos a \[Lam]todos\[Fa] que Él nos ama de ver\[Sol]dad.
		\[Do]Vamos y no tengas mie\[Sol]do; \[Fa]sabes que Él te \[Sol]cuidará
		\[Do]que nos \[Mim]quiere como a \[Lam]hijos, \[Fa]debemos gritar\[Sol]lo ya.
	\endverse

	\beginchorus
		VEN CONMIGO...
	\endchorus

	\beginverse
		Predicare\[Fa]mos al mu\[Sol]ndo que \[Do]por es\[Mim]to hay que lu\[Lam]char
		\[Fa]que no podemos guardarnos nuestra alegría de \[Sol]amar.
		\[Do]Sal afuera \[Lam]y mira lo alto, hay mu\[Fa]cho que conquistar, \[Sol]
		\[Do]no hay fronteras que p\[Lam]aren ya \[Fa]el amor\[Sol] univer\[La]sal.
	\endverse
\endsong

\beginsong{Norte y Sur}[by={CD Vida}]
\beginverse
En al\[Do]gún lu\[Sol]gar le\[Lam]jano\[Fa],
donde el tie\[Do]mpo \[Sol]se ha\[Lam]ce largo\[Fa],
donde el \[Do]bosque \[Sol]no ter\[Lam]mina \[Fa]
y se extien\[Rem]de con su man\[Sol]to
\endverse

\beginverse
Entrégate a ti mismo y escucha tu
corazón, entrega todo lo que eres,
y te \[Rem]senti\[Sol]rás me\[Do...7]jor
\endverse

\beginchorus
IRÉ VO\[Fa]LANDO HAS\[Do]TA TU PUE\[Do]BLO,
DONDE SE\[Fa] UNEN \[Sol]NORTE Y \[Do]SUR,
DONDE EL \[Fa]VERDE ES \[Sol]EL COLOR
\[Do]DEL SUE\[Lam]LO
DONDE \[Fam]DIOS PO\[Sol]NE SU LUZ\[Do] (2)
\endchorus

\beginverse
Voy con ganas de darlo todo,
y volveré con más de lo que di,
Te llevo mi pobre riqueza,
Tú me darás tu gran pobreza
\endverse

\beginverse
Te ayudaré a hacer tu casa,
Tú a construir mi corazón
Descubriré allí tan lejos
que cerca está el amor
que cerca está Dios
\endverse

\transpose{2}
\beginchorus
IRÉ VO\[Fa]LANDO HAS\[Do]TA TU PUE\[Do]BLO,
DONDE SE\[Fa] UNEN \[Sol]NORTE Y \[Do]SUR,
DONDE EL \[Fa]VERDE ES \[Sol]EL COLOR
\[Do]DEL SUE\[Lam]LO
DONDE \[Fam]DIOS PO\[Sol]NE SU LUZ\[Do] (2)
\endchorus

\endsong


\beginsong{El amor (Maite López)}[]
    \beginverse
        E\[Sol]l amor ha \[Do]de traducirse \[Re]en hechos,
        es \[Mim]mucho más que palabras,
        \[Sim]mucho más que sentimientos
        \[Do]obras son amores \[Sol]y no buenas razones,
        \[Lam]el amor no fa\[Re]lla nunca.
    \endverse
    \beginverse
        El amor ha de traducirse en hechos,
        es mucho más que palabras,
        mucho más que sentimientos
        obras son amores y no buenas razones,
        el amor no falla nunca.
    \endverse
    \beginchorus
        EL \[Do]AMOR ES EL \[Si7]SENTIDO \[Mim]DE LA VIDA
        EL \[Do]AMOR ES UN \[Re]DERROCHE \[Sol]DE ALEGRÍA
        EL \[Re]AMOR ES \[Si7]LA CRUZ DE CA\[Mim]DA DÍA
        EL \[Lam]AMOR ES DARLO TODO \[Re]SIN MEDIDA (bis)
    \endchorus
    \beginverse
        El amor es vivir comprometido,
        el amor no es egoísta,
        el amor no es orgulloso,
        el amor todo lo cree,
        el amor todo lo espera,
        el amor no falla nunca.
    \endverse
\endsong
\beginsong{Vivo por ti}[]
\beginverse
\[Dol]Nunca pude imagi\[Re]nar que fueras \[Mim]tú aque\[Sim]lla voz
que lla\[Do]maba \[Sol]desde dentro y es\[Lam]peraba u\[Do]na
respuesta.
\endverse

\beginverse
Has entrado en mi vida y has cambiado todo
aquello que impedía que mi corazón amase de
verdad.
\endverse

\beginverse
\[Mim]Desde entonces soy \[Sim]feliz, y \[Do]aunque \[Sol]no
siempre comprendo
tan sólo \[Do]vivo, tan sólo \[Lam]vivo por tu am\[Re]or
\[Mim]Y a pesar de los e\[Sim]rrores que \[Do]pue\[Sol]da
cometer
mi vida \[Do]gira, mi vi\[Lam]da gira en tor\[Re]no a \[Mi]ti
\endverse

\beginchorus
POR\[La]QUE TÚ ME HAS DADO LA \[Re]VIDA
ME H\[Mi]AS DADO EL A\[La]MOR,
ME HAS DADO RA\[Re]ZONES PARA \[Mi]LUCHAR.
\[La]TÚ LLENASTE \[Re]MI ALMA \[Mi]DE FELI\[La]CIDAD,
ME DISTE UNOS \[Re]OJOS
PA\[La]RA MIRAR AL \[Re]CIELO Y RE\[Mi]ZAR
\endchorus

\beginverse
\[La]Pero si algún \[Mim]día dejo \[Fa#m]de oír tu \[Do#m]voz
y \[Re]pienso que mi \[La]vida es \[Re]una simple ilu\[Mi]sión,
\[Fa#m]no dejes que te \[Mi]olvide, que no \[Re]deje de \[La]
creer.
Mi \[Re]vida gira, mi \[Sim]vida gira en \[Mi]torno a ti. \[Fa#m]
\endverse

\transpose{2}{
\beginchorus
POR\[La]QUE TÚ ME HAS DADO LA \[Re]VIDA
ME H\[Mi]AS DADO EL A\[La]MOR,
ME HAS DADO RA\[Re]ZONES PARA \[Mi]LUCHAR.
\[La]TÚ LLENASTE \[Re]MI ALMA \[Mi]DE FELI\[La]CIDAD,
ME DISTE UNOS \[Re]OJOS
PA\[La]RA MIRAR AL \[Re]CIELO Y RE\[Mi]ZAR
\endchorus
}
\endsong
      \beginsong{Amando hasta el extremo}[]
\beginverse
\[Lam]Déjame, Se\[Do]ñor, mirarte bien por \[Lam]dentro,
entrar en tu Co\[Do]razón y dejarme se\[Fa]ducir
y que aumenten mis de\[Rem]seos de querer ser como \[Sol]Tú,
conocerte intern\[Do]amente, ama\[Mim]rte y seguirte \[Lam]más,
apostar mi \[Fa]vida junto a \[Sol]ti, déjame verte, Señor,
\endverse
\beginchorus
\[Lam]AMANDO HAS\[Sol]TA EL EX\[Fa]TREMO, DEJÁNDOTE LA \[Do]PIEL,
ENTREGANDO LAS EN\[Sol]TRAÑAS, TUS \[Mi]ENTRAÑAS DE MUJER\[Lam],
\[Lam]EN UNA TOALLA Y \[Sol]UN LEBRI\[Fa]LLO, EN UN ACARICIAR LOS PIES,\[Do]
EN UN MIRARNOS HASTA EL FON\[Sol]DO SIN NADA \[Mi]QUE REPROCHAR \[Lam]
Y SIN \[Sol]NADA QUE PE\[Fa]DIR, Y CON T\[Sol]ANTO PARA DAR
\endchorus
\beginverse
Yo, el Maestro y el Señor, ya no puedo amaros más,
pues como el Padre me ha amado, así os he amado yo.
Os dejo mi vida entera en este Vino y este Pan,
este Pan que soy yo mismo que me parto y que me doy.
Mi deseo es que os améis de corazón, yo también os quiero ver
\endverse
\beginchorus
AMANDO HASTA EL EXTREMO, DEJANDOOS LA PIEL,
ENTREGANDO LAS ENTRAÑAS COMO LO HACE UNA MUJER,
EN UNA TOALLA Y UN LEBRILLO, EN UN ACARICIAR LOS PIES,
EN UN MIRAROS HASTA EL FONDO SIN NADA QUE REPROCHAR
Y SIN NADA QUE PEDIR Y CON TANTO PARA DAR.
\endchorus
\beginverse
Sí, te doy todo lo que soy para que sigas amando.
La lucha por la justicia entra en esta intimidad,
que se llena de personas y rostros que acariciar,
que me impulsa desde dentro a comprometerme más.
Todos caben en tu Corazón, quiero seguirte, Señor.
\endverse
\beginchorus
AMANDO HASTA EL EXTREMO, DEJANDOME LA PIEL,
ENTREGANDO LAS ENTRAÑAS, MIS ENTRAÑAS DE MUJER,
EN UNA TOALLA Y UN LEBRILLO, EN UN ACARICIAR LOS PIES,
EN UN MIRARLOS HASTA EL FONDO SIN NADA QUE REPROCHAR
Y SIN NADA QUE PEDIR Y CON TANTO PARA DAR.
\endchorus
        \endsong

\beginsong{Ven tras de mí}[]

\beginverse
{Todo el rato: \nolyrics \[Re Sim Sol La]}
\endverse

\beginverse
Hace ya un tiempo que ocurrió, y todavía
no sé cómo pasó. Sólo recuerdo que lo vi y
sin pensarlo hacia Él acudí.
Hasta aquel día yo pensaba que era feliz
bastó con ver su mirada algo falló en mí
\endverse

\beginverse
Me acerqué, le escuché, le pregunté (x2)
\endverse

\beginchorus
DÉJALO TODO Y VEN TRAS DE MÍ
(Ven tras de mí)
DEJA TU CASA Y VEN TRAS DE MÍA
(Señor yo no puedo, ven tras de mí)
DÉJALO TODO Y VEN TRAS DE MÍ
\endchorus

\beginverse
Hace ya un tiempo que ocurrió y todavía
no sé cómo pasó. Sólo recuerdo que huí, yo
no podía desprenderme de mí. Me
equivoqué lo que tenía no me hizo feliz
Ahora sé lo que su mirada me quiso decir.
\endverse

\beginverse
Acércate, escúchale, pregúntale (x2)
\endverse

\beginchorus
DÉJALO TODO Y VEN TRAS DE MÍ
(Ven tras de mí)
DEJA TU CASA Y VEN TRAS DE MÍA
(Señor yo no puedo, ven tras de mí)
DÉJALO TODO Y VEN TRAS DE MÍ
\endchorus
\endsong
\beginsong{Principio y fundamento}[]

\beginverse
\[Do]Quiero alabarte por el \[Sol]don de la vida,
\[Lam]por poder luchar en mil bata\[Mim]llas,
\[Fa]porque en cada gi\[Fam]ro del cami\[Do]no
\[Rem]Tú estás donde quiera que yo vaya. \[Sol]
\endverse

\beginchorus
CA\[Do]MINANDO CONFIADO \[Fa]QUIERO
\[Do]ALABARTE.
EN MI \[Lam]VIDA DIARIA \[Fa]QUIERO
\[Sol]ADORARTE.
AL \[Do]AMAR A MI HERMANO, \[Fa]QUIERO
\[Do]SERVIRTE.
SEÑOR, \[Lam]DIOS DE ESTE MUNDO,
\[Fa]YO QUIERO AMAR\[Do]TE.
\endchorus

\beginverse
Quiero adorarte en la fe y en la duda,
poner a tus pies fuerza o flaqueza,
atravesar el mundo y en cada destello
adivinar la huella de tu presencia.
\endverse

\beginverse
Quiero servirte sin poner condiciones,
quiero alcanzarte en todas las caricias,
con mis gestos de siempre cantar tu nombre,
\endverse

\beginverse
\[Fa]proclamar tu \[Lam]Reino, que es buena noticia.\[Sol]
\endverse
\endsong


\beginsong{Yo Celebraré}[]
            \beginchorus
             \[Mim]YO CELEBRARÉ \[Do]DELANTE DEL SEÑO\[Re]R,
             CANTARÉ UN CANTO \[Mim]NUEVO. (BIS)
            \endchorus
            \beginverse
            \[Do]Yo le alabaré,
            porq\[Re]ue Él ha hecho g\[Mim]randes cosas. (bis)
            \endverse
        \endsong
\beginsong{Canción al Corazón de Jesús}[]
    \beginverse
        \[Mi]Quiero hab\[La]lar de \[Si]un am\[Mi]or infinito
        \[Si7]que se vuel\[Mi]ve niño, \[Mi7]frágil
        \[La]amor d\[Fa#]e hombre \[Si7]humillado
        \[La]quiero h\[Re]abla\[Fa#m]r
        de un am\[Si7]or apas\[Mi]ionado.
    \endverse
    \beginverse
        Con dolor carga nuestros pecados
        siendo rey se vuelve esclava
        fuego de amor poderoso
        salvador, humilde, fiel, silencioso.
    \endverse
    \beginchorus
        \[Mi]Amor qu\[Si7]e abre sus brazos de acogida
        \[Re]quiero hablar \[La]del cam\[Re]ino ha\[La]cia la vida
        \[Mi]corazón \[La]pacientRe\[]e, amo\[Mi]r ardiente,
        \[Fa#]quiero hab\[Si]lar de aquel que
        \[Sol#]vence la mu\[Do# Fa# (Si7)]erte.
    \endchorus 
    \beginverse
        Quiero hablar de un amor generoso
        que hace y calla amor a todos
        buscándonos todo el tiempo
        esperando la respuesta al encuentro.
    \endverse
    \beginverse
        Quiero hablar de un amor diferente
        misterioso inclaudicable
        amor que vence en la cruz
        quiero hablar del corazón de Jesús.
    \endverse
    \beginverse
        Quiero hablar hoy de un amor
        quiero hablar hoy del Señor
        corazón paciente, amor ardiente,
        quiero hablar de aquel que
        vence a la muerte.
    \endverse
\endsong
\beginsong{Que tu palabra nos cambie}[]

\beginchorus
\[Re]QUE TU PALABRA NOS \[Sol]CAMBIE
EL CORA\[Re]ZÓN,
QUE TU PA\[Sim]LABRA NOS CAM\[Mim]BIE EL
EL \[La]CORAZÓN.
TRANSFORMA EN NIE\[Sim]VE EL \[Fa#m]CARBÓN
DE NUESTRO\[Sol] GRIS\[Re] CORAZÓN.
QUE TU \[Mim]PALABRA NOS CAMBIE
\[La]EL CO\[Re]RAZÓN.
\endchorus

\beginverse
1. Que tu pala\[Rem]bra sea \[Si&]como \[Rem]la linterna
Que nos ense\[Solm]ñe por \[Do]dónde ca\[Fa]minar
Y transfor\[Si&]mar el \[Do]dolor
en \[Fa]una suave can\[Rem]ción
Para lle\[Mi]var al mundo tu ca\[La]lor.
\endverse

\beginverse
2. Que tu palabra sea como la comida
Que nos sostenga y nos haga vivir.
El pan para compartir,
la fuerza para vivir.
Palabra, pan que Dios va a repartir.
\endverse

\beginverse
3. Que tu palabra sea la Buena Noticia
Que todo el mundo sueña con escuchar.
Palabra para salvar,
palabra para curar.
Palabra con poder que nos puede cambiar.
\endverse
\endsong
\beginsong{Gloria a Dios (j.m. cub)}[]
\beginverse
\[Re]No sé como al\[La]abarte, \[Sim]ni que decir, \[Fa#m]Señor.
\[Sol]confío en tu mi\[Re]rada \[Mim]que me abre el co\[La7]razón
\[Re]toma mi pobre \[La]vida q\[Sim]ue, sencilla ante \[Fa#m]ti,
\[Sol]quiere ser ala\[Re]banza \[Mim]por lo que ha\[La]ces en \[Re]mí.
\endverse
\beginchorus
\[Re]GLORIA, \[La]GLOR\[Sim]IA A \[Fa#m]DIOS
\[Sol]GLOR\[Re]IA, G\[Mim]LORIA A \[La]DIOS
\[Re]GLORIA, \[La]GLOR\[Sol]IA A \[Fa#m]DIOS
\[Sol]GLOR\[Re]IA, \[Mim]GLORI\[La7]A A \[La]DIOS
\endchorus
\beginverse
Siento en mí tu presencia, soy como tú me ves bajas a mi miseria, me llenas de tu paz.
Indigno de tus dones, mas por tu gran amor tu espíritu me llena; ¡gracias te doy, Señor!
Gracias por tu palabra, gracias por el amor,
Gracias por nuestra madre, gracias te doy, Señor.
Gracias por mis hermanos, gracias por el perdón.
Gracias porque nos quieres juntos en ti, Señor.
\endverse
        \endsong

\beginsong{Gloria (Yo soy libre)}[]
\beginchorus
\[Re]Gloria, gloria\[Sol Re], Aleluya\[Sol Re], yo soy \[Sol]libre,
oh gloria a \[Re]Dios ah, ah.
Gloria, \[Sim]gloria, Ale\[Fa#m]lu\[Sol]ya, yo soy \[Re]libre,
\[Sol]oh Gloria a \[Re]Dios \[Sol]ah, \[Re]ah.
\endchorus

\beginverse
Mi pasado, se ha borrado, yo soy libre,
oh gloria a Dios ah, ah.
Con el tiempo he cambiado, yo soy libre,
oh gloria a Dios ah, ah.
\endverse

\beginverse
Ni rencores, ni temores, yo soy libre,
oh gloria a Dios ah, ah.
Mis mañanas de colores, yo soy libre,
oh gloria a Dios ah, ah.
\endverse
\endsong
\beginsong{Aleluya (Jesús es el Señor)}[]
\beginchorus
\[Do]Aleluya, a\[Fa]lelu\[Do]ya\[Sol]
\[Fa]Aleluya, \[Sol]alelu\[Do Mim Lam]ya
\[Fa]aleluya, \[Sol]alelu\[Do]ya
\endchorus
\endsong


\beginsong{Aleluya}[]
    \beginchorus
        \[Do]Alelu\[Sol]ya, alel\[Fa]u-aleluy\[Do]a
        \[Fa]Aleluy\[Do]a, alel\[Sol]uya
        \[Do]Alelu\[Sol]ya, ale\[Fa]lu-aleluy\[Do]a
        \[Fa]Aleluya, \[Do Sol Do]Aleluya
    \endchorus
\endsong
\beginsong{Dios sin su pueblo}[]
\beginverse
\[Do]Dio Dios su\[Sol] vida a los \[Lam]niños,
su ama\[Fa]do y querido pue\[Sol]blo.
\endverse

\beginverse
\[Do]En todos dejó \[Sol]su \[Lam]ser,
\[Fa]Él iba en todos \[Sol]ellos.
\endverse

\beginverse
\[Lam]El Dios-Uno-Trinidad
\[Sol]soñó, en la tierra, el cielo
\[Fa]y dejó de ser Dios Solo
\[Mi]para ser Dios-con-su-Pueblo.
\endverse

\beginchorus
\[La]Dios, sin su pu\[Mi]eblo, no qui\[Fa#m]so ser Dios,
no \[Re]quiso ser \[Mi]Dios sin su pue\[La]blo,
no, \[Mi]¡no quiso ser \[La Mi Sol7]Dios!
\endchorus

\beginverse
Por eso dijo a los hombres que creían conocerlo:
"Si no seguís siendo niños
no construiréis mi Reino.
Ya está en vuestro corazón y si no os
parece cierto, olvidasteis el ser niño
porque no os sentís pequeños".
\endverse
\endsong

\beginsong{Aleluya}[by={Kairoi}]
\beginverse
\[Re]Junto a ti, Je\[La]sús reza\[Sim]ré
\[Sol]Tu Palabra \[Mim]quiero vi\[La]vir
\[Re]Tú me haces \[La]sencillo Se\[Sim]ñor
\[Sol]En tus manos \[Mim]me deja\[La]ré
\endverse

\beginchorus
Alelu\[Re]ya \[Fa#m], \[Sol]alelu\[La]ya
Alelu\[Re]ya \[Fa#m],     \[Sol]a\[La]lelu\[Re]ya
\endchorus
\endsong
\beginsong{Aleluya de la tierra}[]
	\beginchorus
	\[Re]Ale\[Sol]luya \[La]canta\[Re]rá
	quien per\[Sol]dió \[La]la espe\[Re]ranza
	y la \[Sol]tierra \[La]sonrei\[Re]rá,
	al\[Sol]eluy\[La]a. \[Re]
	\endchorus
\endsong

\beginsong{Canta Aleluya}[]
\beginchorus
\[Sim]Canta ale\[Fa#m]luya al Se\[Sim]ñor\[Fa#m].
\[Sim]Canta ale\[Fa#m]luya al Señor .
\[Sim]Canta ale\[La]luya, \[Sol]canta a\[Re]leluya
\[Sim]Canta ale\[Fa#m]luya al \[Sim]Señor
\endchorus
\endsong
\beginsong{¿Que te puedo dar?}[]
    \beginverse
        \[Sol]¿Qué te puedo \[Re]dar,
        que no me \[Mim]hayas dado \[Do]Tú?
        \[Mim]¿Qué te puedo \[Re]decir,
        que \[Do]no me hayas dich\[Re7]o Tú?
    \endverse
    \beginverse
        \[Sol]¿Qué puedo ha\[Re]cer por ti?
        \[Mim]Si yo no puedo hacer n\[Do]ada,
        s\[Sol]i yo no puedo \[Re]hacer nada
        si no \[Mi]es por ti.
    \endverse
    \beginchorus
        \[Do]Todo \[Re]lo que \[Sol Sim]sé, todo lo \[Mim]que soy
        \[Do (Lam) Re]Todo lo que tengo \[Fa Mim (Sol)]es tuyo. (bis)
    \endchorus
\endsong
\beginsong{Saber que vendrás}[]
\beginverse
\[Re]En este \[Sol]mundo que \[Fa#m]Cristo nos da\[Sim],
ha\[e]cemos la ofren\[Sol]da del \[La]pan,
\[Re]el pan de nues\[Sol]tro traba\[Fa#m]jo sin fin, \[Sim]
y el \[Re]vino de \[Mim]nuestro can\[La7]tar.
\[Re]Traigo ante \[Sol]Ti nuestra \[Fa#m]justa inquie\[Sim]tud:
\[Re]amar la justi\[Sol]cia y la \[La]paz.
\endverse

\beginchorus
Sa\[Mim]ber que ven\[La7]drás, sa\[Re]ber que \[Sim]estarás
par\[Mim]tiendo a los \[La7]pobres \[Re]tu pan.(Bis)
\endchorus

\beginverse
La sed de todos los hombres sin luz,
la pena y el triste llorar.
El odio de los que mueren sin fe,
cansados de tanto luchar.
En la patena de nuestra oblación,
acepta la vida, Señor.
\endverse
\endsong
\beginsong{Santo (Beatles)}[]
\beginverse
\[Re]Santo, Santo, Santo, Santo,
Santo es el \[Fa#m]Señor, \[Sim]
llenos están el cielo y \[Sol]tierra 
\[Do]de tu amor.\[Re] (Bis)

\[La7]Bendito el que viene en el nom\[Sol]bre
el que \[La]viene en el nombre 
\[Sol]del \[La7]Señor del \[Re]Señor.
\endverse
        \endsong

\beginsong{Santo (Madrid)}[]
\beginchorus
\[Re]Santo eres, Señ\[La]or, Dios nuestro,
\[Sim]rey del cielo y de \[Fa#m]la tierra,
\[Sol]luz de vida eter\[Re]na,
pan de sal\[La]vación.

\[Re]Bendito el que a \[La]ti te busca
\[Sim]y el que tu Pala\[Fa#m]bra escucha;
\[Sol]luchando en tu nom\[Re]bre
la sem\[La]bramos de am\[Re]or.
\endchorus
\endsong
\beginsong{Santo (Hosanna)}[]
    \beginverse
        \[La]Santo, S\[Do#m]anto, \[Fa#m]Santo,
        \[Re]Santo \[Sim]es el Señ\[Mi]or
        \[La]llenos es\[Do#m]tán los c\[Fa#m]ielos
        y la ti\[Re]erra \[Sim]de su gl\[Mi]oria.
    \endverse
    \beginchorus
        \[La]¡Hosanna,(hosanna)
        \[Do#m]Hosanna, (hosanna)
        \[Re]Hosanna \[Mi]en el cie\[La]lo! (b\[Mi]is)
    \endchorus
    \beginverse
        \[La]Bendit\[Do#m]o el que vien\[Fa#m]e
        en el n\[Re]ombre d\[Sim]el Señ\[Mi]or
    \endverse
\endsong
\beginsong{Santo (Jesuitas)}[]
\beginchorus
\[Do]Santo, Santo, Santo-o-o
\[Sol]Santo es el Señor, Dios del Uni\[Fa]vers\[Sol]o
Dios del Univer\[Do]s\[Sol]o. (Bis)
\endchorus

\beginverse
\[Lam]Llenos están el cielo y la tierra de tu \[Sol]gloria
¡Hosanna en el cie\[Fa]lo! (¡Hosanna en el cie\[Sol]lo!)
Hosanna en el \[Do]ciel\[Sol]o
\endverse

\beginverse
Bendito el que viene en el nombre del Señor
¡Hosanna en el cielo! (¡Hosanna en el cielo!)
Hosanna en el cielo—o
\endverse

\beginchorus
\[Do]Santo, Santo, Santo-o-o
\[Sol]Santo es el Señor, Dios del Uni\[Fa]vers\[Sol]o
Dios del Univer\[Do]s\[Sol]o. (Bis)
\endchorus
\endsong
\beginsong{Santo (Rey León)}[]
\beginverse
\[Re]Santo, Santo, \[Sol]Santo, Santo
\[Re]Santo es el Se\[La]ñor,
\[Re]Llenos están el \[Sol]cielo y la tierra \[Re]de su glo\[La]ria

¡Hosanna en el cielo! ¡Hosanna en el cielo!
El nombre del Señor Hosanna en el cielo
¡Hosanna en el cielo! ¡Hosanna en el cielo!
El nombre del Señor Hosanna en el cielo

\endverse
        \endsong

\beginsong{Adelante}[]
    \beginverse
    \nolyrics Intro: \[Rem] \[Sib] \[Do] \[Rem] \[Sib] \[La7] \[Rem]
    \endverse
    \beginverse
       \[Rem]Justo en el momento en que \[Sib]empezaba
       a encontrar oscu\[Do]ridad
       hasta en el sol de mi \[La]ciudad.
       \[Rem]Justo en el momento
       \[Sib]En el que la resignación
       consum\[Rem]ía cada día mi \[La]ilusión.
    \endverse
    \beginverse
        Aparec\[Sib]es tú y me das \[Fa]la mano y
        sin mirarme te acercas a mi lado
        y despa\[Sib]cito me dices susurra\[Do]ndo
        que escuche tu\[Sib] voz.\[Do]
    \endverse
    \beginchorus
        \[Sib]Adelante por los sueños que aún n\[Fa]os quedan,
        \[Sib]adelante por aquellos que está\[Fa]n por venir\[Do],
        \[Sib]adelante porque no importa\[Rem] la meta
        \[Sib]el dest\[Mi7]ino es la prom\[La]esa de \[Rem]seguir.
        Adelante...
    \endchorus
    \beginverse
        Justo en el momento en que empezaba a
        sospechar que la ilusión me abandonó sin
        avisar. Justo en el instante en que
        empezaba a olvidar, a atreverme a
        imaginar, a inventar.
    \endverse
    \beginchorus
        Apareces tú y me das la mano...
    \endchorus
\endsong
\beginsong{Padre Nuestro (mejicano)}[]
\beginverse
Padre \[La]nuestro, que estás en el cie\[Fa#m]lo,
santifi\[Re]cado \[Sim]sea tu N\[Mi]ombre;
\endverse

\beginverse
Venga a \[La]nosotros, venga tu rei\[Fa#m]no;
hága\[Re]se \[Sim]tu volunt\[Mi]ad
\endverse

\beginchorus
Así en la \[La]tierra como en el cie\[Fa#m]lo,
el pan \[Re]nuestro, \[Sim]dánosle ho\[Mi]y;
\endchorus

\beginverse
Y per\[La]dona nuestras ofen\[Fa#m]sas,
así \[Re]como \[Sim]perdona\[Mi]mos
\endverse

\beginverse
No nos \[La]dejes caer, en tentación,\[Fa#m]
líbra\[Re]nos, \[Sim]de todo \[Mi]mal;
\endverse

\beginchorus
A\[La]mén, Amén\[Fa#m];
\[Re]Amén, Amén\[Mi];
A\[La Re La]mén
\endchorus

\endsong
\beginsong{Padre Nuestro (S y F)}[]
    \beginverse
        \[Sim]Padre nuestro tú que es\[La]tás
        en los que aman la ver\[Sim]dad,
        haz que el reino que por \[Sol]ti se dio\[Re],
        llegue pronto a nuestro \[Sol]corazón\[Re],
        y el am\[Sol]or que tu hijo nos de\[Re]jó,
        el am\[Sim]or esté \[La]ya con no\[Sim]sotros.
    \endverse
    \beginchorus
        (se recita el padrenuestro)
    \endchorus
    \beginverse
        Y en el pan de la unidad,
        Cristo danos tú la paz,
        y olvídate de nuestro mal,
        si olvidamos el de los demás.
    \endverse
    \beginverse
        No permitas que caigamos en tentación,
        oh Señor, y ten piedad del mundo.
    \endverse
\endsong
\beginsong{Levanta la mano}[]
\beginchorus
\[Do]Levanta una mano con la \[Fa]palma ha\[Do]cia el cielo
Levanta la otra cruza\[Fa]la no muy le\[Sol]jos.
Si en\[Do]lazas \[Do7]uno con otro pul\[Fa#]gar
ve\[Do]rás la palo\[Sol7]ma de la paz.\[Do Do7]
Si en\[Fa]lazas \[Sol]uno con \[Do]otro pul\[Lam]gar
ve\[Fa]rás una palo\[Sol]ma, ve\[Fa]rás una pa\[Sol]loma \[Do]volar.
\endchorus

\beginverse
\[Fa]Enseña a un a\[Sol]migo el truco\[Do] que te ense\[Lam]ñé
me\[Fa]jor toda\[Sol]vía si le en\[Do]señas a \[Lam]diez.
Y a otros \[Fa]diez cada \[Sol]uno cada \[Do]cual a su \[Lam]vez
\[Fa]así en el cielo tan so\[Sol]lo habrá
\[Fa]muy pronto pa\[Sol]lomas de \[Do]paz.
\endverse
\endsong
  \beginsong{Paz en la tierra}[]
\beginchorus
\[Re]PAZ \[La]EN LA \[Sim]TIERRA,\[La]
\[Re]PAZ EN \[La]LAS AL\[Sim]TURAS.\[La]
\[Sol]QUE EL GOZO E\[La]TERNO \[Re]REINE
\[Mim]EN NUESTRO COR\[La7]AZÓN. \[(2º Re)](bis)
\endchorus
\beginverse
Da \[Re]la paz, her\[La]mano, \[Sol]da la \[Re]paz,
\[Sol]constrúye\[La]la en tu cora\[Re]zón
y con tu \[Fa#m]gesto afirma\[Sol]rás
que \[Mim]quieres la \[La]paz.


Que tu paz, hermano, sea don;
es el mejor signo de amor
que tú nos puedes ofrecer:
abrazo de paz.
\endverse
        \endsong

\beginsong{Paz, Señor}[]
\beginverse
\[Re]Paz Señor, en el \[La]cielo y la tierra,
paz Se\[Sim]ñor, en las olas del \[Fa#m]mar,
\[Sol]paz Señor, en las \[Re]flores que mueven
sin sa\[Mi]berlo la brisa al pa\[La]sar.
\endverse

\beginverse
Tú que has hecho las cosas tan bellas
y les das una vida fugaz,
pon Señor tu mirada sobre ellas
y devuelve a los hombres la paz.
\endverse

\beginverse
Hoy he visto, Señor, en el cielo
suspendidas de un rayo de luz,
dos palomas que alzaban el vuelo
con sus alas en forma de cruz.
\endverse

\beginverse
Haz que vuelvan de nuevo a la Tierra
las palomas que huyeron Señor,
y la llama que enciende la guerra,
se confunda con la paz y el amor.
\endverse

\beginchorus
\[Re]PAZ, \[La]PAZ, \[Sim]PAZ SE\[Sim]ÑOR.
\[Sol]PAZ, \[Re]PAZ, \[Mi]PAZ Y A\[La]MOR. (BIS).
\endchorus
\endsong
\beginsong{Eso no es la paz}[]
    \beginverse
        \[Fa]Hemos construido \[Do]todos una vall\[Sib]a de
        \[Do]metal, co\[Fa]n carteles de \[Do]prohibido el paso
        en las c\[Sib]asas al entrar.
    \endverse
    \beginverse
        \[Fa]Hemos puesto much\[Do]as redes a las \[Sib]alas de
        \[Do]volar, y \[Fa]firmado mil \[Do]papeles que dec\[Sib]imos
        \[Re]que es la paz\[Fa].
    \endverse
    \beginchorus
        \[Fa]¡NO, \[Do]NO, \[Sib]ESO NO ES LA PAZ! (Bis\[Do (Fa)])
    \endchorus
    \beginverse
        Porque la paz no es eso, señores no es un
        beso que sabe a hiel; ni un brazo con
        puñales; sonreír y corta la piel
        Porque la paz no es saber quién puede
        aplastar a ese o aquel; destruir las flores
        que nacen así nos luce el pelo de bien
    \endverse
\endsong
 \beginsong{Como hermanos}[]
        	\beginverse
	        		
\[Do]Apretemos \[Lam]nuestras manos \[Fa]
en torno \[Sol]al altar\[Do],
\[Do]es Jesús \[Lam]centro de vida, \[Fa]fuente \[Sol]de
\[Do]verdad. 
\endverse 
\beginchorus 
\[Fa]SOMOS COMO GRA\[Sol]NOS
DE UNA \[Do]ESPIGA \[Mim]LLENA DE \[Lam]SOL,
\[Rem]SOMOS COMO\[Sol7] UN NUEVO AMA\[Do]NEC\[Do7]ER. 
\[Fa]SEGUIMOS LAS \[Sol]HUELLAS 
DE JE\[Do]SÚS DE \[Mim]NAZARET,  
A \[Rem]TODOS NOS \[Sol7]UNE LA \[Do]FE.
\endchorus
\beginverse
En la vida hay mil motivos para celebrar: un amigo, una mañana, el poder cantar.
La alegría nos desborda, queremos cantar; al sentirnos como hermanos debemos gritar.

        	\endverse
        \endsong

\beginsong{Viva la gente}[]
\beginverse
\[Re]Esta mañana de paseo,
con la \[Sol]gente me encon\[Re]tré.
Al lechero, al cartero, al po\[Mi]licía salu\[La]dé.
De\[Re]trás de cada ventana \[Re7]y pue\[Sol]rta rec\[Re]onocí
a \[Sol]mucha gente \[Re]que an\[Sol]tes ni siquiera \[La]la \[Re]vi.
\endverse
\beginchorus
¡\[Re]VIVA LA GEN\[Sol]TE! LA HAY DONDE \[Re]QUIERA QUE VAS.	
¡VIVA LA GEN\[Mi]TE! ES LO \[La]QUE NOS GUSTA MÁS.
CON \[Re]MÁS GEN\[Re7]TE, A FAVOR \[Sol]DE GENTE
EN CA\[Re]DA PUEBLO Y NACIÓN.
HABRÍA \[Sol]MENO\[La]S GEN\[Re]TE DIFÍCIL
Y MÁS \[Sol]GENTE \[La]CON \[Re]CORAZÓN. HABRÍA...
\endchorus
\beginverse
Gente de las ciudades y también del
interior, la vi como un ejército cada vez
mayor. Entonces me di cuenta de una gran
realidad: las cosas son importantes, pero
la gente lo es más.
\endverse
\beginverse
Dentro de cada uno hay un bien y hay un
mal; mas no dejes que ninguno ataque a la
humanidad. Ámalos como son y lucha por
que sean los hombres y mujeres que Dios
quiso que fueran.
\endverse
   \endsong

\beginsong{Danos tu paz}[]
\beginverse
\[Sol]Danos tu paz, danos la \[Sol7]paz,
danos \[Do]paz Se\[Dom]ñor.
\[Sol]Camina con noso\[Mim]tros 
y encon\[Do]tremos tu a\[Re]mor.
Y uniendo \[Sol]nuestras manos 
\[Sol7]cantaremos \[Do]al cami\[Dom]nar
\[Sol]¡Danos Se\[Re]ñor tu \[Sol]Paz! (bis)
\endverse

\beginchorus
\[Sol]Danos (\[Do]danos), Se\[Sol]ñor (Se\[Do]ñor) \[Re]
\[Sol]Tu paz, \[Sol7]tu paz, \[Paz]tu paz, tu \[Dom]paz,
\[Sol]¡Danos Se\[Re]ñor tu \[Sol]Paz!
\endchorus
\endsong
\beginsong{Un nuevo sitio disponed}[]
\beginverse
\[Fa]Un nuevo sitio disponed para \[Sib]un amigo \[Fa]más.
Un poquitín que os estrechéis
y \[Sol]se podrá se\[Do]ntar.
Pa\[Fa]ra eso sirve la amistad
si \[La7]llega la oca\[Rem]sión:

ha\[Sib]blémosle con \[Fa]libertad y con \[Sib]el cor\[Fa]azón.
Él \[Sib]con su amor nos \[Fa]pagará
y a\[Sib]legrará, alegrará, alegrará \[Sol7]la \[Do7]reunió\[Fa]n.


Un nuevo sitio disponed para un amigo
más,
un poquitín que os estrechéis,
un poquitín que os estrechéis
y se podrá acomodar.
Para eso sirve la amistad,
para estar en reunión,
hablémosle con libertad y con el corazón.
Él con su amor nos pagará y nos dará calor.
Un nuevo sitio disponed, para Él.
\endverse
        \endsong

\beginsong{Blanca Navidad}[]

\beginverse
\[Do]OH, BLANCA NAVI\[La]DAD, sue\[Sol]ño,
\[Fa]y con la nieve alrede\[Do]dor,
blanca es mi quimera,
y es \[Fa]mensajera de \[Do]paz
y \[Rem]de puro amor. \[Sol7]
\endverse

\beginverse
OH, BLANCA NAVIDAD, nieve,
un blanco sueño y un cantar.
Recordar tu infancia po\[Fa]drás,
al lle\[Do]gar la \[Sol]blanca Na\[Do]vidad.
\endverse
\endsong
\beginsong{Canción para la Navidad}[]
    \beginverse
        \[Re]Navidad, es Navida\[Fa#m]d,
        \[Sol]toda la tierra se al\[Re]egra
        y se entris\[La]tece la ma\[Re]r.
        Marinero a dónde vas,
        Deja tus redes y reza,
        Mira la estrella al pasar.
    \endverse
    \beginchorus
        M\[Sol]arinero, marin\[Re]ero
        haz en tu barca un altar,
        marinero, marinero
        porque llegó Navidad. (bis)
    \endchorus
    \beginverse
        Noches blancas de hospital,
        dejad el llanto esta noche,
        que el niño está por llegar.
        Caminante sin hogar
        ven a mi casa esta noche
        que mañana Dios dirá.
    \endverse
\endsong
\beginsong{Noche de paz}[]
\beginverse
1. No\[Sol]che de Dios, noche de paz,
\[Re]claro sol \[Do]brilla \[Sol]ya,
\[Do]y los ángeles cantando \[Sol]están:
¡\[Do]Gloria a Dios, gloria al Rey ce\[Sol]lestial! \[Mim]
\[Re]Duerme el niño Je\[Mim]sús,\[Lam]
\[Sol]duerme el niño \[Re]Jesús. \[Sol]
\endverse

\beginverse
2. Noche de paz, noche de Dios,
al portal va el pastor,
y entre pajas encuentra al Señor,
es el Verbo que carne tomó.
Duerme el niño Jesús,
entre tonadas de amor.
\endverse
\endsong
\beginsong{Adeste fideles}[]
\beginverse
\[Re]Adeste fide\[La]les, lae\[Re]ti tri\[Sol]umph\[Re]ante\[La]s,
\[Sim]venite, v\[La]enite in Betleh\[Mi7]em.
\[La]Nat\[Re]um \[Sol]videt\[Re]e, reg\[La]em ange\[Mi]lorum. \[La]
\endverse
\beginchorus
\[Sim]VENITE ADO\[Re]REMUS, \[Sim]VENITE ADO\[Re]REMUS,
\[Sol]VENITE \[Mi]ADORE\[La]MUS \[Re]DO\[La]MINU\[Re]M
\endchorus
\beginverse
En grege relicto humiles ad cunas
vocati pastores approperant,
et nos ovanti gradu festinemus.

\beginchorus
VENITE...
\endchorus
\beginverse
Aeterni parentis splendorem aeternum,
velátum su carne vidébimus.
Deum infantem pannis involutum
\endverse
        \endsong

\beginsong{Ay del chiquirritín}[]
\beginchorus
\[La]Ay del chiquirritín, chiquirriquitín
Que ha nacido entre \[Mi]pajas.
Ay del chiquirritín, chiquirriquitín,
queridín, queridito del \[La]alma.
\endchorus

\beginverse
2. Por debajo del arco del porta\[Mi]lito
se descubre a María, José y el \[La]Niño.
\endverse

\beginverse
3. Entre un buey y una mula Dios ha nacido
y en un pobre pesebre le han recogido
\endverse
\endsong
\beginsong{Dime Niño}[]
    \beginverse
        Dime niño, \[Sol]¿de quién eres?,
        todo vesti\[Re]dito de blanco. \[Sol](bis)
        \[Do]Soy de la \[Re]Virgen Marí\[Sol]a
        y del \[Lam]Espíritu Santo. \[Re](bis)\[Sol7]
    \endverse
    \beginchorus
        \[Sol]Resuenen con alegría
        \[Re]los cánticos de mi tierr\[Sol]a
        y viva el niño de D\[Lam]ios,
        que ha \[Re]nacido en Noche\[Sol]buena.
        La \[Sol]Nochebuena se vi\[Re]ene,
        la Nochebuena se \[Sol]va,
        y nosotros nos ir\[Re]emos
        y no volveremos m\[Sol]ás. (bis)
    \endchorus
    \beginverse
        Dimo niño, ¿de quién eres
        y si te llamas Jesús?. (bis)
        Soy amor en el pesebre
        y sufrimiento en la cruz. (bis)
    \endverse
\endsong
\beginsong{El tamborilero}[]
\beginverse
1. \[Re]El camino que \[Sol]lleva a Be\[Re]lén
baja hasta el valle que la \[Sol]nieve cu\[Re]brió,
los \[La]pastorcillos quieren ver a su Rey,
le traen re\[Re]galos en su humilde zu\[Sol]rrón,
ropo\[Re]pón, ropo\[La]pón.
\[Re]Ha nacido en un por\[Sol]tal de Be\[Re]lén
\[La]el niño \[Re]Dios\[Sol Re Sol Re Sol La Re].
\endverse

\beginverse
2. Yo quisiera poner a tus pies,
algún presente que te agrade Señor.
Mas tú ya sabes que soy pobre también,
y no poseo más que un viejo tambor,
ropopón, ropopón.
En Tú honor frente al portal tocaré
con mi tambor.
\endverse

\beginverse
3. El camino que lleva a Belén
lo voy marcando con mi viejo tambor,
nada mejor hay que yo pueda ofrecer,
su ronco acento es un canto de amor
al Redentor, al Redentor.
Cuando Dios me vio tocando ante Él,
me sonrió.
\endverse
\endsong
\beginsong{Los peces en el río}[]
\beginverse
\[Mim]La Virgen se está peinan\[Si7]do
entre cortina y cor\[Mim]tina,
sus cabellos son de o\[Si7]ro,
el peine de plata f\[Mim]ina.
\endverse
\beginchorus
PERO MI\[Mim]RA COMO BEBEN
LOS PECES EN EL \[Si7]RÍO,
PERO MIRA COMO BEBEN
POR VER AL DIOS NA\[Mim]CIDO.
BEBEN Y BEBEN Y VUELVEN A BE\[Si7]BER,
LOS PECES EN EL RÍO
POR VER A DIOS NA\[Mim]CER
\endchorus
\beginverse
La Virgen está lavando
y tendiendo en el romero,
los pajarillos cantando,
y el romero floreciendo.
\endverse
\beginchorus
PERO MIRA COMO BEBEN...
\endchorus
\beginverse
La Virgen va caminando
por entre aquellas palmeras,
el Niño mira en sus ojos,
el color de la vereda.
\endverse
\beginchorus
PERO MIRA COMO BEBEN...
\endchorus
        \endsong

\beginsong{Veinticinco De Diciembre}[]
\beginchorus
1. \[Lam]Veinticinco de Di\[Mi]ciembre, \[Lam]fun, \[Mi]fun, \[Lam]fun.
\[Lam]Veinticinco de Di\[Mi]ciembre, \[Lam]fun, \[Mi]fun, \[Lam]fun.
\endchorus

\beginverse
Un ni\[Do]ñito \[Sol7]muy bo\[Do]nito \[Sol7]ha naci\[Do]do en \[Sol7]un por\[Do]tal,
con su ca\[Rem]rita de \[Mi]rosa pa\[Lam]rece una flor her\[Mi]mosa,
\[Lam]fun, \[Mi]fun, \[Lam]fun. (bis)
\endverse

\beginchorus
2. Venid, venid pastorcitos, fun, fun, fun.
Venid, venid pastorcitos, fun, fun, fun.
\endchorus

\beginverse
Venid con la pandereta y castañuelas al portal.
Alegrad al Dios del cielo, que ha aparecido en el suelo.
fun, fun, fun. (bis)
\endverse
\endsong
\beginsong{La marimorena}[]
    \beginverse
        Los pas\[La]tores y zagales
        Camin\[Re]an hacia el porta\[La]l
        Llevando llenos de fruta\[Mi7]s
        Los cestos y el delan\[La]tal.
    \endverse
    \beginchorus
        Ande, ande, ande
        la Marimor\[Mi]ena,
        Ande, ande, ande
        que es la Nochebu\[La]ena. (Bis)
    \endchorus
    \beginverse
        En el portal de Belén, hay estrella, sol y 
        Luna, la Virgen y San José 
        Y el Niño que está en la cuna.
    \endverse
    \beginverse
        En el portal de Belén,
        en invierno es primavera,
        el Mesías esperado
        ha venido ya a la tierra.
    \endverse
    \beginverse
        En el portal de Belén,
        han entrado los ratones 
        y al bueno de San José
        le han roído los calzones.
    \endverse
\endsong
\beginsong{Alegría, Alegría}[]
\beginverse
1. Esta \[Re]noche nace el Niño,
yo no \[La]tengo que llevar\[Re]le:
le llevo mi corazón
que le sir\[La]va de paña\[Re7]les.
\endverse

\beginchorus
Ale\[Sol]gría, alegría, ale\[Re]gría,
Ale\[Sol]gría, alegría y pla\[Re]cer.
Es\[La]ta no\[Re]che nace \[La]el ni\[Re]ño
en el por\[La]tal de Be\[Re]lén.
\endchorus
\endsong
\beginsong{Hacia Belén}[]
\beginverse
\[Re]Hacia Belén va una burra, rin, \[La]rin,
yo me remendaba yo me remendé
yo me eché un remiendo yo me lo quité,
cargada de chocol\[Re]ate.


\[Re]Lleva en su chocolatera rin, \[La]rin
yo me remendaba yo me remendé
yo me eché un remiendo yo me lo quité,
su molinillo y su ana\[Re]fre.
\endverse
\beginchorus
\[La]MARÍA, MA\[Re]RÍA, VEN A A\[La]CÁ CORRIEN\[Re]DO
QU\[Sol]E EL CHOCOLA\[Re]TILLO SE LO ES\[LA]TÁN COMI\[Re]ENDO
MAR\[La]ÍA, MAR\[Re]ÍA, VEN A A\[La]CÁ VOLAN\[Re]DO
QUE \[Sol]EL CHOCOLAT\[Re]ILLO SE LO ES\[Re]TÁN LLEVAN\[Re]DO.
\endchorus
\beginverse
En el portal de Belén, rin, rin,
yo me remendaba yo me remendé
yo me eché un remiendo yo me lo quité,
han entrado los ratones.

Y al bueno de San José, rin, rin
yo me remendaba yo me remendé
yo me eché un remiendo yo me lo quité,
le han roído los calzones.
\endverse
        \endsong

      \beginsong{Como antes}[]
        \beginverse
\[Re]Como antes, más que an\[La]tes, Am\[Sim]or, \[Sol]
sopla tu vien\[Re]to, ven Espíritu \[La]de Am\[Sim]or  \[Sol]
enciende el \[Re]fuego que da \[La]vida en \[Sim]
comunión.
Siembra tu huer\[Re]to, inunda\[La] mi corazón\[Re].

Como antes, más que antes, Amor,
sopla tu viento, ven Espíritu de Amor
alienta el alma, da esperanza, da calor.
Tú me conoces y me amas como soy.

\[Re]Como an\[Re]tes, más que \[Sim]antes \[Sol]Amor\[Re]
Como antes, más que antes Amor.
\endverse
        \endsong

\beginsong{Pastores venid}[]
\beginverse
1. \[Rem]San José al niño Je\[La]sús 
un beso le dio en la ca\[Rem]ra, 
y el niño Jesús le \[Solm]dijo: 
\[La]Que me pinchas con las bar\[Rem]bas. 
\endverse

\beginchorus
\[Rem]Pastores venid, pastores llegad, 
a adorar al \[Lam]niño a adorar al niño 
\[La]que ha nacido \[Rem]ya. 
\endchorus

\beginverse
2. Oiga usted señor José 
no le arrime usted la cara 
que se va a asustar el niño 
con esas barbas largas. 
\endverse

\beginverse
3. Las barbas de San José 
el niño le acariciaba 
y San José sonreía 
cuando de ellas le tiraba.
\endverse
\endsong
\beginsong{}[]
    \beginverse
        \[Re]Me he comprado una zambomba,
        u\[Sol]n pandero y un tambor,
        y \[La]pa’ completar la orqu\[Re]esta
        l\[La]os cacharros del fog\[Re]ón.
    \endverse
    \beginverse
        Coge tú las tapaderas
        que no hay que dejar dormir,
        ni al de arriba ni al de abajo,
        ni al de cien leguas de aquí.
    \endverse
    \beginchorus
        \[Re]Canta, ríe, bebe que hoy es Noche\[Sol]buena
        y en esto\[La]s momentos no hay que ten\[RE]er pena
        dale a la zambomba, dale al al\[Sol]mirez
        y dale a\[La]l puchero que ese sue\[Re]na bien.
    \endchorus
\endsong
\beginsong{Feliz navidad}[]
\beginchorus
\[Sol]Feliz Navi\[Do]dad, Fe\[Re]liz Navidad\[Sol], 
Feliz Navi\[Do]dad 
Próspero \[Re]año y felici\[Sol]dad. 
\endchorus

\beginverse
Aunque haga frío o es\[Do]té lloviendo 
Aunque haga frío o es\[Sol]té lloviendo 
Aunque haga frío o es\[Do]té lloviendo 
Te de\[Re]seo Feliz Navi\[Sol]dad. 
\endverse

\beginverse
Todos los chicos de la Safa, (3) 
Te desean Feliz Navidad.
\endverse
\endsong
 \beginsong{Más allá del rencor}[]
\beginverse
\[Re]Cuando das sin esperar, cuando qui\[Sim]eres de verdad, 
cuando brin\[Sol]das perdón en lugar de \[Mim]rencor, 
hay \[Do]paz en \[La]tu corazón. 
\[Re]Cuando sientes compasión del a\[Sim]migo y su dolor, 
cuando \[Sol]miras la estrella que oculta \[Mim]la niebla, 
hay \[Do]paz en tu cora\[La]zón. 
\endverse
\beginchorus
MÁS \[Re(Mi)]ALLÁ DEL \[Sim(Do#m)]RENCOR, 
DE LAS \[Sol(La)]LÁGRIMAS Y EL DO\[Re(Mi)]LOR 
BRILLA \[Fa#(Sol#m)]LA LUZ DEL AMOR \[Sim(Do#m)]DENTRO 
DE \[Sol(La)]CADA CORA\[La(Si)]ZÓN. 
\[Re(Mi)]ILUSIÓN, NAV\[Sim(Do#m)]IDAD, PON TUS \[Sol(La)]SUEÑOS 
A VO\[Re(Mi)]LAR. 
SIEMBRA \[Fa#(Sol#m)]PAZ, BRINDA A\[Sim(Do#m)]MOR, 
QUE EL MUNDO E\[Sol(La)]NTERO PIDE \[La(Si)]MÁS.
\endchorus
\beginverse
Cuando brota una oración 
cuando aceptas el error, 
cuando encuentras lugar para la libertad, 
hay una sonrisa más. 
Cuando llega la razón 
y se va la incomprensión. 
cuando quieres luchar por un ideal, 
hay una sonrisa más. 
\[Do]Hay un rayo de \[Sol]sol a través del \[La]cristal, 
\[Do]hay un mundo me\[Sol]jor 
cuando aprendes a \[La]amar 
\endverse
\beginchorus
MÁS ALLÁ... 
\endchorus
\beginverse
\[Mi]Cuando alejas el temor 
y pro\[Do#m]digas tu amistad, 
cuando a un \[La]mismo cantar 
has un\[Re]ido tu voz 
hay paz en tu cora\[Si]zón. 

Cuando buscas con ardor 
y descubres tu verdad, 
cuando quieres forjar 
un mañana mejor, 
hay paz en tu corazón.
\endverse
        \endsong

\beginsong{Mi Burrito}[]
\beginverse
\[Sol]Con mi burrito saba\[Do]nero 
voy \[Re]camino de Be\[Sol]lén 
Con mi burrito sabanero 
voy camino de Belén
\endverse

\beginchorus
si me \[Do]ven si me \[Sol]ven 
voy ca\[Re]mino de Be\[Sol]lén 
si me ven si me ven 
voy camino de Belén 
\endchorus

\beginverse
El lucerito mañanero ilumina mi sendero 
El lucerito mañanero ilumina mi sendero 
\endverse

\beginchorus
si me ven si me ven voy camino de Belén 
si me ven si me ven voy camino de Belén 
\endchorus

\beginverse
Con mi cuatrico voy cantando mi burrito 
va trotando con mi cuatrico voy cantando 
mi burrito va trotando 
\endverse

\beginchorus
si me ven si me ven voy camino de Belén 
si me ven si me ven voy camino de Belén 
\endchorus

\beginverse
\[Si7]Tuqui Tuqui Tuqui\[Mim]tuqui 
\[Si7]Tuquituqui Tu qui \[Mim]Tu 
A\[La]púrate mi bu\[La7]rrito 
que ya vamos a lle\[Re]gar 
\endverse

\beginverse
Tuqui Tuqui Tuquituqui 
Tuquituqui Tu qui Tu 
apú\[Do]rate mi bu\[Sol]rrito 
vamos a \[Re]ver a Je\[Sol]sús
\endverse
\endsong
\beginsong{Caminito de Belén}[]
    \beginverse
        \[Do]Ya vienen los Reyes Magos,
        ya vienen los Reyes Magos,
        caminito de Be\[Lam]lén
    \endverse
    \beginchorus
        \[Fa]Olé, olé, Holanda y ol\[Sol]é,
        Holanda, ya se ve, \[Do]ya se ve, ya se ve.
    \endchorus
    \beginverse
        Cargaditos de juguetes,
        cargaditos de juguetes,
        para el Niño de Belén
    \endverse
    \beginchorus
        Olé, olé, Holanda...
    \endchorus
    \beginverse
        La Virgen va caminando,
        la Virgen va caminando,
        caminito de Belén 
    \endverse
    \beginchorus
        Olé, olé, Holanda...
    \endchorus
    \beginverse
        La Virgen va caminando,
        la Virgen va caminando,
        caminito de Belén 
    \endverse
    \beginchorus
        Olé, olé, Holanda...
    \endchorus
    \beginverse
        No pidas agua,mi vida,
        no pidas agua, mi vida,
        no pidas agua, mi bien
    \endverse
    \beginchorus
        Olé, olé, Holanda...
    \endchorus
    \beginverse
        Que los ríos vienen turbios,
        que los ríos vienen turbios,
        y no se puede beber
    \endverse
    \beginchorus
        Olé, olé, Holanda...
    \endchorus
\endsong
\beginsong{Oración del pobre}[]
\beginverse
1. \[Rem]Vengo ante ti, mi Se\[Fa]ñor, 
\[Do]recono\[Rem]ciendo mi cul\[La7]pa;\[Rem]
\[Si&]con \[Do]la \[Fa]fe puesta en tu amor, 
\[Solm]que \[Rem]tú \[La7]me \[Rem]das como a un hijo. 
\[Fa]Te abro mi corazón 
\[Si&]y te \[Do]ofrez\[Fa]co mi miseria; 
\[Solm]des\[Rem]pojado de mis cosas, 
\[La7]quie\[Rem]ro lle\[Re]narme de ti. 
\endverse

\beginchorus
\[Re]QUE\[La7] TU ESPÍRITU, SEÑOR, 
\[Sol]ABRA\[Re]SE TODO MI SER. 
\[La7]HAZME DÓCIL A TU VOZ, 
\[Sol]TRANS\[Re...7]FORMA MI VIDA ENTERA. 
\[Sol]HAZ\[Re]ME DÓCIL A TU VOZ, 
\[Mim]TRANS\[La]FOR\[Re]MA\[Rem] MI VIDA ENTERA. 
\endchorus

\beginverse
2. Puesto en tus manos, Señor, 
siento que soy pobre y débil; 
mas tú me quieres así, 
yo te bendigo y te alabo. 
Padre, en mi debilidad, 
tú me das la fortaleza; 
amas al pobre y sencillo, 
le das tu paz y perdón.
\endverse
\endsong
\beginsong{Amaos, Amaos}[]
            \beginchorus
             \[Do]Amaos (amaos), Ama\[Sol]os (amaos)
             y el m\[Fa]undo creer\[Rem Sol7]á.
             \[Lam]Amaos (amaos) en M\[Do]i (y no temáis)
             porque mi Es\[Sol]píritu os m\[Fa]ueve a la unida\[Do]d.
             (Amaos)
            \endchorus
            \beginverse
            \[Do]Tomad de mí el Amor
            \[Fa]y así lo extenderéi\[Do]s
            \[Lam]vivid, obrad, conta\[Mim]d, sufri\[Fa Rem]d, 
            con comunió\[Sol]n.
            \[Do]Los dos a una, entreg\[Fa]aos a la gran misió\[Do]n:
            \[Lam]irradiar la v\[Mim]ida del Resu\[Fa Rem]citado\[Sol].
            \endverse
        \endsong
\beginsong{Padre me pongo en tus manos}[]
\beginverse
\[Lam]Padre, \[Do]Padre, \[Rem]Padr\[Mi7]e,
me pongo en \[Lam]tus manos.
Haz de mí lo que \[Rem]quieras, sea lo \[Sol]que sea, \[Do] \[Mi7]
te doy las \[Lam]gracias\[Fa]. Lo acepto \[Do]todo\[Sol],
con tal que tu vo\[Do]lunt\[Fa]ad se cumpla en \[Do]mí, \[Sol]
y en todas tus cria\[Do]tura\[Fa]s.
No deseo nada \[Sol]más, Pa\[Fa]dre,
no deseo nada \[Do]más.
Yo te ofrezco mi \[Sol7]alma y te la doy
con todo el amor de que soy ca\[Do]paz,
porque deseo \[Sol7]darme,
ponerme en tus \[Fa]man\[Do]os, sin me\[Fa]dida\[Sol],
con infinita con\[Fa]fianza,
porque tú eres \[Do]mi \[Sol7]Padre.\[Do]
\endverse
        \endsong

\beginsong{Canción de Juan}[]
\beginverse
\[Lam]Tu claridad aquí es os\[Rem]cura
tus cantos son lamentos \[Sol]
tu risa es a\[Do]margura
tu ale\[Mi]gría es hoy tormento. \[Lam]
\endverse

\beginchorus
Hoy ha muerto mi Señor
Que me hacía sonreír,
Hoy se ha ido mi maestro
El que me enseñó a vivir.
\endchorus

\beginverse
Estoy solo, (estoy solo)
Me siento frío, (me siento frío)
Hoy mi vida se derrumba
Han cesado mis suspiros.
\endverse

\beginverse
Ya no hay prisa, (ya no hay prisa)
Por vivir, (por vivir)
Sólo hay miedo en este día 
Que me empuja a desistir.
\endverse

\beginverse
Te me has ido, (te me has ido)
Amigo mío, (amigo mío)
Y toda mi esperanza 
Te la has llevado contigo.
\endverse

\beginverse
Y no puedo, (y no puedo)
Olvidar su rostro, (olvidar su rostro)
Por primera vez fue débil
Su semblante estaba roto.
\endverse

\beginverse
Me ha dejado aquí a su madre 
Me ha otorgado su palabra 
me han clavado en una cruz
María llora y calla
\endverse

\beginverse
Sólo hay llanto en mi alma 
Se ha nublado mi destino 
Solo hay pena en mi esperanza 
Ya no hay luz en el camino
\endverse

\beginverse
Pero Él me dijo, (Él me dijo)
Sé valiente, (sé valiente)
Si caminas con mis pasos 
Yo en ti estaré presente.
\endverse

\beginverse
Hay otros Juanes, (hay otros Juanes)
Entristecidos, (entristecidos)
Que necesitan que les guíes
Que les muestres que estoy vivo
\endverse

\beginverse
No lo dejes, (no lo dejes)
Pon coraje, (pon coraje)
Se requiere mucho valor
Para hacer este viaje
\endverse

\beginverse
Y hoy oigo, (y hoy oigo)
Su silencio, (su silencio)
Que de nuevo me acompaña 
Y me ayuda a seguir viviendo.
\endverse
\endsong
\beginsong{¿Por Qué Lo Condenan?}[]
    \beginverse
        \[Mim]El hombre que dijo: "Dich\[Re]osos los \[Do]pobres".
        Es \[Mim]reo de muerte: ¡\[Sol]crucifica\[Mim]dlo!
    \endverse
    \beginverse
        "Dichosa la gente que perdona".
        Es reo de muerte: ¡crucificadlo!
    \endverse
    \beginverse
        "Feliz el que lucha por la justicia".
        Es reo de muerte: ¡crucificadlo!
    \endverse
    \beginverse
        "Amaos los unos a los otros".
        Es reo de muerte: ¡crucificadlo!
    \endverse
    \beginchorus
        Pero ese h\[Re]ombre, desde \[Do]la cruz
        a todos \[Mim Re Mim]salvará.
        Pero ese hombre, es nuestro Dios,
        y a todos salvará.
    \endchorus
    \beginverse
        El hombre que dijo:"Pobre del rico".
        Es reo de muerte: ¡crucificadlo!
    \endverse
    \beginverse
        "Ay de aquel que siembra el odio".
        Es reo de muerte: ¡crucificadlo!
    \endverse
    \beginverse
        "Ay de aquel que causa escándalo".
        Es reo de muerte: ¡crucificadlo!
    \endverse
\endsong
\beginsong{Eli, Eli}[]
\beginchorus
\[Rem]Eli, \[Lam]Eli, la\[Si&]mma sabacta\[La]ni. 
\[Rem]Eli, \[Lam]Eli, la\[Si&]mma sabacta\[Rem]ni. 
\endchorus

\beginverse
Me \[Si&]han echado por tierra\[Rem], 
me \[Si&]han golpeado con saña\[Rem]. 
Me \[Si&]han arrancado el cabello\[Rem]. 
Me \[Si&]han destruido los ojos\[Rem La]. 
\endverse

\beginchorus
\[Rem]Eli, \[Lam]Eli, la\[Si&]mma sabacta\[La]ni. 
\[Rem]Eli, \[Lam]Eli, la\[Si&]mma sabacta\[Rem]ni. 
\endchorus

\beginverse
Se\[Si&]ñor, me siento mo\[Rem]rir. 
Se\[Si&]ñor, hace frío en mi al\[Rem]ma. 
El \[Si&]rostro cubierto de es\[Rem]putos. 
Los \[Si&]ojos bañados en llan\[Rem La]to 
\endverse
\endsong
\beginsong{Hoy el Señor resucitó}[]
\beginverse
Hoy el \[Do]Señor resucitó
y de la muerte nos \[Sol7]salvó.
\endverse
\beginchorus
¡a\[Do7]legría y paz, he\[Fa]rmano\[Fam]s,
que el S\[Do]eñor \[Sol7]resucitó\[Do]!
\endchorus
\beginverse
Porque esperó, Dios le libró,
y de la muerte lo sacó.
\endverse
\beginverse
El pueblo en Él vida encontró,
la esclavitud ya terminó.
\endverse
\beginverse
La luz de Dios en Él brilló,
la nueva vida nos llenó.
\endverse
\beginverse
Con gozo alzad el rostro a Dios,
que de Él nos llega salvación.
\endverse
\beginverse
Todos cantad: ¡Aleluya!
Todos gritad: ¡Aleluya!
\endverse
\endsong

\beginsong{Alma misionera}[]
    \beginverse
        \[Sol]Señor \[Re]toma mi vida nueva, \[Mim]
        \[Do]antes de que la espera,\[Sol]
        \[Do]desgaste años en mí.\[Re]
    \endverse

    \beginverse
        Estoy \[Sol]dispues\[Re]to a lo que quie\[Mim]ras,
        \[Do]no importa lo que sea, \[Sol]
        tú \[Do]llámame a \[Re]servir. \[Sol Re]
    \endverse

    \beginchorus
        LLEVA\[Sol]ME DONDE LOS HOMBRES \[Re]
        NECESI\[Mim]TEN TUS PALABRAS,\[Do]
        NECESI\[Sol]TEN MIS GANAS DE VIVIR.\[Re...7]
        DONDE \[Sol]FALTE LA ESPERANZA, \[Re]
        DONDE\[Mim] FALTE LA ALEGRÍA,\[Do]
        SIMPLE\[Sol]MENTE POR NO\[Re] SABER DE\[Sol] TI.
    \endchorus

    \beginverse
        Te doy mi corazón sincero,
        para gritar sin miedo,
        lo bueno que es tu amor.
    \endverse

    \beginverse
        Señor tengo alma misionera, condúceme
        a la tierra, que tenga sed de Dios.
        Y así en marcha iré cantando
        por calles predicando
        lo bello que es tu amor
    \endverse

    \beginverse
        Señor, tengo alma misionera,
        condúceme a la tierra
        que tenga sed de Dios.
    \endverse
\endsong

\beginsong{Cristo es la luz}[]
\beginverse
\[Re]De la sombra de la \[Sol]muer\[Re]te (2)
\[Sim]surge Cris\[Re]to nueva \[Sim]vi\[La]da. (2)
\[Re]De la sombra del pe\[Sol]ca\[Re]do (2)
\[Sim]un nuevo \[Re]hombre re\[Re]su\[La]ci\[Re]ta. (2)
\endverse

\beginchorus
\[Re]CRISTO ES LA \[Fa#m]LUZ,
\[Sol]CRISTO ES EL \[Re]SOL.
CRISTO ES EL \[Sol]RESPLAN\[La7]DOR.
\[Re]CRISTO ES LA \[Fa#m]LUZ,
\[Sol]CRISTO ES EL \[Re]SOL
QUE A LA NO\[La7]CHE VEN\[Re]CIÓ.
\endchorus

\beginverse
Las tinieblas de la noche (2)
se han rasgado con el día. (2)
Las tinieblas de la muerte (2)
en Jesús quedan vencidas. (2)
\endverse
\endsong
\beginsong{Jesús resucita hoy}[]

\beginverse
\[Re]Mirad, Je\[Sol]sús resu\[La]cita \[Re]hoy.
\[Sol]Mirad, la \[La]tumba está va\[Re]cía.\[Sim]
El \[La]Padre ha pensado en \[Sim]Él.
De los \[Sol]hombres es Se\[Mim]ñor,
de la \[Sol]vida Salva\[La]dor.
\endverse

\beginverse
Mirad, Jesús resucita hoy.
Mirad, vive a nuestro lado.
La muerte no tiene poder.
Procla\[Sol]mad por la \[Mim]fe que está \[Sol]vivo
y somos \[La]libres porque...
\endverse

\beginchorus
\[Re]ÉL RESUCITA \[Sol]HOY.
\[La]ÉL VIVE ENTRE {NO\[Re La]SOTROS}.
\[Re]ES \[Re7]CRISTO, ES SE\[Sol]ÑOR\[Solm].
ALELU\[Re]YA, \[Sol]ALELU\[La (2ª Re)]YA. (bis)
\endchorus

\beginverse
Mirad, Jesús resucita hoy nos da la paz
con su palabra el gozo vuelve al corazón.
Con su Espíritu de amor nuestra vida
cambiará.
\endverse

\beginverse
Mirad, Jesús resucita hoy su amor no nos
dejará su fuerza nos empujará.
Él será guía y luz,
esperanza y fortaleza
porque...
\endverse

\beginchorus
ESTRIBILLO
\endchorus
\endsong
\beginsong{Pregón Pascual}[]
    \beginverse
        \[Do]Exulten los \[Fa]coros de los \[Do]ángel\[Fa]es
        \[Do]exulten los \[Fa]ministros de \[Do]Dios\[Sol7]
        y que \[Lam]suenen las trom\[Mim]petas de vic\[Fa]toria
        por el tr\[Sol7]iunfo de Jesús, nuestro Se\[Do Sol]ñor.
    \endverse
    \beginverse
        Que se alegre y se goce esta fiesta
        inundada de tanta claridad,
        que se sienta libre de la oscuridad
        porque las tinieblas Él venció.
    \endverse
    \beginchorus
        \[Do]GLORIA, \[Sol]GLORIA, \[Fa]GLORIA, \[Sol]GLORIA,
        \[Do]GLORIA, \[Sol]GLORIA, \[Fa]GLORIA, \[Sol7]GLORI\[Do]A.
    \endchorus
    \beginverse
        La Iglesia también se alegrará,
        revestida de tan brillante luz.
        Que resuene este templo, al aclamar
        nuestras voces cantando su bondad.
    \endverse
    \beginverse
        Esta es la fiesta de la Pascua,
        esta es la noche en que llegó,
        la libertad a su pueblo Israel,
        al sacarlo de su esclavitud.
    \endverse
    \beginverse
        Esta es la noche en que se rompen
        las cadenas del pecado, por Jesús.
        Esta es la noche en que Cristo,
        a la muerte con su muerte derrotó.

    \endverse
\endsong
\beginsong{Pregón Pascual 2}[by={Kairoi}]

\beginchorus
\[Re]Alelu\[La]ya, Cris\[Sol]to ha resucita\[La]do
\[Re]Alelu\[La]ya, Cris\[Sol]to vi\[La]ve \[Re]hoy (Bis)
\endchorus

\beginverse
\[Re]Esta es \[La]la noche de \[Sim]la vida, de \[Re]la luz de la 
\[Sol]resurre\[La]cción, porque el Pa\[Re]dre por la fuerza\[Sol]
del Es\[La]píritu resuci\[Re]tó a Jesús \[Sol]de Nazaret\[La].
Perse\[Sim]guido, tor\[Fa#7]turado, asesinado \[Sol]
por \[La7]la causa de los pobres,
por a\[Re]mor hacia su \[Sol]pueblo Isra\[La]el.
\endverse

\beginverse
Esta es la más santa de las noches, la más 
llena de alegría por Jesús, nuestros cantos
de alegría, nuestros gritos, la Iglesia, con 
su fuerza, llenarán. Es la Iglesia de los 
pobres y oprimidos, la comunidad de 
hermanos es el Reino que se vive hoy aquí
\endverse

\beginverse
Muerta está la muerte cruel e impuesta, y
este grito de esperanza triunfará, por 
montañas de justicia, por torrentes,
repletos de fraterna igualdad. En 
creatividad continua, permanente pues 
queremos ser mejores al servicio de la 
causa popular.
\endverse

\beginverse
Noche de hermandad, noche de vida, la 
semilla de una nueva sociedad. Que se
apaguen ya los odios y rencores y por fin se 
restablezca la amistad. Nos unimos a los
gritos de esta tierra que recibe tu
Evangelio que reclama a gritos la liberación.
\endverse
\endsong
%      \beginsong{Secuencia de Pascua}[]
\beginverse
{\nolyrics  \[Do] \[Sol] \[Do]}
{\nolyrics  \[Do] \[Sol]}
{\nolyrics  \[Lam] \[Mim]}
{\nolyrics  \[Do] \[Fa] \[Sol] \[Di]}
Ofrezcan los cristianos
ofrendas de alabanza
a gloria de la Víctima
propicia de la Pascua.
\endverse
\beginverse
Cordero sin pecado
que a las ovejas salva,
Dios y a los culpables
unió con nueva alianza.
\endverse
\beginverse
Lucharon vida y muerte
en singular batalla,
y, muerto el que es la vida,
triunfante se levanta.
\endverse
\beginverse
¿Qué has visto de camino,
María, en la mañana?
<<A mi Señor glorioso,
la tumba abandonada,
\endverse
\beginverse
los ángeles testigos,
sudarios y mortaja.
¡Resucitó de veras
mi amor y mi esperanza!
\endverse
\beginverse
Venid a Galilea,
allí el Señor aguarda;
allí veréis los suyos
la gloria de la Pascua>>
\endverse
\beginverse
Primicia de los muertos,
sabemos por tu gracia
que estás resucitado
la muerte en ti no manda.
\endverse
\beginverse
Rey vencedor, apiádate
de la miseria humana
y da a tus fieles parte
en tu victoria santa.
\endverse
        \endsong

\beginsong{Sólo tu Espíritu}[]
\beginverse
\[Sol]Perdóname Se\[Re]ñor
\[Do]ando dis\[Re]perso en mis \[Sol Mim]cosas, 
me perdí en la me\[Sim]diocridad, 
\[Do]mi exis\[Lam]tencia está va\[Re]cía. 
\endverse

\beginverse
Perdóna\[Sol]me, Se\[Re]ñor,
\[Do]ni siquiera \[Re]hoy he ora\[Sol Mim]do. 
Lo he dejado por co\[Sim]modidad, 
\[Do]de mi vi\[Lam]da no eres \[Re]centro. 
\endverse

\beginchorus
Sólo tu Es\[Do]píritu \[Re]me mantiene vi\[Sol]vo 
cuando \[Do]todo pa\[Lam]rece que va \[Re]mal. 
Si tu Es\[Sol]píritu es\[Si7]tá con\[Mim Do]migo
soy fuer\[Sol]te en \[Lam7]mi debili\[Re]dad. \[(2ª Sol)] 
\endchorus

\endsong
\beginsong{A tu aire}[]
\beginverse
\[Do]A tu aire se vive me\[Sol]jor pero más ex\[Lam]puesto
a la \[Do]duda, a la incompren\[Sol]sión,
A ver al descu\[Lam]bierto todos \[Fa]esos errores,
incoherencias y men\[Sol]tiras en los que \[Fa]vamos
cayendo a lo largo de la \[Sol]vida.
\endverse

\beginverse
A tu aire el pan se llama pan y el vino, vino.
El futuro es cosa de profetas no de
adivinos.
Y entre el cielo y el suelo hay lugar para
todo lo humano, sin tener que decir o
juzgar si algo es bueno o es malo.
\endverse

\beginchorus
A TU \[Do]AIRE ME \[Fa]LANZO A VO\[Sol]LAR, 
A TU \[Do]AIRE DES\[Fa]PLIEGO MIS VELAS
Y \[Lam]SIENTO QUE ERES UN HURA\[Sol]CÁN 
QUE REMUEVE MI \[Lam]VIDA,
BRISA \[Do]QUE RECONFORTA Y
\[Sol]SANA MIS HE\[Lam]RIDAS.
CUANDO EL \[Fa]VIENTO HABLA
TIENES QUE ESCU\[Sol]CHAR,
DESCU\[Fa]BRIR QUE TUS ALAS SIRVEN
PARA VO\[Sol]LAR.
\endchorus

\beginverse
A tu aire no queda lugar para el disimulo.
Con tus notas se decide lo que es claro u
oscuro.
Y aún con miedo y temblando es mejor
romper los muros que me atan a una vida
gris y seguir desnudo.
\endverse

\beginchorus
A TU \[Do]AIRE ME \[Fa]LANZO A VO\[Sol]LAR
A TU AIRE DESPLIEGO MI SER.
A TU AIRE QUERIENDO SOÑAR.
A TU AIRE VOLVIENDO A NACER.
A TU AIRE.
\endchorus
\endsong
\input{./canciones/26.tex}
\beginsong{Ven, espíritu de Dios}[]
    \beginchorus
        \[Sim]VEN ESPÍRITU DE D\[Re]IOS SOBRE \[La]MÍ,
        ME A\[Mim]BRO A TU PRES\[Sim]ENCIA,
        CAMBIAR\[Sol]ÁS MI \[Re]CORAZ\[La]ÓN. \[Fa#](bis)
        CAMBIAR\[Sol]ÁS MI \[Re La Re]CORAZÓN\[Fa#].
    \endchorus
    \beginverse
        \[Sim]Toca mi debili\[La]dad, \[Mim]toma todo lo que \[Sim]soy.
        \[Sol]Pongo mi \[Fa#]vida en tus \[Re]manos \[Mi Sol La]y mi fe.
        \[Sim]Poco a poco llegar\[La Mim]ás a inundarme de tu \[Sim]luz.
        \[Sol]Tú cambiar\[Fa#]ás mi pa\[Re]sado. \[Mi Sol La Fa#]Cantaré.
    \endverse
    \beginverse
        Quiero ser signo de paz,
        quiero compartir mi ser.
        Yo necesito tu fuerza, tu valor.
        Quiero proclamarte a ti,
        ser testigo de tu amor.
        ¡Entra y transforma mi vida! ¡Ven a mí!
    \endverse
\endsong
\beginsong{Bendigamos al Señor}[]
    \beginverse
        \[Sol]Bendigamos al Señor,
        Dios de \[Mim]toda la creación,
        por ha\[Do]bernos rega\[Lam]lado su \[Re7]amor.
        Su bon\[Sim]dad y su perdón
        y su \[Mim]gran fidelidad
        por los \[Do]siglos de los \[Lam]siglos dura\[Sol7]rán.
    \endverse

    \beginchorus
        \[Do]EL ESPÍRITU DE \[Re7]DIOS
        HOY ES\[Sol]TÁ SOBRE \[Mim]MÍ
        Y \[Lam]ES QUIEN ME HA UNGIDO\[Re7]
        PARA PROCLA\[Sol]MAR\[Mim]
        \[Do]LA BUENA \[Re7]NUEVA A LOS MÁS \[Sol]PO\[Mim]BRES,
        \[Do]LA GRACIA \[Lam]DE SU SALVA\[Re7 (2ª SOL)]CIÓN.
    \endchorus

    \beginverse
        Enviados con poder
        y en el nombre de Jesús,
        a sanar a los enfermos del dolor,
        a los ciegos dar visión,
        a los pobres la verdad
        y a los presos y oprimidos libertad.
    \endverse

    \beginverse
        Con la fuerza de su amor
        y de la Resurrección,
        anunciamos llega ya la salvación;
        que ni el miedo ni el temor,
        ni la duda o la opresión
        borrarán la paz de nuestro corazón.
    \endverse
\endsong
 \beginsong{El espíritu del Señor}[]
\beginchorus
EL \[Sim]SEÑOR OS DARÁ
SU ES\[La]PÍRITU \[Sim]SANTO.
\[Re]YA NO TE\[Sol]MÁIS, AB\[La]RID EL CORA\[Re]ZÓN.
\[Fa#]DERRAM\[Sim]ARÁ \[Mim]TODO SU A\[Fa#]MOR. (bis)
\[Fa#]DERRAMA\[Sim]RÁ \[Fa#]TODO SU A\[Sim]MOR (2ª vez)
\endchorus
\beginverse
\[Sim]Él transforma\[La]rá hoy vuestra \[Sim]vida
os da\[Re]rá la \[La]fuerza para a\[Re]mar.
No per\[Si7]dáis vuestra espe\[Mim]ranza,
\[La]Él os sal\[Re]vará. \[Fa#]
\endverse
\beginverse
Él transformará todas las penas,
como a hijos os acogerá.
Abrid vuestros corazones
a la libertad.
\endverse
\beginverse
Fortalecerá todo cansancio,
si al orar dejáis que os dé su paz.
Brotará vuestra alabanza, Él os hablará.
Os inundará de un nuevo gozo,
con el don de la fraternidad.
Abrid vuestros corazones a la libertad.
\endverse
        \endsong

\beginsong{Bendigamos al Señor}[]
\beginverse
\[Re]Somos un grupo de a\[La]migos
\[La7]que queremos can\[Re]tar
con ilusión y ale\[La]gría
\[La7]a un coleguilla espe\[Re]cial.
Vamos a empe\[La]zar, \[La7]prepára\[Re]te:
Uno, Dos y Tres:
\endverse

\beginchorus
\[Re]LEVÁNTATE, LEVÁNTATE Y SA\[La]LUDA
SIÉNTATE, SIÉNTATE Y SON\[Re]RÍE
SOMOS FE\[La]LICES
\[Sol]PORQUE HOY ES\[La]TAMOS DE \[Re]FIESTA
\[Mim]CON EL HER\[La]MANO GA\[Re]BRIEL
\endchorus

\beginverse
Muchas cosas nos enseña
para vivir como él.
Si tú quieres es fácil
sólo presta atención:
\endverse

\beginverse
Compar\[Sol]tir, am\[La]ar y jug\[Re]ar
perdon\[Mim]ar, traba\[La]jar y rez\[Re]ar.
\endverse
\endsong
\beginsong{Hemos oído}[]
\beginverse
\[Do]Hemos oído hablar \[mim]mucho de \[Lam]ti,
Her\[Fa]mano Gabri\[Do]el.
Ahora conocemos mejor \[Mim]la misión\[Lam]
\[Fa]que el Se\[Do]ñor te \[Rem]encomen\[Sol7]dó.
\endverse

\beginverse
\[Do]Desde muy joven oís\[Mim]te su \[Lam]voz
\[Fa]su luz te \[Do]guió,
\[Fa]con decisión, apo\[Mim]yado en la \[Lam]fe,
respon\[Fa]diste por a\[Sol7]mor.
\endverse

\beginchorus
\[Do]Cuenta con noso\[Lam]tros hoy
en \[Fa]la lucha \[Do]por hacer
de \[Fa]esta tierra un mundo de \[Sol7]paz.
\[Do]cuenta con noso\[Mim]tros hoy\[Lam]
\[Fa]ven y ayúda\[Mi]nos a ha\[Lam]cer
\[Fa]de esta causa \[Sol]nuestro ho\[Do]gar. (bis)
\endchorus
\endsong
\beginsong{Un Sí Generoso}[]
\beginchorus
\[Do]TÚ QUE DISTE UN SÍ\[Sol]
GENEROSO A \[Lam]DIOS;
\[Fa]TÚ QUE FUISTE LUZ
DE AYER PARA \[Sol]HOY;
\[Do]TÚ, LA HUELLA
\[Sol]QUE DEBEMOS SE\[Lam]GUIR,
TÚ, NUES\[Do]TRO FUNDA\[Fa]DOR:
GA\[Sol]BRIEL TABO\[Do]RIN. (bis)
\endchorus

\beginverse
1. \[Lam]Danos, Ga\[Fa]briel, es\[Do]peranza, prote\[Sol]cción;
\[Lam]danos, Ga\[Fa]briel, con\[Do]fianza en el a\[Sol]mor.
\endverse

\beginverse
2. Haznos, Gabriel, ser amigos del Señor;
haznos, Gabriel, ser testigos de ilusión.
\endverse
\endsong
\beginsong{Gracias Hermano}[]
    \beginverse
        \[Do]Gracias por tu \[La]vida,
        \[Fa]gracias por tu f\[Sol]e,
        \[Fa]por todo lo bueno que dejaste
        un d\[Sol]ía para q\[Do]ue,
        \[Fa]el mundo sea más b\[Do]ello
        \[Fa]un poco más fam\[Do]iliar
        \[Fa]y algo más \[Rem]humano y más fe\[Sol]liz.
    \endverse
    \beginverse
        Gracias por tu vida, 
        gracias por tu fe, 
        por escucharle en la oración 
        y hacer su voluntad 
        por no desfallecer, 
        por no mirar atrás 
        por buscar tu fuerza en el Señor.
    \endverse
    \beginchorus
        Y ahora \[Do]compart\[Sol]imos,
        Her\[Lam]mano Gabrie\[Mim]l
        un p\[Fa]oco del re\[Do]ino
        del re\[Fa]ino de Di\[Do]os.
        Que con tus m\[Do]anos vac\[Sol]ías
        y un \[Lam]poco de f\[Mim]e
        nos dej\[Fa]aste un camino
        hacia D\[Sol]ios.
    \endchorus
    \beginverse
        Ahora es nuestro turno hacer tu voluntad
        y desde esta casa Nazarena
        buscar la vocación, que hay mucho por hacer tú puedes ayudar:
        ¡Vamos todos juntos a sembrar!
        \[La7 Re]Vamos a disfrutar, prepárate:
        Uno, Dos y Tres:
    \endverse
\endsong
\beginsong{Amaos}[]
\beginchorus
\[Re]COMO EL P\[La7]ADRE ME A\[Sim]MÓ,
\[Sol]YO \[Mim]OS HE AM\[La7]ADO.
\[Re]PERMANE\[La7]CED EN MI A\[Sim]MOR,
\[Sol]PERMANE\[La7]CED EN MI A\[Re La7]MOR.
\endchorus
\beginverse
\[Sim]Si guardáis mis pa\[Re]labras,
\[Sol]y como her\[La7]manos os a\[Re]máis,
\[Fa#m]compartiréis con ale\[Sim]gría
el \[Sol]don de la \[Mim]fraterni\[La7]dad.
\endverse
\beginverse
Si os ponéis en camino,
sirviendo siempre la verdad,
fruto daréis en abundancia,
mi amor se manifestará.
\endverse
\beginverse
No veréis amor tan grande,
como aquel que os mostré.
Yo doy la vida por vosotros.
Amad como yo os amé.
\endverse
\beginverse
Si hacéis lo que os mando
y os queréis de corazón,
compartiréis mi pleno gozo,
de amar como Él me amó.
\endverse
   \endsong

\beginsong{En el taller}[]
\beginverse
\[Re]En el taller de Nazaret,
p\[La7]equeño y pobre \[Re]taller.
\[Re]En su labor está José,
y el \[La]Niño quiere apren\[Re]der.
\endverse
\beginchorus
LA\[Sol]BORA Y \[Re]CANTA
LA E\[La7]SPOSA DEL CARPIN\[Re]TERO.
Y EL \[Sol]MUNDO EN\[Re]TERO
SON\[]RÍE Y CANTA TAM\[Re]BIÉN. (bis)
\endchorus
\beginverse
En el taller de Nazaret,
pequeño y pobre taller.
Silencio y paz, amor y fe;
\endverse
        \endsong

\beginsong{Haz de tu casa}[]
\beginverse
\[Do]Subiré al des\[Fa]ván mis \[Do]odios y e\[Mim]goís\[Lam]mos,
\[Do]sacaré al jar\[Fa]dín esa \[Do]tonta discu\[Sol7]sión,
\[Fa]la mesa le pon\[Sol]dré, \[Do]hoy Je\[Mim]sús cena \[Lam]conmigo.
\[Do]Converti\[Sol]ré mi \[Lam]casa en Naza\[Fa]ret.
\endverse

\beginverse
A mi puerta llamarás, serás bien recibido,
dejaré que tú seas todo en este hogar,
una luz en mi interior
me dice que es posible.
Seré feliz viviendo en Nazaret.
\endverse

\beginchorus
\[Do]HAZ DE TU \[Fa]CASA UN NAZA\[Do]RET\[Fa]
\[Do]HAZ DE TU \[Fa]CASA UN NAZA\[Sol]RET
\[Fa]HAZ DE TU \[Sol]CASA, \[Do]HAZ DE TU CA\[Lam]SA,
\[Fa]HAZ DE TU \[Sol]CASA UN NAZA\[Do]RET
\endchorus

\beginverse
Si en tu \[Do]casa no hay a\[Sol]mor, ca\[Lam]lor ni fe;
si tu her\[Fa]mano poco o nada impor\[Sol]ta ya;
si tu \[Do]padre es un ex\[Lam]traño
y tu ma\[Fa]dre qué \[Fam]más da...
piensa a\[Do]migo y \[Fa]canta una vez \[Sol]más.
\endverse
\endsong
\beginsong{Padre Nuestro Nazaret}[]
\beginverse
\[Re]Padre de todos no\[Fa#m]sotros,
\[Sol]Tú que desde allí todo \[Re]lo ves,
manda sobre mi co\[Fa#m]razón
\[Sol]el espíritu de la \[La]fe.
\endverse

\beginverse
\[Re]Que mi boca se \[Fa#m]llene de tu nombre,
y \[Sol]mis manos de tu rei\[Re]nar.
Sea tu pala\[Fa#m]bra ley
\[Sol]y por \[La]ella respi\[Re]rar.
\endverse

\beginchorus
\[Re]Nazaret pueblo de los \[Fa]hombres
\[Sol]Nazaret cuna del amor \[Re]
\[Fa]Nazaret humilde y senci\[Sim]lla
\[Sol]Nazaret pue\[La]blo de \[Re (2ª Si)]dios
\endchorus

\beginverse
No me falte nunca alimento terrenal.
Y en la noche oscura no me dejes caer.
Hazme sensible ante el dolor.
Y al enemigo perdonar.
\endverse

\beginverse
Que tus brazos me puedan abarcar,
que tu fuente sea alivio de mi sed.
Tu Iglesia mi pasión,
mi deber continuar.
\endverse

\transpose{2}{
\beginchorus
\[Re]Nazaret pueblo de los \[Fa]hombres
\[Sol]Nazaret cuna del amor \[Re]
\[Fa]Nazaret humilde y senci\[Sim]lla
\[Sol]Nazaret pue\[La]blo de \[Re (2ª Si)]dios
\endchorus
}

\endsong
\beginsong{Al calor de la palabra}[]
	\beginverse
		\[Re]Cerca \[Do]del \[Sol]hogar
		que calienta mi \[Re]alma\[Do], \[Sol]
		quiero yo \[Re]saber lo que en comuni\[Sim]dad 
		Tú quieres de \[Mi]mí. \[La7]
	\endverse
	
	\beginverse
		Sintiendo el calor que me da tu Palabra 
		quiero yo saber a lo que me pides,
		sin que en nada yo pueda temer.
	\endverse

	\beginchorus
		A \[Sim]nada, a na\[Sol]da,
		nunca he de \[Re]temer, 
		yendo junto a \[Sol]ti
		con tus o\[La]jos de fe 
		nunca he de te\[Re]mer. (Bis)
	\endchorus

	\beginverse
		Sólo he de beber de tu fuente de agua
		sé que sólo ella será la que sacie 
		mi hambre y mi sed.
		Tú eres el Señor que alimenta mi alma
		y si hago la opción por seguirte a Ti
		nunca jamás yo temeré
	\endverse

	\beginchorus
		A nada…
	\endchorus

	\beginverse
		Llegan hasta mí momentos sin calma
		que me hacen dudar
		de si mi camino se orienta hacia ti.
		Comienza a faltar la paz en mi alma
		y sin esperarlo apareces Tú
		haciéndome ver que a nada he de temer.
	\endverse
\endsong

\beginsong{Ideal}[]
    \beginverse
        \[Do]Ideal de mi v\[Fa]ida abnegad\[Do]a
        \[Lam]de oració\[Do]n, de trab\[Fa]ajo y de amor\[Sol],
        ha de \[Fa]ser la fam\[Sol]ilia sagr\[Do]ada,
        \[Lam]que delFa \[Fa]hombre al dec\[Sol]hado alberg\[Sol7]ó.
        ha de \[Fa]ser la fam\[Sol]ilia sagr\[Do]ada,
        \[Lam]que del \[Fa]hombre al dec\[Sol]hado albergó\[Sol7].
    \endverse
    \beginchorus
        Nazar\[Fa]et, Nazare\[Sib]t, dulce \[Fa]nido
        de oraci\[Do]ón, de trabajo y de \[Fa]paz;
        vivir quiero a tu sombra tranquilo\[Sib],
        desde \[Fa]hoy tú serás \[Do]mi heredad\[Fa].
    \endchorus
\endsong
\beginsong{Madre}[]
\beginverse
\[Re]Junto a ti, Ma\[La]ría,
como un ni\[Sim]ño quiero est\[Fa#m]ar, 
\[Sol]tómame en tus \[Re]brazos, 
guía\[Mim]me en mi cami\[La]nar. 
\[Re]Quiero que me edu\[La]ques, 
que me ense\[Sim]ñes a re\[Fa#m]zar; 
\[Sol]hazme trans\[Re]parente, 
lléna\[Sol]me \[La]de \[Re]paz. 
\endverse

\beginchorus
\[Re]MA\[La]DRE, \[Sim]MADRE\[Fa#m],
\[Sol]MA\[Re]DRE, \[Mim]MA\[La]DRE. \[(2ª Mi La Re)](bis)
\endchorus

\beginverse
\[Re]Gracias, Madre \[La]mía, 
por lle\[Sim]varnos a Je\[Fa#m]sús.
\[Sol]Haznos más hu\[Re]mildes, 
tan senciMimllos como \[La]tú. 
\[Re]Gracias, Madre \[La]mía, 
por a\[Sim]brir tu cora\[Fa#m]zón, 
\[Sol]porque nos con\[Re]gregas 
y nos \[Sol]das \[La]tu a\[Re]mor. 
\endverse
\endsong
  \beginsong{Magnificat}[]
\beginverse
 \[Mim]Mi alma canta el \[Lam]gesto de amor.
Se \[Re]alegra mi espíritu en \[Sol]Dios Salvador,
pues \[Si]Él se fijó en \[Mim]la sencillez
y humil\[Do]dad de su \[Sol]sierva\[Si]. \[Mim]
\endverse
\beginchorus
SANTO ES \[Re]NUESTRO \[Do]DIOS
Y SU A\[Sol]MOR PERDURA\[Do]RÁ,
\[Do]Y LO RECORDA\[Re]RÁ
\[Si]CADA GENERAC\[Mim]IÓN.
\[Do]Y LO RECORDA\[Re]RÁ
\[Si]CADA GENERAC\[Mim]IÓN. \[Si] \[Mim]
\endchorus
\beginverse
Su brazo es fuerte y justo su obrar,
dispersa a los hombres de mal corazón;
derroca al poderoso, levanta al humilde,
llena al pobre de bienes.
\endverse
\beginverse
Ha protegido a su esclava el Señor.
A nuestros padres ya lo prometió
y se acordó de su amor
a Abraham por siempre.
\endverse

        \endsong

\beginsong{María la Madre Buena}[]
\beginverse
1. \[Sim]Tantas cosas en la vida
nos ofre\[Mim]cen plenitud,
y no \[La]son más que mentiras
que desgas\[Re]tan la inquietud.\[Fa#m]
Tú has \[Sim]llenado mi existencia
al que\[Mim]rerme de verdad.
Yo qui\[Sol]siera, Madre \[Mim]buena, amarte \[Fa#m]más.
\endverse

\beginverse
En silencio escuchabas
la Palabra de Jesús,
y la hacías pan de vida
meditando en tu interior.
La semilla que ha caído
ya germina y está en flor.
Con el corazón en fiesta cantaré.
\endverse

\beginchorus
\[Re]AVE MA\[Sol]RÍA, A\[Mim]VE MA\[Re]RÍA,\[La]
\[Re]AVE MA\[Sol]RÍA, A\[Mim]VE MA\[Re]RÍA.\[La Re]
\endchorus

\beginverse
2. Desde que yo era muy niño,
has estado junto a mí,
y, guiado de tu mano, aprendí a decir sí.
Al calor de la esperanza
nunca se enfrió mi fe,
y en la noche más oscura fuiste luz.
No me dejes, Madre mía,
ven conmigo al caminar.
Quiero compartir mi vida
y crear fraternidad.
Muchas cosas en nosotros
son el fruto de tu amor.
La plegaria más sencilla cantaré.
\endverse
\endsong
\beginsong{Quiero decir que sí}[]
\beginverse
1. \[Do]Quiero decir que \[Rem]sí, como tú Ma\[Sol7]ría,\[Do]
como \[Lam]tú un \[Rem]día, como \[Sol7]tú Ma\[Do7]ría.
\endverse

\beginchorus
\[Do7]QUIERO DECIR QUE \[Fa]SÍ,
\[Sol7]QUIERO DECIR QUE \[Mim]SÍ,
\[La7]QUIERO DECIR QUE \[Rem]SÍ,
\[Sol]QUIERO DECIR QUE \[Do]SÍ.
\endchorus

\beginverse
2. Quiero entregarme a Él, como tú María,
como tú un día, como tú María.
\endverse

\beginchorus
QUIERO ENTREGARME A ÉL (4).
\endchorus

\beginverse
3. Quiero alabarle a Él, como tú María,
como tú un día, como tú María.
\endverse

\beginchorus
QUIERO ALABARLE A ÉL... (4)
\endchorus

\beginverse
4. Quiero ser de Jesús, como tú María,
como tú un día, como tú María.
\endverse

\beginchorus
QUIERO SER DE JESÚS... (4).
\endchorus
\endsong
\beginsong{A ti que estás sentado}[]
            
            \beginverse
            \[Re]A ti que estás se\[La]ntado en tu si\[Sim]llón,
            \[Sol]atento, dist\[Mim]raído, te rec\[La]lamo;
            \[Sol]quiero que hagamos\[La] vida
            esta can\[Re Fa#m Sim]ción,
            \[Sol]y a Dios vayamos jun\[Mim]tos de la m\[La]ano.
            \endverse
            \beginverse
            Hasta ahora sólo hablamos en teoría
            de un mundo que sería más cristiano,
            de un mundo que será
            sólo teoría si tú y yo,
            ahora no hacemos nada por cambiarlo.
            \endverse
            \beginchorus
             \[Sol]De poco servi\[La]rán nuestros esfuer\[Re...7]zos,
             \[Sol]si en la calle no ha\[La]cemos que se no\[Re...7]ten
             \[Sol]las palabras que lan\[La]zamos 
             a los v\[Re Fa#m Sim]ientos;
             \[Sol]que Cristo sigue \[Mim]vivo en cada hom\[La]bre.
            \endchorus
            \beginverse
            Ya no valdrán ni «peros» ni disculpas,
            pues Cristo no es alguien con quien se juega.
            Si este mundo no cambia
            solamente será culpa
            de quien diciendo «sí» luego reniega.
            \endverse
            \beginverse
            A tí que estás sentado en tu sillón,
            quisiera que no olvides este rato.
            Dios ha pasado
            por este rincón,
            y quiere que hoy vayamos de la mano.
            \endverse
        \endsong
\beginsong{Tú eres del Señor}[]
    \beginverse
        \[Re]Eres tan sencilla
        como \[La]luz de amanecer.
        \[Sim]Eres tú, María, forta\[Fa#m]leza de mi fe.
        \[Sol]Tú eres flor, \[Re]eres del Señor,
        te \[Mim]dejas acariciar por su am\[La]or.
    \endverse
    \beginverse
        Eres tan humilde
        como el vuelo del gorrión,
        eres tú, María, el regazo del amor.
        Tú eres flor, eres del Señor,
        te dejas acariciar por su amor.
    \endverse
    \beginchorus
        \[Sim]YO QUIERO ESTA\[Sol Mim]R
        EN LAS M\[Sol]ANOS DEL SEÑ\[Re]OR
        COMO \[Sim]TÚ.
        \[Sim]PARA AM\[Sol Mim]AR,
        EN LAS \[Sol]MANOS DEL SEÑ\[Re]OR
        COMO \[Sim]TÚ. COMO \[La]TÚ, COMO \[Re]TÚ.
    \endchorus
    \beginverse
        Eres tan pequeña como el canto de mi voz,
        eres la grandeza de aquel que te modeló.
        Tú eres flor, eres del Señor,
        te dejas acariciar por su amor.
        Eres tan hermosa como el cielo,
        como el mar,
        eres tú, María, como el gozo de soñar.
        Tú eres flor, eres del Señor,
        te dejas acariciar por su amor.
    \endverse
\endsong
\beginsong{La Virgen sueña caminos}[]
\beginverse
\[Mi]La Vir\[La]gen sueña ca\[Mi]minos, está a la es\[Si7]pera;
la \[La]Virgen sabe que el \[Mi]niño, es\[Si7]tá muy cer\[Mi]ca.
de Nazaret a Belén, hay una \[Si7]senda;
por \[La]ella van los que \[Mi]creen en \[Si7]las prome\[Mi]sas.
\endverse

\beginverse
\[Mi]Los que soñáis y esperáis la Buena Nue\[Si7]va,
A\[La]brid las puertas al \[Mi]niño que es\[Si7]tá muy cer\[Mi]ca.
El Señor \[La]cerca es\[Mi]tá 
Él \[Si7]viene con la \[Mi]paz.
El Señor \[la]cerca es\[Mi]tá 
Él \[Si7]trae la ver\[Mi]dad.
\endverse

\beginverse
En estos días del año
El, el pueblo espera que
Venga pronto el Mesías
A nuestra tierra.
En la ciudad de Belén 
Llama a las puertas,
Pregunta en las posadas y no hay respuesta.
\endverse

\beginverse
La tarde ya lo sospecha: está alerta.
El sol le dice a la luna que no se duerma.
A la ciudad de Belén
Vendrá una estrella,
Vendrá con todo el que quiera
Cruzar fronteras.
\endverse
\endsong
\beginsong{Creo en ti}[]
\beginverse
\[Do]No recuerdo ni \[Fa]cuando te encon\[Do]tré,
Has entrado en mi \[Fa]vida sin que\[Sol]rer,
\[Fa]Eres como el \[Do]viento, \[Fa]que no po\[Sol]demos\[Do]
\[Fa]Nada \[Rem]te puede deten\[Sol]er.
\endverse

\beginverse
Aquí estoy, poco te puedo ofrecer.
Sólo un canto, una gota de ilusión.
como una flor mi vida se abre al sol,
va creciendo, regada por tu amor.
\endverse

\beginchorus
\[Do Sol Lam]CREO EN TI 
\[Fa]AUNQUE SIENTA \[Sol]CANSANCIO
CREO EN\[Do...7] TI 
\[Fa]CREO EN TU \[Sol]MIRADA,
EN TU \[Mi]FUERZA Y EN TU\[Lam] LUZ
CREO EN TI JESÚS.

CREO EN TU MIRADA ME LLAMA
Y QUIERO IR. AUNQUE NO TE
SIENTA A VECES JUNTO A MÍ,
YO…, YO CREO EN TI.
\endchorus

\beginverse
Yo no entiendo la muerte ni el dolor,
Tú me pides un nuevo corazón.
Ayúdame a esperar, una vez más,
que la mañana venza a la oscuridad.
\endverse

\beginverse
De la roca ha brotado un manantial, 
eres agua que corre en libertad,
Cristo del arroyo, brillando bajo el sol,
nuevo día, universo de amor. Amén.
\endverse
\endsong

        \beginsong{Maravillas}[]
\beginverse
\[Do]Maravillas hizo en \[Sol]mí
mi \[Fa]alma canta \[Sol]de gozo
pues \[Lam]al ver mi \[Mim]pequeñez
se \[Fa]detuvieron sus o\[Sol]jos.
Y el que \[Lam]es Santo \[Mim]y poderoso
\[Fa]hoy aguarda p\[Sol]or mi sí.
Mi \[Do]alma canta \[Sol]de gozo
ma\[Fa]ravillas \[Sol]hizo en \[Do]mí.
\endverse
\beginverse
Maravillas hizo en mí
del alma brota mi canto.
El Señor me ha amado
más que los lirios del campo.
Por el Espíritu Santo
Él habita hoy en mí.
No cese nunca mi canto,
maravillas hizo en mí.
\endverse

        \endsong

\beginsong{María mírame}[]
\beginverse
\[Do]María mí\[Sol]rame, Ma\[Lam]ría míra\[Fa]me 
si tú me \[Do]miras, Él tam\[Fa]bién me mira\[Sol]rá 
Madre mía \[Lam]mírame\[Fa], de la mano \[Do]llévame 
\[Fa]muy cerca de \[Do]Él, que ahí \[Fa]me quie\[Sol]ro
que\[Do]dar.
\endverse

\beginchorus
\[Do]María cúbre\[Sol]me con tu man\[Lam]to, 
que tengo mie\[Fa]do, no \[Sol]sé rezar,\[Do] 
que por tus oj\[Fa]os mi\[Sol]sericor\[Lam]diosos 
tendré la fuerza, \[Fa]tendré \[Sol]la paz.\[Do]
\endchorus

\beginverse
Madre, consuélame de mis penas, 
que no quiero ofenderle más, 
que por tus ojos misericordiosos 
quiero ir al cielo y verlos ya.
\endverse
\endsong
\beginsong{AVE MARÍA}[by={Maite López}]
\beginchorus
\[Re]Ave Ma\[La]ría, \[Sol]Ave Ma\[Re]ría\[La]. 
\[Re]Ave \[Sim]María\[Fa#m Sol La].
\endchorus

\beginverse
\[Re]Dios te salve María, \[Fa#m]llena eres de gracia 
\[Sol]el Señor está con\[La]tigo, bendita tú \[Re]eres 
entre \[Fa#m]todas las mujeres 
y ben\[Sol]dito, bendito es el fruto 
de tu \[La]vientre, Jesús. 
\endverse

\beginchorus
Estribillo
\endchorus

\beginverse 
\[Re]Yo me siento tan frágil 
¿\[Fa#m]Cómo vienes a mí? ¿\[Sol]Cómo buscas mi pobre\[La]za? 
Me aban\[Re]dono en tus manos. \[Fa#m]Hágase tu voluntad 
y pro\[Sol]clama, proclama mi alma 
Tu gran\[La]deza, oh, mi Señor.
\endverse
\endsong
\beginsong{Deja que en tu vida cante Él}[]
    \beginverse
        \[Fa]Deja que en tu vida cante Él,
        deja que en tu vida baile Él\[Do],
        \[Sib]y que por tus \[Do]venas s\[Fa]ólo corra Él\[Rem].
        \[Fa]Deja que en t\[Do]u vida cante É\[Sib]l.
    \endverse
    \beginverse
        Ábrele las puertas a Jesús, y cierra los 
        ojos al temor, déjate en sus manos, sólo
        eres de Él. Ábrele las puertas a Jesús.
    \endverse
    \beginchorus
        \[Fa Do]Jesús será
        \[Sib]La fuerza p\[Do]ara amar
        \[Fa Do]Jesús dará \[Sib]sentido a\[Do] tu vivir
        \[Fa]Él será el c\[Do]entro
        y e\[Sib]l eje de tu ser,
        \[Fa]Él cambiará tu exi\[Do]stir.
        \[Fa]Él, siempr\[Do]e joven,
        va\[Sib]lor te exij\[Fa]irá
        \[Fa]para ser \[Do]joven con \[Sib]Él.
    \endchorus
    \beginverse
        Haz del Evangelio una canción, que te 
        acerque a Dios en los demás, cántala con
        fuerza y de corazón. Haz del Evangelio
        una canción.
    \endverse
    \beginchorus
        Sembrarás palabras de verdad que 
        florecerán en torno a ti, llevarán el gozo,
        llevarán amor.
        Sembrarás palabras de verdad.
    \endverse
\endsong
\beginsong{Háblame}[]
\beginverse
\[Re]Yo siento, Se\[Sim]ñor, \[Sol]que tú me \[La7]amas.
Yo \[Re]siento, Se\[Sim]ñor, \[Sol]que te pued\[La7]o amar.
\[Re]Háblame, Se\[Sim]ñor, \[Sol]que tu siervo es\[La7]cucha.
\[Re]Hábla\[Sim]me, ¿\[Sol]qué quieres de \[La7]mí?
\endverse
\beginchorus
\[Re]SEÑOR,
Mim\[]TÚ HAS SIDO GR\[Fa#m]ANDE PARA M\[Sol]Í,
EN EL DESIERTO DE MI \[Sol]VIDA:
¡HÁBLA\[La7]ME!
\[Re]YO QUIERO ES\[Fa#m]TAR 
\[Sol]DISPUESTO A \[La7]TODO,
\[Re]TOMA MI \[Fa#m]SER,
\[Sol]MI CORAZÓN ES PA\[La7]RA TI.
\[Re]POR ESO CANTO \[Fa#m]TUS MARAVILLAS,
\[Sol]POR ESO CANTO TU \[La7]AMOR.
\[Re]POR ESO CANTO \[Fa#m]TUS MARAVILLAS,
\[Sol]POR ESO CANTO TU A\[La7]MOR 
LA, RA, LA, ….
\endchorus
\beginverse
Te alabo, Jesús, por tu grandeza.
Mil gracias te doy por tu gran amor.
Heme aquí, Señor, para acompañarte.
Heme aquí, ¿qué quieres de mí?
\endverse
   \endsong

\beginsong{Jesús te quiere}[]
	\beginverse
		\[Sol]Hermano, tú que buscas
		la \[Sim]paz en tu interior,
		ol\[Do]vida tus problemas 
		y \[Re7]abre tu corazón.
	\endverse

	\beginverse
		Jesús busca tu vida,
		no importa cómo estés,
		cansado o deprimido,
		Hoy puedes renacer.
	\endverse

	\beginchorus
		PORQUE \[Sol]JESÚS TE QUIERE,
		NO IM\[Mim]PORTA COMO SEAS,
		NO IM\[Do]PORTA TU POBREZA,
		JE\[Re]SÚS TE QUIERE \[Re7]ASÍ.
		PORQUE \[Sol]JESÚS
		ESTÁ LOCO DE AMOR \[Mim]POR TI,
		HAZLE UN SITIO EN TU CO\[Do]RAZÓN,
		DÉJALE SER TU \[Re]SEÑOR. (BIS)
	\endchorus

	\beginverse
		Jesús está buscando 
		a alguien como tú,
		dispuesto a darlo todo,
		a amar y ser feliz.
		Él te dará la fuerza, Él te dará la paz,
		y luchara contigo,
		no te abandonará.
	\endverse
\endsong

        \beginsong{Le llamaron loco}[]
            
            \beginverse
            \[Re]Le llamaron \[La]loco y tú te at\[Sim]reves,
            a dar un paso al fr\[Sol Mim]ente de esta locu\[La]ra.
            Le llamaron loco
            y tú te entregas simplemente
            ante un trozo de pan y un poco de vino
            \endverse
            \beginchorus
            Y TU LOCURA ES Q\[Re]UERER
            \[La]AMAR SIN M\[Sim]EDIDA
            AL QUE NADA T\[Sol Mim]IENE,
            AL QUE NADA ES\[La]PERA.
            Y TU LOCURA ES S\[Re]EGUIR
            \[La]A AQUEL QUE DIO SU \[Sim]VIDA
            EN UN TROZO D\[Sol Mim]E PAN
            Y UN POCO DE V\[La]INO.
            \endchorus
            \beginverse
            Le llamaron loco y tú te atreves
            a no mirar atrás pese a todo
            Le llamaron loco
            y algún día oirás también tú
            decir a otros ese mismo nombre.
            \endverse
            \beginverse
            \[Re]Le llamaron \[La]loco y tú te \[Sim]atreves,
            a dar un paso al \[Sol Mim]frente de esta lo\[La]cura.
            de am\[Re]or
            \endverse
        \endsong
\beginsong{Necesito de ti}[]
\beginverse
\[Sol]Hoy, Señor, me he acercado a ti 
\[Mim]y me has dado la mano.
\[Do]Y mecido entre tus brazos vi,
\[Re]que estas a mi lado.
\endverse

\beginverse
\[Sol]Eres todo lo que yo quiero,
\[Mim]todo lo que yo ansío,
\[Do]eres vida y esperanza, co\[Re]razón y camino.
\endverse

\beginchorus
\[Do]NECESITO DE\[Re] TI, 
\[Sol]SÓLO VIVO\[Re] EN\[Mim] TI.
\[Do]YA NO EXISTE \[Re]NADA 
QUE ME A\[Sol]PAR\[Re]TE DE \[Mim]TI.
PUES \[Do]NECESITO \[Re]DE TI,
SÓLO DE\[Sol] TI.
\endchorus

\beginverse
Al mirarte he podido sentir 
que eres tú el que nos ama.
Ahora sé que no es una ilusión
No eran sólo palabras.
\endverse

\beginverse
Ahora entiendo tu mensaje,
Ahora se que no es un sueño,
Ahora solo quiero darte
Todo lo que yo siento.
\endverse
\endsong

\beginsong{El agua del Señor}[]
	\beginchorus
		\[Re]EL AGUA DEL SEÑOR
		SA\[Si]NÓ MI ENFERMEDAD 
		E\[Sol La7 La Re]L AGUA DEL SEÑOR JESÚS (Bis) 
		E\[Sol La Re...7]L AGUA DEL SEÑOR JESÚS 
	\endchorus

	\beginverse
		\[Sol]El\[Re]que quiera y tenga sed
		\[La7]que \[Re...7]venga y beba gratis 
		\[Sol]El\[Re]que quiera y tenga sed
		\[La7]beba \[Re]el agua de la vida 
	\endverse

	\beginchorus
		EL AGUA DEL SEÑOR 
		SANÓ MI ENFERMEDAD
		EL AGUA DEL SEÑOR JESÚS (Bis) 
	\endchorus

	\beginverse
		El que beba de este agua
		jamás tendrá sed. 
		El que beba de este agua 
		jamás tendrá sed. 
	\endverse

	\beginverse
		Sobre ti derramaré
		el agua que es mi vida,
		y tu corazón de piedra 
		en amor transformaré.
	\endverse

	\beginverse
		El que crea en mi Palabra, 
		y se abra a mi fuerza, 
		de su seno brotarán
		torrentes de agua viva. 
	\endverse

	\beginverse
		Te doy gracias, tú me cambias, tú me llenas, me has salvado. (bis)
	\endverse
\endsong
\beginsong{Siento que Tú eres Dios}[]
    \beginverse
        \[Sol]Siento que tú eres Di\[Mim]os,
        que tus ojos me m\[Do]iran con tanto amor\[Re7],
        y que puedo senti\[Sim]r
        el gozo de tu prese\[Mim]ncia,
        que está d\[Do]entro \[Lam]de mi corazón\[Re7].
        \[Sol]Cantaré para t\[Mim]i
        porque a tu la\[Do]do por siempre estaré\[Re7];
        camina\[Sim]ré siempre contig\[Mim]o,
        me goza\[Do Lam]ré en tu verda\[Re7]d.
    \endverse
\endsong
\beginsong{Tan solo he venido}[]
\beginverse
\[Re]No he ve\[La]nido a pedir\[Sol]te
\[Re]como s\[La]uelo, Seño\[Sol]r
\[Re]si antes \[La]de yo clama\[Sol]rte
co\[Mim]noces mi peti\[La]ción
\endverse
\beginverse
Solo quiero escucharte 
pon el tema, Señor
caminar por el parque 
y dedicarte una canción 
\endverse
\beginchorus
TAN SÓLO HE V\[Re]ENIDO
A E\[La]STAR \[Sim]CONTIGO
A \[Fa#m]SER TU \[Sol]AMIGO
A COMPAR\[Mim]TIR CON \[La]MI DIOS
ADO\[Re]RARTE Y\[La] DARTE \[Sim]GRACIAS
POR \[Fa#m]SIEMPRE \[Sol]GRACIAS
POR LO QUE \[Mim]HAS HECH\[La]O, SEÑO\[Re]R CONMIGO
\endchorus
\beginverse
Cuéntame de tus obras
¿Qué hay de nuevo, Señor?
Y de paso pregunto
¿Cómo es la luz del sol?
\endverse
\beginverse
Y yo, sólo quiero abrazarte
bendecirte mi Dios
caminar por las calles
y abrirte mi corazón
\endverse
   \endsong

    \beginsong{Color esperanza}[]
        \beginverse
\[Re]Sé que hay en tus o\[La]jos con \[Fa#]solo mi\[Sol]rar
que estás \[Lam]cansado \[Re7]de andar y de an\[Sol]dar
y cami\[Sim]nar girando \[Mim]siempre en un \[La]lugar

Sé que las ventanas se pueden abrir,
cambiar el aire depende de ti, te ayudara
vale la pena una vez más
\endverse
\beginchorus
Saber que se p\[Re]uede \[La]querer que se \[Sim]pueda 
quitarse los mie\[Mi]dos sacarlos afu\[La]era
pintarse la ca\[Re]ra co\[Fa#]lor espe\[Sim]ranza
tentar al fu\[Sol]turo con el corazón\[La]. \[Re]
\endchorus
\beginverse
Es mejor perderse que nunca embarcar,
Mejor tentarse a dejar de intentar,
Aunque ya ves que no están fácil 
Empezar. Sé que lo imposible se puede
Lograr, que la tristeza algún día se irá y 
Así será la vida cambia y cambiará

\[Sim]Sentirás que \[Fa#]el alma vuela por \[Sol]cantar u\[La]na 
Vez más.
\endverse
\beginchorus
Saber que se puede querer que se pueda…
\endchorus
\beginverse
\[Sim]Vale más \[Mim]poder \[Sim]brillar que \[Sim]solo \[Sol]buscar ver el sol.
\endverse
\beginchorus
Saber que se puede querer que se pueda…
\endchorus
        \endsong

        \beginsong{Sigues llamando}[]
            
            \beginverse
            \[Mim]Te conozco de hace ti\[Re]empo
            sin em\[Do]bargo no rec\[Re]uerdo
            \[Mim]que nos hayan presen\[Re]tado
            en el pasad\[Do Re]o
            \endverse
            \beginverse
            Y aunque siempre te he sentido tan
            cercano y tan amigo sé que lo que ahora
            me pides es distinto
            \endverse
            \beginverse
            \[Sol]Y hoy vienes a \[Re]mi, y mi\[Mim]rándome
            me i\[Re]nvitas a empe\[Do]zar
            a en\[Re]tregarme a los de\[Do]más, a \[Re]amar
            Y hoy estoy aquí, escuchándote
            asustado sin saber
            si seré capaz de serte fi\[Do Re]el.
            \endverse
            \beginchorus
            QUI\[Sol]SIERA SER \[Re]VOZ DE TU \[Mim]MÚSICA 
            YA ES\[Do]CRITA
            EN TU \[Sol]AGUA BE\[Re]BER, Y SA\[Mim]CIAR MI SED 
            DE \[Do]AMOR
            QUI\[Sol]SIERA SER \[Re]LUZ
            NO UNA \[Mim]CHISPA QUE SE A\[Do]PAGA
            Y T\[Sol]OMAR DE TU \[Re]CRUZ,
            ESA \[Mim]FUERZA QUE HOY ME \[Do]FALTA
            PARA SEGU\[Re]IR
            \endchorus
            \beginverse
            Me conoces y ya sabes
            mis excusas y mis dudas
            y aún así sigues llamando al corazón
            \endverse
            \beginverse
            Hoy Señor quiero arriesgarme,
            aunque falle en el camino,
            confiarme en tus brazos
            como un niño
            \endverse
            
            
        \endsong
\beginsong{Evangelio}[]
\beginverse
Evan\[Sol]gelio es decir amigo\[Sim], es\[Fam] decir her\[Mim]mano.
Evan\[Sim]gelio es darte mi tiem\[Do]po, es dar\[Re7]te mi mano.
Evan\[Sol]gelio es mirar\[Sim]te a los\[Sol7] ojos, es reír\[Do] contigo.
\endverse
\beginchorus
Es compar\[Re]tir tu \[Sol]pena, es lle\[Re]varte a Cris\[Sol]to (bis)
\endchorus

\beginverse
Evangelio es llevar la paz siempre contigo.
Evangelio es amar de balde hasta caer rendido.
Evangelio es decir “te amo” a tu enemigo.
\endverse
\beginchorus
Abandonar tu vida en manos de Cristo(bis)
\endchorus

\beginverse
Evangelio es vivir como un pobre que todo lo espera.
Evangelio es mirar al cielo con ojos de niño.
Evangelio es dar gracias al Padre al nacer el día.
\endverse
\beginchorus
Y continuar cantando al hacer camino. (bis)
\endchorus

\beginverse
Evangelio es sembrar libertad, es vivir unidos.
Es llevar la esperanza a un mundo que llora perdido.
Evangelio es romper cadenas, es abrir sepulcros.
\endverse
\beginchorus
No \[Re]lo bus\[Sol]quéis muerto que \[Re]está entre \[Sol]los vivos (4 veces) 
\endchorus
\endsong

\beginsong{Baja a Dios}[]
    \beginverse
        \[Do]Baja a Dios de las nubes,
        llévale a la fáb\[Mim]rica donde trabajas
        \[Do]quita a Dios del retablo
        y grábale \[Re]dentro de tu corazón.
        Y roba, roba, roba
        \[Sol]roba a dios de los templos
        donde lo encerra\[Mim]mos hace tantos años;
        \[Do]déjale libre en las plazas,
        y llévale tambi\[Re]én al mercado del pueblo.
    \endverse
    \beginchorus
        Porque \[Do]Dios
        No es un Dios muerto
        y si\[Sol] pensais que es\[Mim]tá muerto
        \[Do]equivo\[Sol]cados,\[Do]equi\[Sol]vocados
        eq\[Do Re]uivocados e\[Sol7]stáis.
        Porque \[Do]Dios
        no es un Dios m\[Dom]uerto
        y si\[Sol] pensáis que es\[Mim]tá muerto
        \[Do]equivo\[Sol]cados,\[Do]equi\[Sol]vocados
        eq\[Do Re]uivocados e\[Sol7]stáis.
    \endchorus
\endsong
\beginsong{Como corre un río}[]
\beginchorus
\[Re]COMO CORRE \[Sol]UN RÍO \[La]DENTRO DE \[Re]MI SER
COMO CORRE \[Sol]UN R\[La]ÍO DENTRO DE \[Re]MI SER
ASÍ YO \[Sol]CONFIO EN \[La]CRISTO \[Re]MI REY
ASÍ YO \[Sol]CONFIO EN \[La]CRISTO \[Re]MI REY
COMO UN RIO DE AGU\[Sol]A VIVA,
QUE SALTA \[La]PA’ ARRIBA, QUE LLE\[Re]VO DENTRO
\[Fa#m]CONFIRM\[Sim]A, CONFIRMA, EN ESTE M\[Sol]OMENTO 
EL ESPÍRI\[La]TU SANTO, EN SU DERRA\[Re]MAMIENTO
\endchorus
\beginverse
Jesús \[Sol]está pa\[La]sando por \[Re7]aquí
Jesús \[Sol]está pa\[La]sando por \[Re7]aquí
Por donde \[Sol]pasa todo lo transfor\[Re]ma
Se va la trist\[Sol]eza, \[La]llega la aleg\[Re7]ría
Por donde \[Sol]pasa t\[La]odo lo transfo\[Re]rma
Llega la aleg\[Sol]ría, para \[La]ti y para m\[Re]í
\endverse
   \endsong

\beginsong{Cristo está presente en tu corazón}[]
        \beginverse
\[Mi]Cristo está presente en tu corazón                            
\[La]Uoo,uoo  
\[Mi]Cristo está presente en tu corazón 
Uooo  
\[Mi]Grita JESÚS
\[Mi]más fuerte, JESÚS
\[Mi]di que le amas, LE AMO
Cristo está presen\[La]te en tu \[Si]corazón \[Mi]

Abre tu corazón a su vida y amor
Uoo,uoo.
Abre tu corazón a su vida y amor
Uoo,uoo.

Grita JESÚS
más fuerte, JESÚS
di que le amas, LE AMO
Cristo está presente en tu corazón
\endverse
\beginchorus
\[Mi]Levanta las manos,
                     La
mueve la ca\[La]beza, la cintura
                       Si
salta para un \[Si]lado, para el otro,
                           Mi Sim
para alabar a D\[Mi]ios(\[Sim]bis)

\[Mi]Chup-chup-chup
churu-churu-\[La]chup-chup-chup
churu-churu-\[Si]chup-chup-chup
churu-churu-\[Mi]chup
\endchorus
        \endsong

   \beginsong{Hay tantos pies que lavar}[]
            \beginchorus
            \[Do]Sigue habiendo \[Sol]tantos pies que la\[Lam Sol]var
            \[Do]sigue habiendo \[Sol]tanta oscuri\[Lam]dad que ilum\[Sol]inar
            \[Fa]tantas cade\[Sol]nas que romp\[Do Lam]er
            pan y \[Rem]vino para el p\[Fa]obre quiero \[Sol]ser
            forta\[Rem]lece, S\[Fa]eñor, mi poca \[Sol]fe.
            \endchorus
        \endsong
\beginsong{Dependo de ti}[]
\beginverse
\[Do]Dependo de ti, necesito de ti
Como el pez\[Lam] necesita del agua.
\[Fa]Dependo de ti, impres\[Rem]cindible para mí
Como nece\[Sol]sito la comida diaria\[Sol]. (bis)
\endverse

\beginchorus
\[Do]Señor, si no estás con\[Fa]migo,
mi Dios, me siento ven\[Fa]cido,
Jesús\[Rem], me siento perdi\[Sol]do. (bis)
Por eso \[Do]hoy yo le canto, le amo, le alabo
con \[Lam]todas mis fuerzas.
He comprendi\[Sol]do que Él es la energía,
la fuerza div\[Do]ina que me alimenta.
Por eso \[Lam]hoy quiero gritarle 
con lágrimas en los o\[Fa]jos
que estoy enamorado de Él
y que no \[Sol]puedo
pues dependo de \[Do]Él.
\endchorus

\beginverse
Dependo de ti necesito de ti
Como el aire que respiro.
Dependo de ti, imprescindible para mí
Como el niño necesita del cariño.
\endverse
\endsong

\beginsong{Enamorado}[]
    \beginverse
        \[Sol]Aprendí a viv\[Re]ir a tu manera\[Do Re]
        \[Sol]hoy hablo de t\[Re]i por donde quieras\[Do Re]
        \[Mim]Tú cambiaste el \[Re]mundo \[Do]en el cual vivía\[Sol]
        \[Do]me diste la \[Lam]paz que no tení\[Re]a.
    \endverse
    \beginverse
        Cuando estoy contigo estoy seguro
        tú me das la fuerza cuando dudo
        y si el sol se oculta, no me desespero
        yo no he de temer, porque en ti espero.
    \endverse
    \beginchorus
        \[Sol]Enamorad\[Re]o, en\[Mim]amorad\[Sim]o
        \[Do]de mi cristo vivo \[Lam]enamorad\[Re]o.
        \[Sol]Enamorad\[Re]o, en\[Mim]amorad\[Sim]o
        \[Do]tu palabra viva \[Re]me ha salv\[Sol]ado
    \endchorus
    \beginverse
        Quiero agradecerte aquel momento 
        En que te entregaste a ese tormento
        Cuando te morías, tú por mi lo hacías
        Pero regresaste al tercer día.
    \endverse
\endsong
  \beginsong{Misionero serás}[]
\beginverse
\[Do]Dice la gente
que este \[Lam]mundo es muy raro,
que lo que \[Fa]hay que hacer es pasa\[Sol]r,
\[Do]que son muchos los que \[Lam]tienen problemas
y por qué \[Fa]los vas a ayudar, \[Sol]
\[Fa]pero no les hagas \[Sol]caso,
pues en \[Do]ellos no está \[Mim]la verda\[Lam]d,
\[Fa]ese don que ig\[Sol]noran,
en \[Do]ti lo ha\[Mim]n de encontrar, \[La,]
\[Fa]si lo descubr\[Sol]es,
misionero ser\[Do7]ás.
\endverse
\beginchorus
MISIONERO SE\[Fa]RÁS,
MISIONERO SE\[Sol]RÁS,
MISIONERO DE \[Do]TU CIU\[Lam]DAD,
MISIONERO QUE \[Fa]BUSCA
LA ALE\[Sol]GRÍA EN A\[Do7]YUDAR.
\endchorus
\beginchorus
MISIONERO SERÁS,
MISIONERO SERÁS,
MISIONERO DE TU CIUDAD
MISIONERO QUE OFRECE
TODA SU VIDA A LOS DEMÁS
\endchorus
\beginverse
Creen muchos que ser misionero
es solamente viajar 
A todos los pueblos 
que a Dios no conocen
y a Él quisieran llegar,
puedes ser misionero en tu propia ciudad,
Dios te ayudará
si lo consigues, misionero serás
\endverse
   \endsong

\beginsong{Jesús está en ti}[]
    \beginchorus
        \[Re]JESÚS ESTÁ \[Sim]EN TI, ES\[Sol]TÁ EN MÍ,
        ESTÁ EN T\[La]ODOS MIS HERMANOS\[RE].
        ESTÁ E\[Sim]N TI, ESTÁ E\[Sol]N Mí
        ESTÁ EN\[La] TODOS MIS HERM\[Re...7]ANOS
    \endchorus
    \beginverse
        Quien h\[Sol]ace crecer las \[La]flores, 
        quien \[Re]hace des\[Fa#m]cender el a\[Sim]gua,  
        quien \[Sol]hace nuevo tu corazón \[La Re]cada mañana. \[Re7]
    \endverse
    \beginverse
        Quien hace amanecer el sol, 
        quien llena de estrellas el cielo, 
        quien en la noche más oscura 
        te da todo su consuelo
    \endverse
\endsong
\beginsong{Hoy solo te pido}[]
\beginverse
\[Re]Hoy sólo te \[Fa#m]pido, \[Sol]que me des tu a\[La]liento,
\[Re]que me des tu \[Fa#m]mano para \[Sol]amar,
nada \[La]más.
\endverse
\beginverse
Para ver más claro, seguir noche y día,
no temer a nada para andar,
nada más.
\endverse
\beginchorus
DESDE \[Re]HOY QUIERO SEGUIR\[Sol]TE,
DESDE \[Re]HOY QUIERO ALA\[Sol]BARTE,
DESDE \[Re]HOY QUIERO ANUN\[Sol]CIAR
YA TU PALA\[La]BRA.
\endchorus
\beginverse
Hoy quiero vivirte, sentir tu mirada, que Tú
me acompañes al andar nada más.
Quiero hacer camino, desde tu Palabra,
salir de mi tierra y caminar, nada más...
\endverse
\beginchorus
DESDE HOY ...
\endchorus
\beginverse
Siento un nuevo aviso, ¡basta ya de excusas!
hazme despertar a tu verdad,
una vez más.
\endverse

        \endsong

        \beginsong{Yo no puedo solo}[]
            \beginverse
            \[Re]¡Qué difícil \[Sim Sol]es andar en so\[La]ledad...!
            \[Re]¡Qué difícil \[Sim]es crece\[Sol]r cuando no \[La]estás...!
            \[Re]Comprendo que hay m\[Sim]omentos
            en l\[Sol]os que he sido in\[La]fiel.
            \[Re]Hoy te pido: a\[Sim]cércate.
            \[Sol]Tenemo\[La]s mucho de que ha\[Re]blar.
            \endverse
            \beginchorus
            \[Sol]SI ME FAL\[La Re]TAS TU
            NO SABRÉ QUE HA\[Sim]CER.
            \[Sol]SIN TI MI EXIS\[La]TENCIA
            NO ES IGU\[Re...7]AL.
            YO PROCURARÉ
            SIEMPRE CAMINAR
            DE FORMA QUE AUMENTE
            TU AMISTAD.
            \endchorus
            \beginverse
            A tu lado alcanzaré antes mi afán,
            yendo juntos no caeré, Tú me tendrás.
            Será mucho más fácil si vienes junto a mí
            Paso a paso y día a día lejos hemos de llegar.
            \endverse
            \beginchorus
            Estribillo...
            \endchorus
            \beginverse
            Hay tantas cosas que hacer, tanto que dar,
            que me asusta no saber cómo empezar. Si
            tu marcas la ruta sin miedo avanzaré.
            Tan seguro estoy de ti, que nada
            me podrá parar.
            \endverse
        \endsong
\beginsong{YE-YE-YE}[]

\ifchorded
\beginverse
\memorize
\nolyrics Todo el rato: \[Re] \[La] \[Sol] \[La]
\endverse
\fi

\beginchorus
\[Re]YE-YE-YE 
\[La]YE-YE-YE-YE  
\[Sol]YE-YE-YE 
\[La]YE-YE-YE-YE.(Bis) 

\[Re]AH- U-U-U-\[La]U-U-U
\[Sol]U-U-U-\[La]U (Bis)
\endchorus

\beginverse
\[Re]Canta, \[La]canta, \[Sol]canta \[La]
para el \[Re]Señor \[La] \[Sol] \[La]
\[Re]canta, \[La]canta, \[Sol]canta \[La]
para el \[Re]Señor \[La]
1,2,3, y…..
\endverse 

\beginchorus
ESTRIBILLO 
\endchorus

\beginverse
Salta, salta, salta 
para el señor 
Salta, salta, salta 
para el señor 
1,2,3, y…
\endverse
\endsong
\beginsong{Alabo tu bondad}[]
    \beginverse
        \[Do]Todo mi ser canta \[Sol]hoy, 
        \[Fa]por las cosas que hay en \[Do]mí. 
        Gr\[Sol]acias te doy, mi Señor, 
        \[Fa]tú me haces tan fe\[Do]liz. 
        \[Fa]Tú me has rega\[Sol]lado tu amis\[Lam]tad, 
        \[Mi]confío en \[Fa]ti, me \[Do]llenas de tu \[Sol]paz. 
        \[Fa]Tú me haces sen\[Sol]tir tu gran bon\[Lam]dad, 
        \[Mi]yo \[Fa Do Rem]cantaré por \[Sol]siempre \[Do]tu fideli\[Sol7]dad. 
    \endverse
    \beginchorus
        \[Do]Gloria a \[Sol]Ti Se\[Re]ñor
        \[Lam]por tu bondad
        \[Fa Sol]Gloria, \[Lam Mi]Gloria
        \[Fa]simepre canta\[Do]ré
        \[Rem Sol...7]tu fidelidad
    \endchorus
    \beginverse
        Siempre a tu lado estaré
        alabando tu  bondad.
        A mis hermanos diré
        el gran gozo que  hallo en ti. 
        En ti podrán siempre encontrar fidelidad,
        confianza y amistad.
        Nunca fallará tu gran  amor,
        ni tu perdón.
        Me quieres
        tal como soy. 
    \endverse
\endsong
\beginsong{Dios está aquí}[]
\beginchorus
\[Do]DIOS ES\[Sol]TÁ A\[Lam]QUÍ,
\[Fa]TAN CIERTO COMO EL AIRE
\[Sol]QUE RESPI\[Do]RO \[Do7]
\[Fa]TAN CIERTO COMO LA MA\[Sol]ÑANA
SE \[Do]LEV\[Mim]ANTA\[Lam]
\[Fa]TAN CIERTO COMO QUE ESTE
\[Sol]CANTO LO PUEDES \[Do]OIR. \[(2º Sol)](BIS)
\endchorus
\beginverse
Lo puedes s\[Fa]entir movién\[Sol]dose
entre los \[Do]que a\[Do7]man.
Lo puedes \[Fa]oír cantan\[Sol]do con nosotros \[Do]aquí. \[Do7]
Lo puedes ll\[Fa]evar
cuan\[Sol]do por esa puerta \[Do]salgas. \[Do7]
Lo puedes gua\[Fa]rdar
muy \[Sol]dentro de tu co\[Do]razó\[Sol]n.
\endverse
\beginverse
Lo puedes notar junto a ti
en cualquier momento.
Le puedes hablar
de esa vida que le quieres dar
No temas ya más
Él es Dios y nos perdona a todos.
Jesús está aquí,
si tú quieres le puedes seguir.
\endverse

        \endsong

        \beginsong{Laudato Si}[]
            
            \beginchorus
            \[Do]LAUDATO SII, O MI SIGNORE
            \[Lam]LAUDATO SII, O MI SIGNORE
            \[Fa]LAUDATO SII, O MI SIGNORE
            \[Sol]LAUDATO SII, O MI SIGNORE
            \endchorus
            \beginverse
            \[Do]Y por todas las criaturas,
            \[Lam]por el sol y por la luna,
            \[Fa]por la noche y las estrellas,
            \[Sol]por el viento, el agua, el fuego.
            \endverse
            \beginverse
            Por la hermana madre tierra, nos da vida y
            nos sustenta; por las flores y la hierba, por
            los montes y los mares.
            \endverse
            \beginverse
            Porque todo el universo es un canto de
            alabanza hacia ti, Creador del mundo,
            siempre alegres cantaremos.
            \endverse
            \beginverse
            El sentido de la vida es cantarte y alabarte;
            y que toda nuestra vida sea siempre una
            canción.
            \endverse
        \endsong
\beginsong{Celebrarte}[]
    \beginverse
        \[La]Tiemp\[Sol]o atr\[Re]ás vi \[La]un cami\[Sol]no empez\[Re]ar
        \[La Sol Re]Inquietud, sueños y an\[Re]helos\[Sol Re]
        \[La]Pero un nuevo se\[Sol]ndero se abre ant\[Re]e mí
        \[La]Mirándome me pides \[Sol]un paso \[Re]más\[La Sol Re]
        \[Sim]Cuando Tú \[La]gritas mi \[Sol]nombre
        \[Sim]Dudas y m\[La]iedos me in\[Sol]vaden
        \[Sim]¿Crees que s\[Mim]eré yo \[La]capaz?
        \[Sim]¿Y si de\[Mim]cido arriesg\[La]arme?
        \[Sol]Hágas\[Sim]e Tu volunt\[La]ad
    \endverse
    \beginchorus
        \[Mim]Escucharte \[Sol]y sentirte de t\[Re]u mano
        \[La]sostenerme
        \[Mim]celebrarte, \[Sol]compartir,
        \[Re]com mi vida al\[La]abarte.
    \endchorus
    \beginverse
        Siento en mí: ¡quiero saber cómo soy!
        Máscaras, sombras que ocultan.
        Todos aquellos dones que recibí de Ti
        ¡Libera mi ser de verdad!
    \endverse
    \beginverse
        Como Tú quieres, sin prisa
        Moldeas mi corazón
        Siendo pequeño entre todos.
        Me eliges por el amor.
        Hágase Tu voluntad
    \endverse
    \beginchorus
        Escucharte y sentirte de Tu mano
        sostenerme
        celebrarte, compartir,
        com mi vida alabarte.
    \endchorus
\endsong
\beginsong{Jesús está entre nosotros}[]
\beginchorus
\[La]JESÚS ES\[Mi]TÁ \[Re]ENTRE \[La]NOSOTROS. \[Mi]
\[La]ÉL VIVE H\[Do#m]OY 
\[Re]Y SU ESPÍRITU A TODOS \[Mi]DA. 
\[La]JESÚS, RAZ\[Mi]ÓN \[Re]DE NUESTRA \[La Mi]VIDA, 
\[La]ES EL SE\[Do#m]ÑOR, 
\[Re]NOS REÚNE EN PUEBLO DE A\[Mi]MOR. 
\endchorus
\beginverse
\[Fa#m]Cambia nues\[Re]tras vidas con tu \[La Mi]fuerza, 
\[Fa#m]guárdanos por \[Re]siempre en tu \[La Mi]presencia. 
\[Fa#m]Tú eres \[Do#m]verdad, \[Re]tú e\[La]res la paz. \[Mi7] 
\endverse
 \beginverse
Rompe las cadenas que nos atan, 
llénanos de gracia en tu palabra; 
gracias, Señor, gracias, Salvador. 
\endverse
\beginverse
Nuestras existencias hoy te alaban, 
nuestros corazones te dan gracias. 
Tú eres amor, tú eres canción.
\endverse
   \endsong

 \beginsong{Te alabo en verdad}[]
\beginverse
\[Re]Aún en la tormen\[La]ta,
\[Re]aún cuando arre\[Re7]cie el \[Sol]mar
\[Re]te alabo, te a\[La]labo en ver\[Re]dad
\endverse
\beginverse
Aún lejos de los míos
aún en mi soledad
te alabo, te alabo en verdad
\endverse
\beginchorus
PUES \[Sol]SÓLO A TI \[La]TE TENGO
PUES \[Fa#m]TÚ ERES MI HE\[Sim]REDAD
TE A\[Sol]LABO, TE A\[La]LABO EN \[Re]VERDAD (BIS)
\endchorus
\beginverse
Aún sin muchas palabras,
aún cuando no sé alabar,
te alabo, te alabo en verdad
\endverse

        \endsong

%TODO
\beginsong{Al amor más sincero}[]
    \beginverse
        Al amor más sincero,
        a ese amor sin fronteras,
        a ese amor que dio su vida por amor

        Le entregué mi vida entera
    \endverse
\endsong
\beginsong{Sé que voy contigo}[]
        	\beginverse
	        		
\[Do]Apretemos \[Lam]nuestras manos  \[Fa]
en torno \[Sol]al altar, \[Do]
      Fa 	      Sol
\[Do]es Jesús cen\[Lam]tro de vida, \[Fa]fuente \[Sol]de
\[Do]verdad. 
\endverse
\beginchorus
			     
\[Fa]SOMOS COMO GRA\[Sol]NOS 
DE UNA ES\[Do]PIGA \[Mim]LLENA DE \[Lam]SOL,
Rem 		Sol7 		       Do…7 
\[Rem]SOMOS COMO \[Sol7]UN NUEVO AMA\[Do]NEC\[Do7]ER. 
Fa		    Sol
\[Fa]SEGUIMOS LAS \[Sol]HUELLAS 
           Do       Mim         Lam
DE JE\[Do]SÚS DE \[Mim]NAZAR\[Lam]ET, 
        Rem         Sol7 	     Do 
A TO\[Rem]DOS NOS \[Sol7]UNE LA \[Do]FE.
\endchorus
\beginverse
En la vida hay mil motivos para celebrar: un amigo, una mañana, el poder cantar.
La alegría nos desborda, queremos cantar; al sentirnos como hermanos debemos gritar.
        	\endverse
        \endsong

\beginsong{Extiende tu mano, Señor}[]
    \beginverse
        \[Sol]Extiende tu \[Re]mano, Se\[Mim]ñor,
        \[Do]pues todo l\[Re]o puede tu \[Sol7]amor.
        \[Do]Hoy te co\[Re]nfieso m\[Mim]i Dios\[Sim], mi Señor,
        sa\[Re]lvador.
        \[Sol]Ya no hay \[Re]tinieblas en \[Mim]mi
        \[Lam]solo hay \[Do]luz.
    \endverse
    \beginchorus
        Ahora e\[Lam]xtiende t\[Re]u mano, Je\[Sol]sús.
        \[Sol]Escucha Señor mi orac\[Re]ión,
        \[Mim]hoy te entrego mi co\[Do]razón.
        \[Sol]Borra mis faltas, Di\[Re]os salvador,
        \[Do]hoy te co\[Re]nfieso, Se\[Sol]ñor. (Bis)
    \endchorus
\endsong
 \beginsong{La razón escogida}[]
\beginverse
\[Re]Aprendí a seguirte, hace ya mucho ti\[Sim]empo
risas, juegos, can\[Sol]cion\[Mim]es, todo en su mo\[La]mento
\endverse
\beginverse
Han pasado los años, y ahora comprendo
que mi vida es más, mucho más que eso
\endverse
\beginverse
Levanto hoy la vista, no sé qué me pasa,
que es lo que me ocurre,
que es lo que me espanta
\endverse
\beginchorus
AUNQUE ME \[Re]CUESTE OÍR TU \[La]VOZ
POR CULPA DE ESTA LL\[Sol]UVIA
YO TE \[Re]QUIERO SEGUIR
A P\[La]ESAR DE MIS \[Sol]DUDAS
PORQUE TÚ ER\[Re]ES LA RAZ\[La]ÓN ESCOGID\[Sol]A
PORQUE \[Re]TÚ LE DAS
SEN\[La]TIDO A MI \[Sol]VIDA
PARA SEG\[La]UIR...
SÓLO PIENSO EN \[Re]TI, \[La]
SOLO PIENSO EN \[Sol]TI... \[La]
SOLO PIENSO EN \[Re]TI... \[La]
SOLO PIENSO EN \[Sol]TI... \[La]
\endchorus
\beginverse
\[Re]Surgen dudas, y mo\[La]mentos de confu\[Sol]sión
t\[Re]ú eres la luz, tú eres \[La]mi salida
sin \[Sol]ti estoy perdida
\[Re]Cuando todo es\[La]tá en silencio, \[Sol]puedo oír
\[Re]que algo en mi me \[La]dice que te \[Sol]he de
seguir
\endverse
\beginverse
\[Mim]No voy a de\[Re]jar que me in\[Sol]unde el te\[La]mor
\[Mim]he encont\[Sol]rado. Mi vo\[La]cación
mi vocación
\endverse
        \endsong

% TODO repetida con 73
\beginsong{Loco por ti}[]
    \beginverse
        \[Do]He visto guerras y \[Rem]engaños,
        \[Sol]gente luchando por \[Do]drogas, y a\[Mim Lam]ún peor,
        niños muriendo de \[Fa]hambre,
        miradas pro\[Rem]fundas, llenas de te\[Fa Sol]mor.
    \endverse

    \beginchorus
        QUIERO SE\[Fa]NTIRTE \[Sol]CERCA SEÑ\[Do]OR
        DAME LA \[Mi]VIDA Y DAME CA\[Lam]LOR
        SOLO D\[Fa]IME QUÉ ESPERAS,
        HA\[Sol]RÉ LO QUE QUIERAS MI D\[Do Do7]IOS
        SABES ME G\[Fa]USTA ESTAR L\[Sol]OCO POR \[Do]TI
        DIME SI \[Mi]NO LO QUÉ HARÍA YO A\[Lam]QUÍ
        AHORA T\[Fa]ENGO MOTIVOS
        Y LA V\[Sol]IDA UN SENTIDO,
        DANDO UN \[Do]“SÍ”
    \endchorus

    \beginverse
        Es como un pozo que se ha de llenar
        sendero que se ha de cruzar.
        Este es mi mundo sediento de amor
        y yo estoy dispuesto Señor.
    \endverse

    \beginverse
        Ponme en camino, no debo tardar
        mi barca ya se hace a la mar
        solo tus vientos me pueden llevar
        por las rutas del Reino, ABBA.
    \endverse    
\endsong
\beginsong{Tengo cara de niño}[]
            \beginverse
            \[Do]Tengo c\[Fa]ara de \[Do Fa]niño, 
            \[Do]tengo o\[Fa]jos de n\[Sol]iño, 
            teng\[Lam]o alma de n\[Mim]iño 
            y p\[Fa]uedo soñ\[Do Fa]ar como n\[Sol]iño que so\[Do Fa]y. 
            \endverse
            \beginverse
            Tengo manos pequeñas
            y las piernas de un niño,
            pero puedo saltar
            y hasta puedo correr
            y a\[Fa]brazar como un ni\[Sol]ño.
        	\endverse
                \beginchorus
            \[Fa]Y ES QUE PARA SE\[Do]R FELIZ
            \[Fa]Y ES QUE PARA AMA\[Do]R
            \[Lam]HAY QUE SER UN N\[Mim]IÑO
            \[Fa]HAY QUE VOLVER A NACE\[Sol (2ºDo)]R (2)
        	\endchorus
                \beginverse
                Disfrazado de hombre
                tú también puedes ser
                un niño que salta,
                un niño que ríe, un hombre feliz.
                \endverse
                \beginchorus
                \endchorus
        \endsong
\beginsong{Punto y Seguido}[]
    \beginverse
        \[Do]De nuevo hoy me siento a pensar
        (de nuevo hoy me \[Mim]siento a pen\[Lam]sar el mañana)
        el mañana,
        buscan\[Fa]do en la oscuridad, que \[Rem]sólo romperá
        esta ll\[Sol]ama
    \endverse
    \beginverse
        No sé si podré responder
        (no sé si podré responder, tu llamada)
        tu llamada
        no sé si podré acoger, el mensaje de tu fe
        tu Palabra
    \endverse
    \beginchorus
        \[Do]Y ahora al ver que esa l\[Fa]uz coje fuerza
        \[Lam]e ilumina tod\[Sol]a mi existencia
        \[Do]ya lo veo no \[Fa]me he equivocado
        \[Lam]mi futuro dib\[Sol]ujo a tu lado
        \[Do]me levanto, no \[Fa]soy el primero
        \[Lam]otros antes tus h\[Sol]uellas siguieron
        \[Do]porque siento su m\[Fa]isma esperanaza
        \[Lam]decidido ten\[Sol]go confianz\[Do Fa Lam]a.
        \[Sol]Confian\[Do Fa Lam Sol]za...
    \endchorus
    \beginverse
        Con firmeza ya me hizo andar
        (con firmeza ya me hizo andar tu llamada)
        tu llamada
        sin miedo a despertar, a esta nueva realidad
        del mañana
    \endverse
    \beginverse
        No me escondo quiero llevar
        (no me escondo quiero llevar tu Palabra)
        tu Palabra
        y es mi fuente la oración,
        la que un día me inspiró
        esa llama
    \endverse
    \beginchorus
        ESTRIBILLO
    \endchorus
    \beginverse
        ¡Y a\[SOl]hora te siento, Se\[Do]ñor!
        En la so\[Sol]nrisa de aq\[Fa]uél
        que me acom\[Sol]paña al an\[Do]dar.
        En el ca\[Mim]mino que ha\[Fa]cemos
        juntos \[Sol]con humil\[Do]dad.
        En esa \[Mim]mano que \[Fa]llega
        donde \[Sol]no sé lle\[Fa]gar.
    \endverse
    \beginverse
        Te veo, te sie\[Sol]nto...
        (te veo, te \[Fa]siento)
        te veo, te \[Sol]siento...
        (te veo, te \[Fa]siento)
        te veo, te \[Sol]siento...
        (te siento)
        ¡Y ahora te sigo, Se\[Do Fa Lam Sol]ñor!
    \endverse
\endsong
\beginsong{Sabe de amor}[]
\beginverse
\[Sol]Quien tiende \[Mi]las manos, \[Lam]quien da el \[Si]valor
\[Do]y quien no pi\[Lam]de nada, \[Sol]sabe \[Do]de am\[Sol]or.
\endverse
\beginverse
Quien cada mañana saluda alegre al sol,
quien es fuerte, quien vive, sabe de amor.
\endverse
\beginverse
Quien siempre espera del otro lo mejor,
y quien nunca se cansa, sabe de amor.
\endverse
\beginverse
Quien tiene escondida su fuerza en una
cruz es que ha recibido de Dios esa luz.
\endverse
\beginverse
Quien fuerte proclama que Dios es
Salvador es que ha conocido que Dios es
amor, que Dios es amor...
\endverse
        \endsong

  \beginsong{Volver a ti}[]
            \beginverse
            \[Re]Si cuando \[Do]miro al cielo só\[Mim]lo lluvia v\[Sol La]eo aquí
            \[Re]y bajo \[Do]techo me r\[Mim]efugio n\[Sol]o sé a d\[La]onde i\[Re]r
            \endverse
            \beginverse
            Y me marché de casa la herencia malgasté
            pensé que otra vida sin mi Padre
            solo yo podría hacer.
            \endverse
            \beginverse
            \[Sol]Me di \[Re]cuenta que a\[Do]sí ya no \[Mim]podía conti\[Sol]nuar
            pienso que a\[Re]hora de re\[Do]greso p\[Mim]erdón le pe\[Sol]diré
            Y le d\[Re]iré, lo \[Do]mucho que le \[Mim]quiero
            a mi Padre del \[Mim]cielo, con s\[Do]u ayuda av\[La La7]anzaré.
            \endverse
            \beginchorus
            \[Re]TAL VEZ, S\[Sim]ERÁ,
            \[Sol]QUE EL MUNDO \[La]OLVIDE QUE ERES \[Re]TÚ
            LA LU\[Sim]Z QUE ME HACE BRI\[Sol]LLAR
            EN LOS MOM\[La]ENTOS DE OSCURI\[Re]DAD
            \[Re]SÉ QUE P\[Sim]OR TI
            \[Sol]MI CORAZÓ\[La]N AHORA \[Re]LATE
            BUSC\[Sim]ARÉ TU VOLUN\[Sol La]TAD
            IRÉ HACIA DELA\[Re]NTE
            \endchorus
            
        \endsong
\beginsong{Agua lávame}[]
\beginverse
\[Sol]Agua lávame, puri\[Do]fícam\[Sol]e.
Dame agua, tu es\[Mim]píritu.
\[Do]Agua \[Re]lávam\[Sol]e.
\endverse
   \endsong

\beginsong{Caminad}[]
\beginverse
\[Mim]Caminad \[Si7]mientras ten\[Mim]éis luz,
\[Sol]antes que os \[Do]envuelva la t\[Sol Re]iniebla.
\[Mim Si7 Mim]Caminad.
\endverse
\endsong

\beginsong{Compartir}[]
\beginverse
\[Do]Com\[Sol]parti\[Lam]r, hay que \[Fa]com\[Do]parti\[Sol]r
para \[Do]ser fe\[Lam]liz
hay que \[Rem Sol Do]compartir
\[Do Sol Lam]Compartir, hay que \[Fa Do Sol]compartir
que no te \[Do]falte \[Lam]a ti
lo que me \[Rem Sol]sobra a \[Do]mi.
\endverse
   \endsong

\beginsong{Dame vida}[]
\beginverse
\[Re]Dame vida, \[Ka]dame la vida,
\[Sol]dame tu vida, mi \[La]Dios
\endverse
\beginverse
(Conviérteme, Señor,
conviértenos a tu Reino)
\endverse
\beginverse
Enséñame tu tierra,
guíame a tu tierra,
llévame a tu tierra, Señor.
\endverse
   \endsong

\beginsong{De noche iremos}[]
\beginverse
\[Lam]De noche iremos, de \[Rem]noche.
Que \[Do]para encontrar la \[Mi]fuente
\[Lam]solo la \[Sol]sed nos \[Do]alumb\[Mi]ra,
\[Fa]solo la sed nos \[Mi]alumbra.
\endverse
   \endsong

\beginsong{El Señor es mi pastor}[]
\beginverse
\[La]El Señor es mi pas\[Mi]tor, nada me fa\[Fa#m]lta
El Se\[Re]ñor es mi pas\[Mi]tor.
\[La]El Señor es mi pas\[Mi]tor, nada me fa\[Fa#m]lta
El Se\[Re]ñor es \[Mi]mi pas\[La]tor.
\endverse
   \endsong

\beginsong{A nadie diré}[]
\beginverse
\[Do]A \[Sol Lam]nadie diré: ¡\[Fa]te a\[Do Sol]doro!
\[Do]A \[Sol Lam]nadie diré, \[Fa]sino \[Sol]solo a \[Do]ti.
\endverse
\beginverse
A nadie diré más te adoro sino sólo a ti, Señor.
A nadie diré más te adoro sino sólo a ti.
A nadie diré más te adoro sino sólo a ti, Señor,
sino sólo a ti.
\endverse
\endsong

\beginsong{El Dios de la vida}[]
\beginverse
\[La]Yo \[Re]amo al \[Mi]Dios de la vi\[La]da
Y en su \[Fa#m]nombre su\[Rem]plico \[Mi]miserico\[La]rdia.
\endverse
   \endsong

\beginsong{En mi debilidad}[]
\beginverse
\[Re]En mi debilidad, Tú me haces \[Sim]fuerte
En mi debili\[Sol]dad, \[Sol]me haces f\[Re]uerte.
Sólo en tu amor me haces \[Re7]fuerte,
sólo en tu \[Sol]vida me haces \[Solm]fuerte.
En mi debili\[Re]dad, te \[Sol]haces \[La]fuerte en \[Re]mí
\endverse
\endsong

\beginsong{En todo amar y servir}[]
\beginverse
\[La]En todo \[Re]amar y ser\[La]vir
\[La]En todo \[Re]amar y ser\[La]vir
\[La]En todo \[Re]amar y ser\[Mi]vir
\endverse
   \endsong

\beginsong{La misericordia}[]
\beginverse
\[Rem]La miseri\[La]cordia \[Rem]del Señ\[Do]or
\[Fa]cada \[Do]día \[Rem La Rem]cantaré
\endverse
   \endsong

    \beginsong{En nuestra oscuridad}[]
\beginverse
\[Mi]En nuestra oscuri\[Lam]dad
enciende la llama de tu amor \[Sol Do]Señor,
de tu a\[Fa]mor \[Do Sol]Señor.
En nu\[Do]estra os\[Lam]curida\[Fa]d
enciende la \[Mi7]llama de tu a\[Lam]mor \[Rem Si7]Señor,
a\[Lam]mor \[Rem Mi7]Señor
\endverse
   \endsong

\beginsong{Escucha tú}[]
\beginverse
\[Do]Escucha \[Sol]tú la pa\[Lam]labra de \[Fa]Dios
No \[Do]solo con tus o\[Lam]ídos ta\[Fa]mbién
Con tu cora\[Sol]zón
\[Do]Escucha \[Sol]tú, la pa\[Lam]labra de \[Fa]Dios
\[Do]Y estate siempre \[Fa]atento
\[Sol]A su \[Do]voz.
\endverse
   \endsong

\beginsong{Hola Dios}[]
\beginverse
\[Sol]Hola, D\[Re]ios, es\[Do]toy a\[Sol]quí.
\[Do]Gracias te \[Sol]doy por da\[Do]rme la vi\[Re]da.
\[Sol]Hazla n\[Re]ueva to\[Do]dos l\[Sol]os días.
\[Do]Buenos dí\[Sol]as, \[Re]mi Seño\[Sol]r.
\endverse
   \endsong

\beginsong{A pesar de mis pobres talentos}[]
\beginverse
\[Re]A pesar de mis \[Sim]pobres talentos
\[Sol]Dios ha querido ser\[La]virse de \[Re]mí.
A pesar de mis \[Sim]pobres talentos
\[Sol]Dios ha querido ser\[La]virse de \[Sim]mí.
\endverse
\beginverse
En a\[Fa#m]mor, o\[Sol]rando con la \[La]vida,
\[Sim]en la la\[Fa#m]bor, ba\[Sol]jo el humilde \[Mim]techo
de la \[Sol]Santa Fa\[La]milia,
la \[Re]paz.
\endverse
   \endsong

  \beginsong{Iubilate Deo}[]
\beginverse
\[La]Iu\[Mi]bi\[Fa#m]late \[Re]Deo, \[La]iubi\[Mi Fa#m]late \[Mi]Deo,
¡\[La]A\[Mi7]lelui\[La]a!
\endverse
\endsong

\beginsong{Jesús es Señor}[]
\beginverse
\[Do]Jesús es, \[Fa]Jesús es Se\[Do]ñor \[Sol]
\[Fa]Jesús es, Je\[Sol]sús es Se\[Do]ñor \[Mim Lam]
\[Fa]Jesús es, \[Sol]Jesús es Se\[Do]ñor
\endverse
   \endsong

\beginsong{Laudate omnes gentes}[]
\beginverse
\[Do]Laudate \[Fa]omnes \[Sol]gent\[Lam]es,
\[Do]laudate \[Fa]Domi\[Sol]num.
\[Do]Laudate \[Fa]omnes \[Sol]gent\[Lam]es,
lau\[Fa]date \[Sol]Domi\[Do]num.
\endverse
   \endsong

\beginsong{Lo que agrada a Dios}[]
\beginverse
\[Re]Lo que agrada a Dios,
\[Sol]de mi pequeña \[Re]alma,
es que ame mi \[Sol]pequeñe\[La]z y mi po\[Re Re7]breza (Bis)
Y es la espe\[Sol]ranz\[La]a, \[Re Fa#m Sim]plena,  
q\[Sol]ue teng\[La]o en su miseri\[Re]cordia (Bis)
\endverse
   \endsong

\beginsong{Magnificat}[]
\beginverse
\[Re]Magnifi\[Sol]cat, \[La]magnifi\[Re]cat,
magnifi\[Sol]cat, anima \[La]mea Domi\[Re]num.
Magnifi\[Sol]cat, \[La]magnifi\[Re]cat.
Magnifi\[Sol]cat anima \[La Re]mea.
\endverse
   \endsong

\beginsong{Más allá}[]
\beginverse
\[Re]Más allá de mis \[La]miedos,
más \[Sim]allá de \[Fa#m]mi inseguri\[Sol]dad,
quiero \[Mim]darte mi res\[La]puesta.
Aquí \[Re]estoy para ha\[La]cer tu volun\[Sim]tad
para que mi a\[Fa#m]mor sea de\[Sol]cirte Sí\[La]
hasta el f\[Re]inal.
\endverse
   \endsong

\beginsong{Mi alma se reposa}[]
\beginverse
\[Mim]Mi a\[Si7]lma \[Mim]se r\[Do]eposa \[Sol]en \[Re]paz \[Mim]
sólo en \[Si7]Dios. Él \[Mim]es \[Do]mi \[Re]salvac\[Sol]ión.
\[Mim]Si, \[Do]sólo en \[Sol]Dios \[Re]mi \[Mim]alma \[Lam]se re\[Si7]posa, \[Mim]
se re\[Si7]posa en \[Mim]paz.
\endverse
   \endsong

\beginsong{Mi fortaleza}[]
\beginverse
\[Re]Te amo, Se\[La]ñor, fortaleza \[Sim]mía,
ro\[Sol]ca mía, casti\[La]llo mío, mi libera\[Re]dor.
\[Sol]Dios mío en \[La]ti confia\[Sim]ré.
Mi \[Sol]escudo eres Tú
y la \[La]fuerza de mi salva\[Re]ción
\endverse
   \endsong

\beginsong{Muéstranos, Señor}[]
    \beginverse
        \[Do]Muéstran\[Sol]os Se\[Lam]nor\[Mim],
        \[Fa]Tus \[Do]caminos, \[Fa]tus ca\[Sol]minos,
        \[Do]para no hace\[Do]r nada m\[Lam]ás\[Mim]
        \[Fa]que lo que s\[Do]es t\[Fa]u volunt\[Sol]as.
    \endverse
\endsong
\beginsong{Muéveme}[]
\beginverse
\[Sol]Muéveme mi \[Do]Dios hacia \[Sol]ti,
que no me mu\[Do]evan los hi\[Re]los de este mun\[Mim]do
Oh,oh
\[Do]Muéveme, at\[Lam]ráeme hacia ti\[Do]
desde lo pro\[Re]fundo
\endverse
   \endsong

\beginsong{Abre mis labios}[]
\beginverse
{\nolyrics Intro: \[Re] \[Re7]}
Abre mis \[Sol]labios, S\[La]eñor,
para po\[Re]dert\[Fa#m]e alaba\[Sim]r.
Abre mi \[Sol]coraz\[La]ón,
para po\[Sol]derte \[La]ador\[Re]ar.
\endverse
   \endsong

\beginsong{Nada has pensado}[]
    \beginverse
        \[Do]Nada \[Fa]has pens\[Fa]ado
        que no puedas hac\[Sol]er.
        \[Do]Nada \[Fa]has so\[Do]ñado imposible de \[Sol]ver.
        Ni el m\[Fa]ar, ni el ciel\[Rem]o,
        \[Fa]nada está \[Sol]lejos de \[Do]ti.
        Todo es\[Sol]tá en tu adent\[Lam]ro,
        \[Rem]todo pue\[Sol]des hacerlo
        \[Do]si te \[(Fa)]fías de mi.
    \endverse
\endsong
\beginsong{Alabaré}[]
\beginchorus
Alaba\[Re]ré, alabaré, alabaré, alabaré,
A\[La]labaré a mi Se\[Re]ñor
Alaba\[Re]ré, alabaré, alabaré, alabaré,
A\[La]labaré a mi Se\[Re]ñor.
\endchorus
\endsong
\beginsong{Nada te turbe}[]
\beginverse
\[Lam]Nada te \[Rem]turbe, \[Sol]nada te es\[Do]pante
\[Fa]quien a Dios \[Sol]tiene, \[Mim]nada le \[Lam]falta.
\[Lam]Nada te \[Rem]turbe, \[Sol]nada te es\[Do]pante.
\[Fa]Sólo \[Sol]Dios \[Mim Lam]basta.
\endverse
   \endsong

\beginsong{Nada nos separará}[]
\beginverse
\[Do]Nada nos se\[Fa]para\[Do]rá, (nada nos se\[Fa]para\[Do]rá)
\[Do]nada nos se\[Fa]para\[Do]rá, (nada nos se\[Fa]para\[Do]rá)
\[Do]nada nos \[Mim]separa\[Lam]rá \[Fa]del am\[Sol]or de Di\[Do]os.
\endverse
   \endsong

\beginsong{Nadie como Tú}[]
    \beginverse
        \[Do]Nadie \[Sol]como \[Lam]Tú,
        \[Fa]nadie \[Rem]más llena mi \[Sol]ser como Tú
        \[Do]por la \[Sol]eternidad po\[Lam]dría buscar...
        más n\[Fa]o hay nadie \[Sol]como \[Do]Tú.
    \endverse
\endsong
\beginsong{Nazaret}[]
\beginverse
\[Lam]En amor, si \[Mi]oras y tra\[Lam]bajas, tendrás \[Do]paz.
Naza\[Fa]ret mo\[Sol]delo para \[Do Mi]ti.
En a\[Lam]mor, si \[Mi]oras y tra\[Lam]bajas, tendrás \[Do]paz.
Na\[Fa]zaret, mo\[Sol]delo de tu \[Lam]sí.
\endverse
   \endsong

\beginsong{Oh Señor, ten piedad}[]
\beginverse
\[Do]Oh, Señor ten piedad, de los hom\[La]bres
ten piedad
\[Fa]mirada suplicante levan\[Sol]tamos hacia ti.
\[Do]Oh, Señor, ten piedad, de los \[La]hombres
ten piedad
\[Fa]escucha nuestros gritos en auxilio \[Sol]ten
pie\[Do]dad
\endverse
   \endsong

\beginsong{Óyeme, te basta mi gracia}[]
    \beginverse
        \[Sol]Óyeme, te b\[Do]asta mi graci\[Sol]a
        déjate ha\[Do]cer por mi gracia\[Re],
        \[Do]que en tu debilida\[Re]d,
        se mue\[Do]stra \[Re]mi fuerza.
    \endverse
\endsong
\beginsong{Paz}[]
\beginverse
\[Mi]Paz, os doy la \[Do#m]paz,
os doy la \[La]luz, os doy \[Si]la libertad
\endverse
   \endsong

\beginsong{Que sea tu vida}[]
\beginverse
\[Do]Que sea tu vida la sal,
que sea tu vida la \[Sol]luz,
\[Fa]sal que \[Do]sala, \[Mim]luz que \[Lam]brilla
\[Fa]sal y \[Sol]fuego es Je\[Do]sús
\endverse
   \endsong

\beginsong{Quedaos aquí}[]
    \beginverse
        Que\[Rem]daos aq\[Solm]uí y ve\[Rem]lad co\[Solm]nmigo.
        \[Rem]Vela\[Do]d y \[Fa]ora\[Do]d, \[Rem]velad y \[Lam]orad.
    \endverse
\endsong
 \beginsong{Quiero alabarte}[]
\beginverse
Quiero a\[Re]labarte más y más aún
quiero ala\[Sim]barte más y más aún
bus\[Sol]car tu voluntad, tu \[Mim]gracia conocer
quiero ala\[Re]bart\[La]e
\endverse
\beginverse
Las aves del cielo cantan para ti las
bestias del campo, reflejan tu poder,
quiero cantar, quiero levantar
mis manos a ti
\endverse
   \endsong

\beginsong{Recíbeme}[]
\beginverse
\[Re]Recíbeme,
con todo lo que \[La]tú pusiste en \[Sim]mi
con todas esas \[Fa#m]ganas de vi\[Sol]vir
con \[Mim]toda mi mi\[La]seria
\endverse
   \endsong

\beginsong{Sed sencillos}[]
    \beginverse
        \[Re]Sed sencillos, \[Sol]sed humildes.
        \[Re]Cuanto más sen\[Sol]cillos s\[La]ois
        \[Re]mayor bi\[Re7]en har\[Sol]éis al \[Solm]mundo,
        \[Re]mayor \[La]bien ha\[Re]réis.
    \endverse
\endsong
\beginsong{Sé mi luz}[]
\beginverse
\[La]Sé mi lu\[Re]z, enciende mi \[La]noche
sé mi \[Fa#m]luz, enc\[Re]iende mi \[Mi]noche.
\[La]Sé mi lu\[Re]z, enciende mi \[La]noche
mi \[Re]noch\[Mi]e, sé mi \[La]luz.
\endverse
   \endsong

\beginsong{Si no tengo amor}[]
\beginverse
\[Do]Si no tengo amor, \[Fa]no soy na\[Do]da
nada \[Lam]soy, \[Fa]si no tengo am\[Sol]or
nada \[Fa]soy, \[Sol]nada \[Mim]soy, si \[Fa]no teng\[Sol]o am\[Do]or
\endverse
   \endsong

\beginsong{Siento tu llamada}[]
    \beginverse
        \[Re]Sient\[La]o tu ll\[Sim]ama\[Fam]da, y co\[Sol]nfío \[Re]en \[La7]ti.
    \endverse
\endsong
\beginsong{Shema Israel}[]
\beginverse

\[Sim]Shema, Israel:
Ado\[Re]nai elohenu, \[La]Adon\[Fa#m]ai ehad\[Sim]
Sh\[Sim]ema, Israel:
Ado\[Re]nai elohenu, \[La]Adon\[Fa#m]ai ehad\[Sim]
She\[Re]ma, Israel:
Ado\[La]nai elohenu, \[Sol]Adonai eh\[Re]ad
Shema, Israel:
Ad\[La]onai elohenu, \[Sol Fa#m]Adonai eha\[Sim]d
\endverse
\beginverse
Escucha Israel:El Señor es nuestro Dios,
uno es el Señor (4)
\endverse
   \endsong

 \beginsong{Sois la sal}[]
\beginverse
\[Do]Sois \[Sim]la \[Lam]sal
que puede \[Solm]dar sab\[Fa]or
a la vida\[Sol7]
\endverse
\beginverse
Sois la luz que tiene que alumbrar
llevar a Dios.
\endverse
   \endsong

\beginsong{Adoramus te domine}[]
\beginverse
\[Do]Uh, \[Sol]uh, \[Lam]uh, \[Fa]adoramus \[Sol]te, Domi\[Re]ne
\endverse
   \endsong

\beginsong{Somos una familia}[]
    \beginverse
        \[Mi]Somos una familia
        un auténtico mogollón
        \[La]una igle\[Si]sia diverti\[Mi]da
        que donde \[La]vamos armam\[Si]os la de Di\[Mi]os.
    \endverse
\endsong
\beginsong{Te quiero}[]
\beginverse
\[Do]Enséñame a decir “te quiero”,
en\[Lam]séñame a decírtelo.
En\[Fa]séñame a decir “te quiero”, Padre,
en\[Sol]séñame a decírtelo.
\endverse
   \endsong

\beginsong{Ubi caritas}[]
\beginverse
\[Do]Ubi \[Sol]carit\[Lam]as, et \[Fa]a \[Do]- \[Sol]mor.
\[Do]Ubi \[Sol]carit\[Lam]as, De\[Fa]us \[Sol]ibi \[Do]est.
\endverse
   \endsong

\beginsong{Tu fidelidad}[]
\beginverse
\[Do]Tu fidelidad es \[Rem]grande
\[Sol]Tu fidelidad in\[Do]compa\[Mim]rable \[Lam]es
\[Do]Nadie como tú, \[Rem]bendito Di\[Fa]os
Grande \[Sol]es tu fideli\[Do]dad \[(Fa Sol)]
\endverse
   \endsong

\beginsong{Tú, Señor, quieres ser}[]
\beginverse
\[Do]Tú, Se\[Sol]ñor, quieres \[Lam]ser
\[Fa]fuego en \[Rem]nuestros cora\[Fa]zon\[Sol]es.
\[Do]Tú, Se\[Sol]ñor, quieres \[Lam]ser
\[Fa]alas \[Sol]en nuestros \[Do]pies.
\endverse
   \endsong

\beginsong{Ven y descánsate}[]
    \beginverse
        \[Lam]Ven y descánsate
        \[Do]ven y descánsate e\[Rem]n Dios, en \[Mi]Dios (bis)
    \endverse
    \beginverse
        Y \[Fa]deja que \[Do]Dios sea Dios,
        deja que Dios sea Dios.
        \[Fa]Tú s\[Mi]ólo ad\[Lam]órale.
        \[FA]Tú s\[Mi]ólo ad\[Lam]órale
    \endverse
\endsong
\beginsong{Yendo contigo}[]
\beginverse
\[Re]Yendo con\[La]tigo\[Sim], nada me inq\[Fa#m]uieta,
\[Sol]marcho con \[Mim]paz y fuer\[La]za.
\[Re]Yendo con\[La]tigo\[Sim], todo se esp\[Fa#m]era,
\[Sol]cada ma\[La]ñana es \[Re]nueva.
\endverse
   \endsong

\beginsong{Ayúdame a caminar}[]
\beginverse
\[Do]Ayúdame a caminar
contigo iré sobre las olas de la \[Sol]mar.
Y canta\[Mim]ré quien eres \[Lam]Tú
tú eres Jesús, \[Do]mi Dios, mi \[Fa]bien mi \[Sol]liberta\[Do]d.
\endverse
   \endsong

\beginsong{Yo te doy lo que soy}[]
\beginverse
\[Do]Yo te doy lo \[Rem]que soy,
lo que \[Lam]tengo y qu\[Fa]iero
y cuan\[Rem]to pue\[Sol]do ten\[Do]er.
\endverse
   \endsong

\beginsong{Amén}[]
    \beginverse
    \[Mi Si7 Mi Si7]A – mén.
    \[Mi Si7 Mi Si7]A- mén.
    \[Mi Sol#m Do#m Sol#m La Si7 Mi]A - mén
    \endverse
\endsong
\beginsong{Bonum est confidere}[]
\beginverse
\[Rem]Bonum est con\[La7]fidere in \[Rem]Do\[Do]mi\[Fa]no
\[Solm]bonum\[Do] spe\[Fa]rare \[Rem]in \[Sol7]Do\[Lam]mi\[Rem]no
\endverse
   \endsong

\beginsong{Bonum est confidere}[]
\beginverse

{\nolyrics Intro: \[Do]\[Sol]\[Lam]\[Sol]}
Surrexit Dominus vere, alleluia alleluia
Surrexit Christus hodie, alleluia alleluia
\endverse
   \endsong

\beginsong{Hossana Hey}[]
    \beginverse
        \[Re]Hosanna Hey, hosanna ha
        \[Re7]Hosanna hey, \[Sol]hosanna hey, hosanna \[La Re]ha
        El es el \[Re]Santo es el Hijo de Marí\[Sol]a
        es el Dios de Isra\[La]el, es el Hijo de Dav\[Re]id.
        Vamos a él con espigas de mil trigos y
        con mil ramos de olivos,
        siempre alegres, siempre en paz.
        El es el Cristo, es el unificador,
        es hosanna en las alturas es hosanna en el amor.
        Es la alegría la razón de mi existir
        es la vida de mis días
        es consuelo en mi sufrir.
    \endverse
\endsong
\beginsong{Bendice mi alma al Señor}[]
\beginverse
\[Sim]Bendice mi \[Mim]alma al Se\[Sim]ñor,
y tam\[Sol]bién a su \[La]santo \[Re Fa#]nombre.
Ben\[Sim]dice mi \[Mim]alma al \[Sim]Señor.
Él \[Sol]me li\[La]bró de la \[Sim]muerte.
\endverse
   \endsong

\beginsong{Yo quiero ser un vaso nuevo}[]
\beginverse
\[Do]Gracias quiero \[Lam]darte por \[Fa]amar\[Sol]me,
\[Fa]gracias qu\[Sol]iero darte yo a ti \[Do]Señor. Hoy
soy \[Do7]feliz \[Fa]porque \[Fam]te conocí, \[Do]gracias por
\[Sol7]amarme a mi \[Do]también.
\endverse
\beginchorus
\[Do]YO QUIERO \[Sol7]SER, SEÑOR AMADO COMO EL
\[Fa]BARRO \[Sol7]EN MANOS DEL \[Sol]ALFARERO; TOMA MI
\[Do7]VIDA, H\[Fa]AZLA D\[Fam]E NUEVO. \[Do]YO QUIERO SER
\[Sol7]UN VASO \[Do]NUEVO.
\endchorus
   \endsong

 \beginsong{Te presentamos el vino y el pan}[]
\beginchorus
\[Sol]TE PRESENTAMOS EL \[Do]VINO Y \[Sol]EL PAN 
BE\[Do]NDITO \[Sol]SEAS POR \[Re7]SIEMPRE SE\[Sol]ÑOR.
\endchorus
\beginverse
\[Sol]Bendito \[Do]seas Se\[Sol]ñor por \[Do]este Pan que nos \[Sol]diste
\[Do]Fruto de la tierra y del tra\[Sol]bajo de los ho\[Re7]mbres.
\endverse
\beginchorus
\[Sol]TE PRESENTAMOS EL V\[Do]INO Y EL P\[Sol]AN 
BE\[Do]NDITO \[Sol]SEAS POR S\[Re7]IEMPRE SEÑOR. \[Sol]
\endchorus
\beginverse
\[Sol]Bendito \[Do]seas Se\[Sol]ñor el \[Do]vino Tú nos lo \[Sol]diste 
fr\[Do]uto de la tierra y del tra\[Sol]bajo de los hom\[Re7]bres.
\endverse
   \endsong

\beginsong{Tal como soy}[]
\beginverse
\[Sol]Tal como soy Señor, sin nada que ofr\[Mim]ecer más que mi canción. No tengo más que \[Do]darte pues todo es tuyo Se\[Re7]ñor.
\endverse
\beginverse
\[Sol]Tal como soy Señor, sin nada que entre\[Mim]gar más que el corazón. Me rindo todo a \[Do]Ti, tómame Señor, tal como \[Re7]soy.
\endverse
\beginchorus
\[Sol]ACÉPTAME COMO OFRENDA DE AMOR, COMO \[Mim]UN SACRIFICIO A\[Do]GRADABLE EN TU \[Sol]HONOR. GRATO PERFUM\[Re]E, YO QUIERO SER \[Do Re]SEÑOR.
\endchorus
        \endsong

\beginsong{Santa María del Camino}[]
    
    \beginverse
        \[Do]Mientras recorres la \[Fa]vida,\[Do] 
        \[Sol]Tu nunca solo estás\[Do Do7], 
        \[Fa]Contigo por el ca\[Do Lam]mino, 
        \[Rem]Santa María va\[Sol7 Do Do7].
    \endverse

    \beginchorus
        \[Fa]Ven con nosotros a ca\[Do]minar, 
        \[Sol]Santa maría \[Do Do7]ven. (2)
    \endchorus

\endsong
\beginsong{Resucitados - Kénosis by RC}[]
    \beginverse
        \nolyrics Cej 7
    \endverse
    \beginverse
        Enséña\[Mim7]me a abraz\[Do9]ar
        Mi c\[Sol]ruz, Se\[Re]ñor
        A no escandali\[Mim7]zarme
        De \[Do9]todas mis heridas
        \[Sol]Así me has hecho tú
        \[Re]Así me amas, mi Dios
    \endverse
    \beginverse
        Enséñame a abrazar
        Tu voluntad, Señor
        Si no puedo apartar
        El cáliz de mi realidad
        Renuncio a mi plan
        Acojo tu verdad
    \endverse
    \beginchorus
        Quiero \[Mim7]ser, ser la\[Do9]vado
        \[Sol]Ser curado, y quiero
        \[Re]Morir a mí mismo
    \endchorus
    \beginverse
        Enséñame a abrazar
        Mi condición, Señor
        No puedo dar la talla
        No soy suficiente
        Quiero dejar de hacer
        Para dejarme amar (por ti)
    \endverse
    \beginchorus
        Quiero ser, ser lavado...
    \endchorus
    \beginverse
        Quiero ser, ser lavado
        Y quiero estar arrodillado
        Quiero vaciarme
        Y, en el fondo, siento
        Que no soy valorado
        Estoy herido
        Soy rechazado
        Soy pecado
    \endverse
    \beginverse
        Pero, soy amado
        Soy aceptado
        Soy rescatado
        No soy esclavo
    \endverse
    \beginchorus
        Quiero ser, ser lavado...
    \endchorus
    \beginverse
        Y quiero estar arrodillado
        Quiero vaciarme
        Quiero vivir como
        Resucitado, resucitado
        Resucitado, resucitado
    \endverse
    \beginchorus
        Quiero ser, ser lavado...
    \endchorus
    \beginverse
        Y quiero estar arrodillado
        Quiero vaciarme
        Quiero vivir como
        Resucitado
    \endverse
\endsong
\beginsong{Quien amó primero}[]
\beginverse
\[Do]Quién \[Sol]amó pr\[Lam]imero,
\[Do]quién nos \[Sol]dijo “A\[Lam]mad como amo yo”. 
\[Fa]Quién am\[Sol]ó a \[Do]todas, a \[Lam]todos, 
\[Fa]quién \[Sol]en todo \[Lam]amó.
\endverse
   \endsong

 \beginsong{¿Qué más quieres?}[]
\beginverse
\[Sol]¿Y qué más quieres que alumbrar?      
¿Y qué más quieres que \[Do]dar vida?
¿Qué más quieres que crear alg\[Re Do]o 
que en ti está como se\[Sol]milla?

\endverse
        \endsong

\beginsong{Pongo mi vida en tus manos}[]
\beginchorus
\[La]Pongo mi \[Re]vida en \[La]tus manos,
Padre mío me aban\[Re]dono a \[Mi]ti.
\[Re]Haz de \[Do#m]mí lo que quie\[Fa#m]ras
estoy dis\[La]puesta a acep\[Re]tarlo to\[Mi]do...
\endchorus
\endsong
\beginsong{Pescador de hombres}[]
\beginverse
\[Do]Tú \[Sol]has venido a \[Lam]la o\[Do7]rilla,
no has\[Fa] bus\[Rem]cado ni a sabios ni a ri\[Sol]cos.
\[Sol]Tan só\[Do]lo \[Sol]quieres que \[Lam]yo te \[Do7]siga.
\endverse

\beginchorus
Se\[Fa]ñor, \[Sol]me has mi\[Do]rado \[Lam]a los ojos,
son\[Rem]rien\[Do]do has di\[Do]cho mi \[Do7]nombre.
En \[Fa]la a\[Sol]rena he de\[Do]jado \[Lam]mi barca
junto\[Rem] a \[Sol]Ti buscaré \[Do]otro mar.
\endchorus

\beginverse
Tú sabes bien lo que tengo,
en mi barca no hay oro ni espada,
tan sólo redes y mi trabajo.
\endverse

\beginchorus
Se\[Fa]ñor, \[Sol]me has mi\[Do]rado \[Lam]a los ojos,
son\[Rem]rien\[Do]do has di\[Do]cho mi \[Do7]nombre.
En \[Fa]la a\[Sol]rena he de\[Do]jado \[Lam]mi barca
junto\[Rem] a \[Sol]Ti buscaré \[Do]otro mar.
\endchorus

\beginverse
Tú necesitas mis manos,
mi cansancio que a otros descanse,
amor que quiera seguir amando.
\endverse

\beginchorus
Se\[Fa]ñor, \[Sol]me has mi\[Do]rado \[Lam]a los ojos,
son\[Rem]rien\[Do]do has di\[Do]cho mi \[Do7]nombre.
En \[Fa]la a\[Sol]rena he de\[Do]jado \[Lam]mi barca
junto\[Rem] a \[Sol]Ti buscaré \[Do]otro mar.
\endchorus

\beginverse
Tú, pescador de otros lagos,
ansia eterna de hombres que esperan.
Amigo bueno que así me llamas.
\endverse

\beginchorus
Se\[Fa]ñor, \[Sol]me has mi\[Do]rado \[Lam]a los ojos,
son\[Rem]rien\[Do]do has di\[Do]cho mi \[Do7]nombre.
En \[Fa]la a\[Sol]rena he de\[Do]jado \[Lam]mi barca
junto\[Rem] a \[Sol]Ti buscaré \[Do]otro mar.
\endchorus
\endsong
\beginsong{Océanos - Hillsong}[]
\beginverse
{\nolyrics Intro: \[Sim] \[La/Do#] \[Re] \[La] \[Sol]}
\[Sim]Tu voz me llama\[La/Do#] a las \[Re]aguas
Donde \[La]mis pies pued\[Sol]en fallar
\[Sim]Y ahí te encuentro\[La/Do#] en lo in\[Re]cierto
Camin\[La]aré sob\[Sol]re el mar
\endverse
\beginchorus
\[Sol]Y A TU N\[Re]OMBRE CLAM\[La]ARÉ
\[Sol]EN TI MIS \[Re]OJOS \[La]FIJARÉ
EN TEMPES\[Sol]TAD
DESCANSA\[Re]RÉ EN TU \[La]PODER
PUES TUYO \[Sol]SOY \[La]HASTA EL \[Sim]FINAL
\endchorus
\beginverse
{\nolyrics \[Sim] \[La/Do#] \[Re]}
\[Sim]Tu gracia abunda \[La/Do#]en la tor\[Re]menta
Tu mano \[La]Dios, me \[Sol]guiará
\[Sim]Cuando hay temor \[La/Do#]en mi ca\[Re]mino
Tú e\[La]res fiel y no ca\[Sol]mbiarás
\endverse
\beginchorus
\[Sol]Y A TU N\[Re]OMBRE CLAM\[La]ARÉ
\[Sol]EN TI MIS \[Re]OJOS \[La]FIJARÉ
EN TEMPES\[Sol]TAD
DESCANSA\[Re]RÉ EN TU \[La]PODER
PUES TUYO \[Sol]SOY \[La]HASTA EL \[Sim]FINAL
\endchorus
\beginverse
\[Sim]Que tu Espíritu me \[Sol]guíe sin fronteras
Más al\[Re]lá de las barreras
A \[La]donde Tú me llames
\[Sol]Tú me llevas más a\[Re]llá de lo soñado
Donde \[La]puedo estar confiado
Al \[Mim]estar en tu presencia
\endverse
\beginverse
\[Sim]Que tu Espírit\[La/Do#]u me guíe\[Re] sin fronteras
Más a\[La]llá de las barreras
A \[Mim]donde Tú me llames
\[Sim]Tú me llevas \[La/Do#]más allá \[Re]de lo soñado
Dond\[La]e puedo estar confiado
Al \[Mim]estar en tu presencia  (2 veces)
{\nolyrics \[Sim] \[La/Do#] \[Re] \[La] \[Mim]}
\endverse
\beginchorus
\[Sol]Y A TU N\[Re]OMBRE CLAM\[La]ARÉ
\[Sol]EN TI MIS \[Re]OJOS \[La]FIJARÉ
EN TEMPES\[Sol]TAD
DESCANSA\[Re]RÉ EN TU \[La]PODER
PUES TUYO \[Sol]SOY \[La]HASTA EL \[Sim]FINAL
PUES TUYO \[Sol]SOY \[La]HASTA EL \[Sim]FINAL
\endchorus
   \endsong

\beginsong{Vida en abundancia}[]
\beginverse
Los \[Mi]lirios del campo \[La]y las aves del \[Si]cielo
\[Mi]No se preocupan \[La]porque están en mis \[Si]manos
T\[Do#m]ené confianza en \[La]mi 
Acá es\[Mi]toy junto a \[Si]vos
\endverse
\beginverse
\[Mi]Amá lo que sos \[La]y tus circuns\[Si]tancias
Es\[Mi]toy con vos, \[La]con tu cruz en mi esp\[Si]alda
\[Do#m]Todo terminará b\[La]ien
Yo hago \[Mi]nuevas todas las \[Si]cosas
\endverse
\beginchorus
YO \[Do#m]VENGO\[La] A TRAER\[Mi]TE \[Si]VIDA
VIDA EN AB\[Do#m]UNDANCI\[La]A, EN A\[Si]BUNDANCIA
YO SOY EL \[Do#m]CAMIN\[La]O LA VE\[Mi]RDA\[Si]D Y LA VIDA
VIDA EN AB\[Do#m]UNDANC\[La]IA, EN AB\[Si]UNDANCIA
\endchorus
\beginverse
No hice al hombre para que esté solo
Caminen juntos como hermanos
Sopórtense mutuamente
Ámense unos a otros
\endverse
\beginverse
La felicidad    de la vida eterna
Empieza   conmigo en la tierra
Sentite   vivo
La fiesta del reino comienza acá
\endverse
        \endsong

\beginsong{Noche (Hakuna)}[]
\beginverse
	{\nolyrics Intro: \[Rem] \[La#] \[Fa] \[Lam]}
\endverse

\beginverse
Por tu \[Rem]Iglesia que te espera a os\[La#]curas
Por tu \[Fa]pueblo 
que te \[Lam]reza aguardando la aurora
\[Rem]Te roga\[La#]mos, \[Fa]te ro\[Lam]gamos.
\endverse

\beginverse
Por las \[Rem]naciones paganas,
que \[La#]tienen sed de ti sin sa\[Fa]berlo
Ten pie\[Lam]dad, ten piedad.
\endverse

\beginverse
Por los \[Rem]pueblos oprimidos 
por \[La#]el totalitarismo, \[Fa]   
Y la opre\[Lam]sión de la mentira
\endverse

\beginverse
Por \[Rem]aquellos perseguidos por tu \[La#]nombre
Que se ocultan para \[Fa]orar
Y aquellos extraí\[Lam]dos de su hogar
\endverse

\beginverse
Por \[Rem]sus perseguidores, \[La#]cegados por el odio
Per\[Fa]dónales, Padre,
no \[Lam]saben lo que hacen
\endverse

\beginverse
Por \[Rem]los que no nos aman
Por los que \[La#]no sabemos amar 
Por los que \[Fa]sufren y agonizan
Y hoy duermen en \[Lam]el hospital
\endverse

\beginverse
Por los \[Rem]que hoy es su última noche
Y cuyos \[La#]ojos no verán el nuevo \[Fa]día
Ten piedad, ten \[Lam]piedad
\endverse

\beginverse
Por \[Rem]todos los que sufren 
la tenta\[La#]ción del suicidio,
Por los dis\[Fa]puestos a dejar ganar al \[Lam]mal
\endverse

\beginverse
Por \[Rem]aquellos cuyas noches 
\[La#]son interminables
Y a \[Fa]los que la angustia 
les ha \[Lam]quitado la \[La7 La Rem]paz
\endverse

\beginchorus
KYRIE E\[La#]LEISON, KYRIE E\[FA]LEISO\[Lam La7]N,
KYRIE E\[Rem]LEISON, KYRIE EL\[La#]EISON,
KYRIE E\[Fa]LEISON\[Lam La7]
\endchorus

\beginverse
Por \[Rem]aquellos que trabajan 
en la prostitución
Y se \[La#]ven obligados a vender su amor
Por \[Fa]los que caen en la trampa
Del \[Lam]vicio y las drogas
\endverse

\beginverse
Por \[Rem]los que hoy duermen en prisión
Por los que \[La#]hoy esperan su ejecución
Por \[Fa]aquellos a los que torturan
\endverse

\beginverse
Por cr\[Rem]iminales, por los ladrones
Por los que \[La#]erran en soledad
Por los que \[Fa]sufren 
la indife\[La#]rencia de los \[La7 La Rem]demás
\endverse

\beginchorus
KYRIE E\[La#]LEISON, KYRIE E\[Fa]LEISON\[Lam La7]
\endchorus

\beginverse
Por la \[Rem]ciudad, por todos sus hab\[La#]itantes
Que en sus \[Fa]sueños 
solo existas \[Lam]Tú (solo existas Tú) 
\[Rem]Por nuestros difuntos,
que aún no \[La#]han visto tu rostro
Por los \[Fa]alejados entre la multi\[Lam]tud
\endverse

\beginverse
Por \[Rem]los niños que descansan 
en el \[La#]seno de su madre
Por las muje\[Fa]res que van a dar a \[Lam]luz
\endverse

\beginverse
Para que \[Rem]reine tu paz en cada hog\[La#]ar
Por los que \[Fa]quieren 
saciar su sed de \[Lam La Rem]amar
\endverse

\beginchorus
KYRIE EL\[La#]EISON, KYRIE EL\[Fa]EISON, \[Lam La7]
KYRIE E\[Rem]LEISON, KYRIE ELE\[La#]ISON,
KYRIE ELE\[Fa]ISO\[Lam La7]N
(Ten piedad, ten piedad, ten piedad) 
\endchorus

\beginverse
Ten pied\[Rem]ad, ten pie\[La#]dad, ten pie\[Fa]dad \[Lam]
Ten pie\[Rem]dad, ten pie\[La#]dad, ten pie\[Fa]dad \[Lam]
Ten pie\[Rem]dad, ten pie\[La#]dad, ten pie\[Fa]dad \[Lam]
Ten pie\[Rem]dad, Señor, \[La#]ten piedad. \[Rem]
\endverse
\endsong

\beginsong{No temas}[by={Jaime Olguín}]
\beginverse
\[Sol]Si aún no lo \[Do]ves, \[Sol]pronto lo ve\[Re]rás,
\[Mim]todo un \[Do]Dios enamo\[Sol]rado de \[Re]ti.
\[Sol]Habla bien de \[Do]Él, que \[Sol]no menti\[Re]rás, 
\[Mim]todo lo hi\[Do]zo para ha\[Sol]certe fe\[Re]liz.
\endverse

\beginverse
Lo que ^pasa es que el ^mundo no ^quiere escu^char, 
no ^vaya a ser que ^vea, que en^tienda y se em^pape
de ^esta ale^gría que nos ^quiere rega^lar,
no ^vaya a ser que en^tienda, se con^vierta y se sa^ne
\endverse

\beginchorus
No ^temas a la ^luz, no ^temas a la ^paz, 
No ^temas a la ale^gría, no tengas ^miedo a ser fe^liz.
Porque ^Él es el ca^mino que se ha a^bierto para ^ti
Porque ^Él es el que es, ^ahora puedes ^ser feliz. ^
\endchorus
\endsong
\beginsong{Nada Cambiará (Aunque los Montes)}[]
    \beginverse
        \[Fa]Aunque los m\[Do]ontes ca\[Sol]mbien de \[Lam]lugar
        y \[Fa]se derru\[Do]mben las col\[Sol]inas
        \[Fa]no cambi\[Do]ará mi Alia\[Sol]nza de Pa\[Lam]z,
        dice Ya\[Fa]hveh, que est\[Sol]á enamora\[Do]do de ti.
    \endverse
\endsong
\beginsong{Me sobrepasa - Hakuna}[]
\beginverse
{\nolyrics Cejilla 7º traste}
{\nolyrics INTRO: \[Sol] \[Mim] \[Do] \[Re7]}
\[Sol]Tú me sondeas y me conoces,
\[Mim]cuando me siento o me levanto,
\[Do]de lejos penetras mis pensamientos;
\[Re7]distingues mi camino y mi descanso.
\endverse
\beginverse
Todas mis sendas te son familiares,
conoces mis palabras.
Me estrechas, me cubres con tu palma.
Tanto saber me sobrepasa.
\endverse
\beginchorus
\[Do]¿A \[Re]DÓNDE IRÉ LEJOS DE TU ALIENTO,
\[Sol]A \[Sol/Fa#]DÓNDE ES\[Mim]CAPARÉ DE TU MIRADA?
\[Lam]ALLÍ \[Re7]ESTÁS TÚ, ALLÍ TE ENCUENTRO.
\endchorus
\beginverse
{\nolyrics \[Sol] \[Mim] \[Do] \[Re7]}
\[Sol]Si vuelo hasta el margen de la aurora,
\[Mim]si emigro hasta el confín del mar,
\[Do]allí me alcanzará tu mano izquierda,
\[Re7]o tu derecha me agarrará.
\endverse
\beginverse
Mas si digo: «Que al menos la tiniebla
me encubra, que la luz se haga noche
en torno a mí».
La noche es clara como el día.
La tiniebla no es oscura para ti.
\endverse
\beginverse
\[Do]Tú \[Re]has creado mis entrañas,
\[Sol]en \[Sol/Fa#]el seno \[Mim]materno me tejiste.
\[Lam]Te d\[Re7]oy gracias.
\endverse
\beginchorus
\[Do]¿A \[Re]DÓNDE IRÉ LEJOS DE TU ALIENTO,
\[Sol]A \[Sol/Fa#]DÓNDE ES\[Mim]CAPARÉ DE TU MIRADA?
\[Lam]ALLÍ \[Re7]ESTÁS TÚ, ALLÍ TE ENCUENTRO. (X2)
{\nolyrics \[Sol] \[Mim] \[Do] \[Re7]}
\endchorus
   \endsong

\beginsong{Más allá de tu ventana}[]
\beginverse
{\nolyrics Cejilla 2º traste}
{\nolyrics\[La] \[La/Do#] \[Re] \[Fa#m] \[Re] \[La]}
Ooh, oooh, oooh, oooh
Sue\[La]na en mí una melodía
Que \[Re]da vida \[La]al corazón
Se \[Fa#m]levanta como un himno
Que \[Re]arde en \[La]mi interior
\endverse
\beginverse
To\[La]das las generaciones
Le\[Re]vantan\[La]do una voz
C\[Fa#m]ada hijo y cada hija
Ca\[Re]ntarán es\[La]ta canción
\endverse
\beginchorus
NO \[La]TE QUEDES CALLADO
YA \[Re]NO TE ES\[La]CONDAS MÁS
MÁS \[Mi]ALLÁ DE TU VENTANA
HAY \[Re]UN MUN\[La]DO QUE ALCANZAR
ALZA \[Fa#m]AL CIELO TUS MANOS
ALZA \[Re]AL CIE\[La]LO TU VOZ
POR \[Mi]LOS SIGLOS CANTAREMOS
AL \[Re]SEÑOR DE LA CREACIÓN
\endchorus
\beginverse
{\nolyrics\[La] \[La/Do#] \[Re] \[Fa#m] \[Re] \[La]}
Ooh, oooh, oooh, oooh
Del \[La]famoso al ignorado
Del pa\[Re]lacio al \[La]callejón
Se \[Fa#m]escucha este latido
Que \[Re]atrae mi \[La]corazón
\endverse
\beginverse
Oigo \[La]el mundo despertarse
Un \[Re]sonido ce\[La]lestial
Ca\[Fa#m]da tribu y cada lengua
Can\[Re]ta Cristo \[La]es la verdad
\endverse
\beginchorus
NO \[La]TE QUEDES CALLADO
YA \[Re]NO TE ES\[La]CONDAS MÁS
MÁS \[Mi]ALLÁ DE TU VENTANA
HAY \[Re]UN MUN\[La]DO QUE ALCANZAR
ALZA \[Fa#m]AL CIELO TUS MANOS
ALZA \[Re]AL CIE\[La]LO TU VOZ
POR \[Mi]LOS SIGLOS CANTAREMOS
AL \[Re]SEÑOR DE LA CREACIÓN
\endchorus
\beginverse
{\nolyrics\[La] \[La/Do#] \[Re] \[Fa#m] \[Re] \[La]}
Ooh, oooh, oooh, oooh
{\nolyrics\[La] \[La/Do#] \[Re] \[Fa#m] \[Re] \[La]}
Ooh, oooh, oooh, oooh
Pue\[La]do el mundo ver volver a brillar
Al \[Re]oír su voz las murallas caerán
El \[Mi]pecado no me puede atar
\[Fa#m]Nuestro Dios \[Re]Él vivo está
El \[La]perdido vuelve de la oscuridad 
El ca\[Re]utivo libre por siempre será
Nin\[Mi]gún ojo ha visto lo que Dios hará
La c\[Fa#m]reación despertará
{\nolyrics\[La] \[La/Do#] \[Re] \[Fa#m] \[Re] \[La]}
Ooh, oooh, oooh, oooh
{\nolyrics\[Fa#m] \[La/Do#] \[Re] \[Fa#m] \[Re] \[La]}
Ooh, oooh, oooh, oooh
\endverse
\beginchorus
NO \[Fa#m]TE QUEDES CALLADO
YA \[Re]NO TE ES\[La]CONDAS MÁS
MÁS \[Mi]ALLÁ DE TU VENTANA
HAY \[Re]UN MUN\[La]DO QUE ALCANZAR
ALZA \[Fa#m]AL CIELO TUS MANOS
ALZA \[Re]AL CIE\[La]LO TU VOZ
POR \[Mi]LOS SIGLOS CANTAREMOS
AL \[Re]SEÑOR DE LA CREACIÓN
\endchorus
   \endsong

\beginsong{La verdad}[by={Siloé}]

\capo{2}

\ifchorded
\beginverse
{\nolyrics Intro: \[Sol] \[Re] \[Lam7] \[Lam7]}
\endverse
\fi

\beginchorus
\memorize
\[Lam7]Sería más fá\[Sol]cil comenzar 
\[Re]por la ver\[Lam7]dad
\[Lam7]Que contar his\[Sol]torias 
que nos \[Re]puedan ale\[Lam7]jar
\endchorus

\beginverse
^A pesar de ^todo, 
me da i^gual voy a empe^zar
^A seguir cre^yendo
que has ve^nido a la ciu^dad
\endverse

\beginverse
Levanto las copas 
de gloria y redención
Cada vez que me insistes 
en volver a caminar
\endverse

\beginverse
Tus ojos se han clavado 
en la mitad de esta canción
Y lejos de escucharme, 
te pones a divagar
\endverse

\beginverse
No hay manera de engancharnos, 
es la única ambición
Que perseguimos en el centro 
de un viejo salón
\endverse

\beginverse
Y quiero que te acuerdes 
y volver a desgastar
El suelo en cuatro tiempos 
y me invites Tú a bailar
\endverse

\beginverse
Has convertido el agua en vino, 
vaya decepción
La gente no te va a creer 
y alguno lo intentó
\endverse

\beginverse
El éxito los sepultó 
y no quiso parecer
Bandera de las cosas, 
que nos dan y quitan fe
\endverse

\beginverse
En la noche en que sucumbes 
a mis ganas de intentar
Convencerte de que no conozco 
ni el bien ni el mal
\endverse

\beginverse
En la noche en la que 
nos encomendamos al azar
En el suelo de aquel cuarto 
me invitaste tu a bailar
\endverse

\beginchorus
Sería más fácil 
comenzar por la verdad x4
\endchorus
\endsong
\beginsong{Huracán}[]
    \beginverse
        \[Sol]Me he hecho tantas pregun\[Re]tas
        \[Mim]Intentando entende\[Do]r
        \[Sol]Me he lanzado a busc\[Re]arte
        Sin sab\[Mim]erte ve\[Do]r
        \[Sol]Me he asomado al abis\[Re]mo
        \[Mim]Me he atrevido a salta\[Do]r y caer
    \endverse
    \beginchorus
        Y un hura\[Sol Re]cán
        Romperá el c\[Mim]ielo desde mi gargan\[Do]ta
        Gritándo\[Sol]te: ``¿dónde está\[Re]s cuando me haces falt\[Mi Do]a?''
        \endchorus
    \beginverse
        Y me han dado respuestas
        Pero no sé qué hacer
        He prometido seguirte
        Sin entender
        Y hay un eco en lo hondo
        Que me empuja hacia ti
        Y aunque sea sin sentirte
        Te buscaré
    \endverse
    \beginchorus
        Y un huracán ... (bis)
    \endchorus
    \beginverse
        Estoy aquí, en el silencio
        Estoy aquí, en este viento
        Estoy aquí, soy este trozo de pan
        Estoy aquí, en tu lamento
        Estoy aquí, en este eco
        Estoy aquí, soy este trozo de pan
    \endverse
    \beginchorus
        Y un huracán ... (bis)
        (Entra voz de chico)
        Y un huracán ... (bis)
    \endchorus
    \beginverse
        Y tu hurac\[Sol]án
        Romper\[Re]á el cielo \[Mim]desde mi gargant\[Do]a
        Gritándom\[Sol]e cuánto me \[Re]haces falt\[Mim]a
    \endverse
\endsong
\beginsong{Venid conmigo}[]
\beginchorus
\[Re]Venid, venid conmigo a un lugar tranqui\[Sim]lo
y descan\[Sol]sad en mí vuestro cansan\[La]cio.
De\[Sol]jad que os cu\[Fa#m]re las heri\[Sim]das
que el traba\[Sl]jo por el Reino os ha deja\[La]do.
Re\[Sol]poned con mi \[La]Pan vuestras fuer\[Sim]zas,
con mi \[Sol]Vino alegrad el cora\[La]zón.
\endchorus

\beginchorus
Y ahora, ve\[Re]nid...
\endchorus
\endsong
\beginsong{Enciendeme - Hakuna}[]
    \beginverse
        \[Mi]Hoy quiero, Señor, ponerlo \[La]todo en tu presencia
        \[Do#m]Darme hasta gastarme contig\[Si]o y por Ti
        \[Mi]Hoy quiero, Señor, ponerlo \[La]todo ante tu puerta
        \[Do#m]Para en todo amarte y servir\[Si].
    \endverse
    \beginchorus
        Enciénd\[Mi]eme y déjame arder donde hag\[La]a falta,
        Enciénd\[Mi]eme y déjame ser tu \[La]luz,
        Y así po\[Mi]der llevarte hasta todas la\[La]s almas,
        Saciar la \[Mi]sed que tiene\[Si]s Tú desde l\[Mi]a cruz.
    \endchorus
    \beginverse
        Hoy quisiera madre, poner todo en tu presencia,
        Darme hasta gastarme, decirle que sí,
        Hoy te pido madre, que dejes mi puerta abierta,
        Para en todo amarle y servir
    \endverse
    \beginverse
        Enciéndeme...
    \endverse
\endsong
\beginsong{El rey de mi vida}[]
    \beginverse
        Quier\[Sol]o alabarte sin para\[Do9]r todos los días
        Que t\[Sol]u presencia sea el \[Do9]anhelo de mi vida
        Yo qu\[Sol]iero hacer tu voluntad
        Señor \[Do9]yo te quiero agradar
        Yo qui\[Mim]ero darte siempre e\[Do9]l primer lugar
        Yo qui\[Mim]ero darte siempre e\[Do9]l primer lugar
    \endverse
    \beginchorus
        Si, tu ere\[Sol]s el rey el re\[Mim]y de mi vida
        Núm\[Do9]ero uno en \[Mim7]mi corazón
        A ti yo \[Lam]te rindo todo \[Resus4]lo que soy
    \endchorus
\endsong
\beginsong{Desitjo pau}[]
\beginverse
\[Do]Ja n'hi ha \[Sol7-Fa]prou de vio\[Do-Fa]lência! \[Do]Ja no puc suporta\[Lam]r-ho \[Fa Sol7]més!
Ja n'hi ha \[Mi]prou de \[Lam-Fa]matances. Les \[Do]morts no porta\[Sol7]ran la \[Do-Sol7]pau.
\[Do]Ja n'hi ha \[Sol7-Fa]prou de les \[Do-Fa]armes, \[Do]negoci del qui en \[Lam]vendr\[Fa-Sol7]à més.
\[Do]Ja n'hi ha \[Sol7-Fa]prou de me\[Do-Fa]ntides. \[Do]Escolta 'm, desi\[Lam]tjo pau! \[Fa-Sol7]
\endverse
\beginchorus
\[Do-Sol]PAU, PA\[Mi-Lam]U PER LA \[Fa]GENT QUE PATEIX EN SOL\[Mim]EDAT,
I QUE \[Rem]VIU SENSE \[Sol]TENIR:
\[Do-Sol]PAU, PA\[Mi-Lam]U, PEL MON \[Fa]NOSTRE QUE PL\[Fam]ORA...
DESITJO \[Do]PAU.
\endchorus
\beginverse
\[Do]Ja n'hi ha \[Sol7-Fa]prou de tor\[Do-Fa]tures. La \[Do]llei del qui ha \[Lam]pogut \[Fa-Sol7]més.
Ja n'hi ha \[Mi]prou de \[Lam-Fa]cadenes que il\[Do]iguen nostra ili\[Sol7]bertat. \[Do-Sol7]
\[Do]Ja n'hi ha \[Sol7-Fa]prou de mis\[Do-Fa]eries. \[Do]Els rics potser no \[Lam]em co\[Fa-Sol7]mprendran,
\[Do]Veritat? Ja \[Sol7-Fa]n'hi ha prou \[Do-Fa]d'injustí\[Do]cies. Escolta'm, \[Lam]desitj\[Fa-Sol7]o pau.
\endverse
   \endsong

\beginsong{Dejarme hacer}[]
\beginverse
\[Mi Si7]Dejarme \[Mi Si7]hacer, dejarme \[Mi Si7]hacer,
dejarme h\[Mi]acer es cuanto \[La]pides de \[Si7]mí
Dejarme h\[La]acer de n\[Si7]uevo por \[Do#m]ti
Dejarme h\[Mi]acer en tus m\[La]anos, \[Si7 Mi]Señor.
\endverse
   \endsong

\beginsong{Amén – Ricardo Montaner, Mau y Ricky, Camilo y Eva Luna}[]
    \beginverse
        \[Fa7m]Hey, estés dond\[Sol]e estés
        \[Mim]Necesito que vengas ahora
        Que mi mu\[Lam]ndo está al revés
    \endverse
    \beginchorus
        Hey, eso ya lo sé
        Que solo te busco y te llamo
        Cuando no me dan los pies
        Hey, dime si es verdad
        Que por la noche me cuidas
        En medio de la oscuridad
    \endverse
    \beginchorus
        Hey, soy yo otra vez
        Solo sé que tengo miedo
        Pero no sé bien de qué
    \endverse
    \beginchorus
        Que nad\[Fa]ie me cuida
        Como T\[Sol]ú me cuidas
        Que nad\[Lam]ie me ama
        Como \[Mim]Tú me amas
    \endverse
    \beginchorus
        Que nadie se mueve
        Sin que Tú estés viendo
        Que nadie se salva
        Si no está en tu reino
    \endverse
    \beginchorus
        A\[Lam]mén, A\[Fa]mén
        A\[Do]quieta mi
        In\[Sol]quietud, te pido
        A\[Lam]mén, A\[Fa]mén
        Y \[Do]quítame \[Sol]el dolor
    \endchorus
    \beginverse
        \[Fa7]Es tanto lo que 
        Quiero preguntarte
        Respuesta por si
        Vuelvo a equivocarme
        \[Mim]Y si me caigo, estás
        Para levantarme
        Aunque mis hermanos
        Sean los pr\[Fa7]imeros en
        Juzgarme y señalarme los
        Los errores que Tú 
        Ya me perdonaste 
        \[Mim]Que sí me caigo, Tú estás
        Para levantarme una
        Y otra vez
        (Una y otra vez)

    \endverse
    \beginverse
        \[Fa]Y si pierdo \[Sol]el foco
        Voy contra la
        Ma\[Lam]rea y sin \[Mim]reposo
        Sé que aunque
        Esté en \[Fa]un hueco
        No est\[Sol]oy solo
        Sé que aunque esté
        En \[Lam]un hueco
        No \[Mim]estoy solo
        Y yo
    \endverse
    \beginchorus
        Amén...
    \endchorus
    \beginverse
        \[Fa]Cuida de mí, cuida de \[Sol]mí
        Deja mi alma llena de \[Lam]ti
        Cuida de mí cuida de \[Do]mí
        ¿A qué le temo?
        Si cuidas de mí
        \[Fa]Cuida de mí, cuida de \[Sol]mí
        Deja mi alma llena de \[Lam]ti
        Cuida de mí cuida de \[Do]mí
        ¿A qué le temo?
        Si cuidas de mí
    \endchorus
    \beginverse
        Cuida de mí, cuida de mí
        Deja mi alma llena de ti
        Cuida de mí cuida de mí
        ¿A qué le temo?
        Si cuidas de mí
        Cuida de mí, cuida de mí
        Deja mi alma llena de ti
        Cuida de mí cuida de mí
        ¿A qué le temo?
        Si cuidas de mí 
    \endverse
    \beginchorus
        Amén... (x3)
    \endchorus
\endsong
\beginsong{Aleluya (Harijans)}[]
\beginverse
\[Mi]Tengo una sensación, 
que no da \[Si]miedo, no,
\[Do#m]puede ser ilusión, 
pero no es \[Sol#m]cierto, es Él.
\[La]No quedará de mí, 
\[Si]una porción sint ti, 
que \[Mi]no tenga que a\[Si]grade\[Do#m]certe a \[Si]ti, Se\[La]ñor,
sé que can\[Si]taré, Alelu\[Mi]ya.
\endverse
\beginchorus
\[Mi]A - \[Mi7]LE - \[La]LU - \[Lam]YA
ALELU\[Mi]YA, ALELU\[Si]YA, ALE\[Mi]LUYA
\endchorus
   \endsong

\beginsong{Aleluya (cantad al Señor)}[]
    \beginverse
        \[Re]Aleluya (can\[Sol]tad al Se\[Re]ñor)
        Aleluya (al \[Sol]rey del a\[Re]mor)
        Ale\[Sim]lu - \[Fa#m]u - \[Sol]u \[La]ya \[Re]Aleluya 
        \[Sol]Alelu\[La]ya      Ale\[Sol]lu - \[La]u - \[Ra]ya
    \endverse
\endsong
\beginsong{Tú sabes quién soy}[by={Ixcis}]
\beginverse
\[Do]Tú\[Sol] lo sabes \[Lam]todo\[Sol]
Tú \[Fa]sabes\[Rem] quién \[Sol]soy
\[Do]Tú\[Sol] sabes mi his\[Lam]toria\[Sol]
\[Fa]Mucho me\[Sol]jor que \[Do]yo\[Sol]
\endverse

\beginchorus
Y no me \[Do]juzgas \[Sol]ni me con\[Lam]denas
\[Sol]Solo pre\[Fa]guntas: ¿\[Rem]me amas o \[Sol]no?
Y no me \[Do]juzgas \[Sol]ni me con\[Lam]denas
\[Sol]Solo me a\[Fa]brazas
¡\[Rem]Cuánto \[Sol]te quiero, Se\[Do]ñor!
\endchorus
\endsong
\beginsong{Tu palabra es mi gozo}[]
    \beginverse
        \[Mim]Tu Palabra es mi \[Do-Re]gozo,
        da luz a mis entr\[Mim]añas
        Tu Palabra es mi go\[Do-Re]zo,
        me alegra el cora\[Mim (Re-Mim)]zón.
    \endverse
\endsong
\beginsong{Tu modo}[]
\beginverse
\[Sol]Jesús, al conte\[Sim]mplar en tu \[Do]vida \[Re]
el \[Sol]modo que Tú \[Sim]tienes de tratar \[Do]a los dem\[Re]ás
\[Sol]me dejo interpe\[Sim]lar por tu tern\[Do]ura,
tu forma de am\[Re]ar nos mueve a a\[Sol]mar;
tu trato es como el a\[Sim]gua cris\[Do]talina
que limpia y acom\[Re]paña el cam\[Sol]inar.
\endverse
\beginchorus
\[Sim]JESÚS, ENSÉÑA\[Do]ME TU MODO
DE HA\[Re]CER SENTIR AL OTRO M\[Sol]ÁS HUMANO,
QUE TUS \[Sim]PASOS SEAN \[Do]MIS PASOS;
MI \[Re]MODO DE PRO\[Sol]CEDER.
\endchorus
\beginverse
Jesús, hazme sentir con tus sentimientos, 
mirar con tu mirada, comprometer mi acción, 
donarme hasta la muerte por el reino, 
defender la vida hasta la cruz, 
amar a cada uno como amigo, 
y en la oscuridad llevar tu luz
\endverse
\beginverse
Jesús, yo quiero ser compasiva con quien sufre, 
buscando la justicia, compartiendo nuestra fe, 
que encuentre una auténtica armonía 
entre lo que creo y quiero ser, 
mis ojos sean fuente de alegría, 
que abrace tu manera de ser.
\endverse
\beginverse
Quisiera conocerte, Jesús tal como eres
Tu imagen sobre mi es lo que transformará
Mi corazón en uno como el tuyo
Que sale de sí mismo para dar
Capaz de amar al padre y los hermanos
Que va sirviendo al reino en libertad
\endverse
   \endsong

 \beginsong{Tú eres mi fuerza}[]
\beginverse
\[Re]Tú \[La]eres mi \[Sol La]fuerza, tú me \[Re]das pi\[La]ernas de 
g\[Sol La]acela. Y me \[Sim]haces cami\[La]nar por las 
\[Sol]alturas, y me \[Sim]haces cami\[La]nar por las \[Sol]alturas
\endverse


\beginsong{Tomad y comed}[by={Hakuna}]
\beginverse
\[Do]Quiero alcan\[Re]zar el \[Sol]cielo con tus \[Do]Pasos,
\[Do]Quiero abra\[Re]zar la \[Sol]vida con tus Ma\[Do]nos
\[Do]Decir al \[Re]mundo con tu \[Sol]Voz hasta mo\[Do]rir
Qué \[Lam]todo un Cielo es \[La7]solo para \[Re]mí...
\endverse

\beginverse
Quiero estrenar mis ojos en tus brazos,
Quiero besar el suelo con tus Labios,
Borrar la noche y dar a la herida cicatriz.
Quiero empezar de nuevo solo en Ti.
\endverse

\beginchorus
To\[Sol]mad\[Sim], co\[Mim]med,
me en\[Do]trego por A\[Lam]mor hasta el \[Re]final
To\[Do]mad\[Re7], be\[Mim]bed,
el \[Fa]Cielo entero sabe a Vino y \[Re]Pan
To\[Sol]mad\[Sim], co\[Mim]med,
me en\[Do]trego por A\[Lam]mor hasta el \[Re]final
To\[Do]mad\[Re7], be\[Mim]bed,
el \[Do]Cielo entero \[Re]sabe a Vino y \[Sol]Pan
\endchorus
\endsong
\beginsong{Tengo un gozo en el alma}[]

\beginchorus
\[Mi]Yo tengo gozo en mi alma, gozo en mi alma,
gozo en mi alma y en mi ser,\[Si7]
¡Aleluya, gloria a Dios!
Son como \[Mi]ríos de agua viva, ríos de agua viva,
ríos de agua viva en \[Si7]mi \[Mi]ser.
\endchorus

\endsong

	\end{songs}
\end{document}